% =============================================================
% Bölüm 1: Bilgisayar Mimarisine Giriş
% =============================================================
\chapter{Bilgisayar Mimarisine Giriş}
\label{chap:giris}
\bolumfoto[Yaklaşık 1500 yıldır ayakta duran Ayasofya, çağının en ileri mühendislik bilgisiyle inşa edilmiştir. Bilgisayar mimarisi de benzer bir mühendislik disiplinidir: sınırlı kaynaklarla en yüksek başarımı elde etmeye çalışır.]{gorseller/bolum01-ayasofya.jpg}{Ayasofya -- İstanbul}

\begin{tcolorbox}[
    colback=blokyesil!15,
    colframe=sinyalyesil!60!black,
    title={\textbf{Öğrenme Hedefleri}},
    fonttitle=\bfseries,
    rounded corners,
    boxrule=0.8pt
]
Bu bölümü tamamladığınızda aşağıdakileri yapabileceksiniz:
\begin{itemize}[nosep, leftmargin=*]
    \item Bilgisayar tarihinin beş neslini ve her neslin ayırt edici teknolojisini açıklamak.
    \item Moore Yasası ve Dennard ölçeklemesinin ne olduğunu ve neden sona erdiğini tartışmak.
    \item Bir bilgisayarın beş ana bileşenini sıralamak ve her birinin işlevini tanımlamak.
    \item Yonga üretim sürecini (fotolitografi, dağlama, kaplama) ana hatlarıyla betimlemek.
    \item Üst düzey dilden makine koduna uzanan derleme zincirini açıklamak.
    \item RISC ile CISC tasarım felsefelerini karşılaştırmak.
    \item RISC-V buyruk kümesinin neden tercih edildiğini gerekçelendirmek.
    \item Güç duvarının mimari tasarıma etkisini ve heterojen hesaplama kavramını tartışmak.
\end{itemize}
\end{tcolorbox}

\section{Giriş}
\label{sec:giris}

Cebinizdeki akıllı telefon, 1969'da insanı Ay'a götüren Apollo rehberlik bilgisayarından milyonlarca kat daha güçlü. Bugün bilgisayarlar yalnızca masaüstünde değil, saatimizde, buzdolabımızda, otomobilimizde ve veri merkezlerindeki devasa sunucu çiftliklerinde karşımıza çıkıyor. Bu yaygınlığın arkasında yarım yüzyılı aşan sürekli bir mimari yenilenme var.

\emph{Bilgisayar mimarisi}\index{bilgisayar mimarisi}, donanım ile yazılım arasındaki sınırı tanımlayan ve işlemci, bellek ve giriş/çıkış birimlerinin nasıl bir araya geleceğini belirleyen mühendislik dalıdır. Verimli yazılım geliştirmek isteyen bir programcı için donanımın çalışma ilkelerini anlamak, yüksek başarımlı donanım tasarlamak isteyen bir mühendis için de yazılım gereksinimlerini kavramak vazgeçilmez. Üstelik alan, tarihsel bir dönüm noktasında: transistörler küçüldükçe güç yoğunluğunun sabit kalacağını öngören Dennard ölçeklemesinin\index{Dennard ölçeklemesi} sona ermesi, yapay zekânın yükselişi ve açık kaynak donanım hareketleri bilgisayar mimarisini her zamankinden daha heyecanlı bir araştırma alanı yapıyor.

Bu bölümde önce bilgisayarların tarihsel gelişimini, ardından bir bilgisayarın temel bileşenlerini ve işlemcinin nasıl üretildiğini göreceğiz. Sonra programların bu donanım üzerinde nasıl çalıştırıldığını, mimariyi şekillendiren etkenleri, güç tüketimi sorununu, RISC-V buyruk kümesinin neden tercih edildiğini ve yapay zekânın bilgisayar mimarisine getirdiği yeni ufukları ele alacağız.



\section{Bilgisayarların Gelişimi}
\label{sec:bilgisayar-gelisimi}

Abaküsle başlayan hesap yapma gereksinimi, önce mekanik, sonra da elektronik hesap makineleriyle karşılanmaya çalışıldı. Bilinen anlamda ilk mekanik hesap makinesini Blaise Pascal\index{Pascal, Blaise} 1642'de tasarladı. Pascal, toplama ve çıkarma işlemlerini yapabilen bu makineyi vergi tahsildarı babasına yardım etmek amacıyla geliştirdi. Şekil~\ref{fig:pascal-makine}'de Pascal'ın yaptığı bu ilk makinenin bir resmi görülüyor.

\begin{figure}[H]
\centering
\includegraphics[width=0.85\textwidth]{gorseller/pascaline.jpg}
\caption{Blaise Pascal'ın 1642'de yaptığı mekanik hesap makinesi (Pascaline)}
\label{fig:pascal-makine}
\end{figure}

19.~yüzyılın sonlarından itibaren hesaplama makineleri hızla gelişti. 1890'da Herman Hollerith\index{Hollerith, Herman} delikli kart sistemiyle çalışan bir hesaplama düzeneği geliştirdi; bu sistem ABD nüfus sayımında başarıyla kullanıldı.

Bu gelişmelere koşut olarak Alan Turing\index{Turing, Alan} 1936'da yayımladığı makalesinde~\cite{turing1936}, sonradan \emph{Turing makinesi}\index{Turing makinesi} olarak anılacak soyut hesaplama modelini tanımlayarak hesaplanabilirliğin\index{hesaplanabilirlik} teorik temelini attı. Bu model, bir algoritma ile çözülebilecek her problemin basit bir bant, okuma/yazma kafası ve sonlu durum tablosundan oluşan bir makineyle çözülebileceğini gösteriyordu. Turing makinesi kavramı, bilgisayar biliminin kurucu taşlarından biri olarak kabul edilir ve günümüzde bile algoritma karmaşıklığı ve hesaplanabilirlik kuramının temelini oluşturur.

2.~Dünya Savaşı döneminde hesaplama ihtiyacı iyice arttı ve birçok öncü makine üretildi. Turing, İngiltere'de Bletchley Park'ta\index{Bletchley Park} Alman Enigma\index{Enigma} şifre makinesinin kırılması için elektro-mekanik Bombe\index{Bombe} makinesini tasarladı. İngilizler daha sonra daha gelişmiş şifrelerin çözümü için Colossus\index{Colossus}'u geliştirdi. Aynı dönemde IBM ile çalışan Howard H.~Aiken\index{Aiken, Howard H.} 1944'te rölelerden oluşan Mark~I\index{Mark I}'i tasarladı. Bilgisayar tarihinin sonraki dönüm noktaları aşağıda nesiller halinde ele alınıyor.

\subsubsection*{Birinci Nesil: Vakum Tüpleri (1940'lar--1950'ler)}
Günümüz bilgisayarlarının atası sayılan ENIAC\index{ENIAC} (Electronic Numerical Integrator And Computer), 1945'te Pennsylvania Üniversitesi'nde tamamlandı. Yaklaşık 18.000 vakum tüpü\index{vakum tüpü} kullanan, 30~tonluk ve yaklaşık 150~kW güç tüketen bu dev makine yalnızca dört işlem ve basit karşılaştırma yapabiliyordu. Kısa süre sonra Von Neumann'ın\index{Von Neumann} teorik katkılarıyla tasarlanan EDVAC\index{EDVAC}, programı ve veriyi aynı bellekte saklayarak çığır açtı. Birinci nesil bilgisayarlar makine dili\index{makine dili} ile programlanıyordu; vakum tüplerinin kısa ömrü, büyük boyutu ve yüksek ısı üretimi en önemli sınırlılıkları oluşturuyordu.

\subsubsection*{İkinci Nesil: Transistörler (1950'ler--1960'lar)}
1947'de Bell Laboratuvarları'nda William Shockley, John Bardeen ve Walter Brattain tarafından icat edilen transistör\index{transistör}, vakum tüplerinin yerini aldı. Transistörler çok daha küçük, güvenilir ve ucuzdu; güç tüketimleri de düşüktü. Bu dönemde çevirici dil\index{çevirici dil} (Assembly\index{Assembly}) kullanılmaya başlandı ve programlama önemli ölçüde kolaylaştı. Fortran ve COBOL gibi ilk üst düzey diller de bu nesilde ortaya çıktı.

\subsubsection*{Üçüncü Nesil: Tümleşik Devreler (1960'lar--1970'ler)}
Texas Instruments'ta Jack Kilby\index{Kilby, Jack}~\cite{kilby1964} ve Fairchild Semiconductor'da Robert Noyce\index{Noyce, Robert} birbirinden bağımsız olarak tümleşik devreyi\index{tümleşik devre} geliştirdi. Tüm bileşenlerin ve bağlantıların aynı yarıiletken malzeme üzerinde üretilmesi, devrelerin boyutunu ve maliyetini büyük ölçüde düşürdü. Bu dönemde işletim sistemleri\index{işletim sistemi} yaygınlaştı ve birden çok programın aynı belleği paylaşması mümkün hale geldi.

\subsubsection*{Dördüncü Nesil: Mikroişlemciler (1970'ler--2000'ler)}
1971'de Intel'in 4004\index{Intel 4004} yongası, işlemcinin tüm bileşenlerini tek bir tümleşik devre üzerinde toplayan ilk \emph{mikroişlemci}\index{mikroişlemci} oldu. Moore yasasının\index{Moore yasası} öngördüğü transistör artışı sayesinde işlemciler hızla güçlendi; 1981'de kişisel bilgisayar\index{kişisel bilgisayar} (PC) kavramı doğdu ve bilgisayarlar geniş kitlelere ulaştı. Bu nesil boyunca saat sıklığı sürekli artırılarak başarım kazanımı sağlandı.

\subsubsection*{Beşinci Nesil: Çok Çekirdekli ve Heterojen Mimariler (2000'ler--günümüz)}
2000'li yılların ortalarından itibaren güç duvarı\index{güç duvarı} ve Dennard ölçeklemesinin\index{Dennard ölçeklemesi} sona ermesi saat sıklığını artırma stratejisini sürdürülemez kıldı. Tasarımcılar, tek bir yonga üzerine birden çok çekirdek\index{çok çekirdekli} yerleştirerek başarım artışını koşut hesaplamaya kaydırdı. Günümüzde MİB\index{MİB} çekirdeklerinin yanına GPU\index{GPU}, NPU\index{NPU} ve diğer özelleşmiş hızlandırıcılar eklenerek \emph{heterojen} mimariler oluşturuluyor. Açık kaynaklı RISC-V\index{RISC-V} buyruk kümesi mimarisi de bu dönemin önemli gelişmeleri arasında yer alıyor.

\begin{bilginot}[Robert Dennard ve Dennard Ölçeklemesi]
IBM araştırmacısı Robert Dennard\index{Dennard, Robert}, 1974'te yayımladığı makalede~\cite{dennard1974} transistör boyutlarının küçültülmesinin etkilerini inceledi. Dennard'ın ortaya koyduğu ölçekleme kuralına göre transistörler küçüldükçe çalışma gerilimleri ve akımları da orantılı biçimde azalır; dolayısıyla birim alan başına güç tüketimi yaklaşık sabit kalır. Bu ilke sayesinde 1970'lerden 2000'lerin başına kadar her yeni üretim teknolojisiyle hem daha fazla transistör sığdırılabildi hem de saat sıklığı artırılabildi; çünkü ek transistörler fazladan güç tüketmiyordu. Ancak 2006 civarında transistör boyutları 65~nm'nin altına indiğinde sızıntı akımları\index{sızıntı akımı} baskın hale geldi ve Dennard ölçeklemesi geçerliliğini yitirdi. Bu dönüm noktası, işlemci tasarımında saat sıklığı yarışının sona ermesine ve çok çekirdekli mimarilere geçişe yol açtı.
\end{bilginot}

\begin{terimnotu}
\textbf{donanım} (\emph{hardware})\,: \emph{Donanmak} (to equip) fiilinden \emph{-ım} ekiyle. Bilgisayarın fiziksel ``donanımı'', tıpkı bir geminin donanması gibi.

\textbf{yazılım} (\emph{software})\,: \emph{Yazılmak} (to be written) fiilinin edilgen biçiminden \emph{-ım} ekiyle. Donanımın üzerinde çalışan ``yazılmış'' programlar.

\textbf{yonga} (\emph{chip})\,: \emph{Yontmak} (to carve, to whittle) fiilinin kökü \emph{yon-} ve \emph{-ga} ekiyle; ağaçtan yontularak çıkarılan ince parça demektir. İngilizce \emph{chip} de ``koparılan küçük parça'' anlamındadır. Paketlenmiş tümleşik devreyi ifade eder.

\textbf{kırmık} (\emph{die})\,: \emph{Kırmak} (to break) fiilinden \emph{-ık} ekiyle; silisyum dilimden (wafer) kesilerek ``kırılıp'' ayrılan, henüz paketlenmemiş yarıiletken parçasıdır.

\textbf{bellek} (\emph{memory})\,: \emph{Bellemek} (to memorize) fiilinden \emph{-ek} araç yapım ekiyle; ``belleme aracı'' demektir. Yazar Nurullah Ataç'ın önerisiyle Türkçeye kazandırılmış, Arapça kökenli \emph{hafıza} sözcüğünün yerini almıştır.

\textbf{tümleşik devre} (\emph{integrated circuit})\,: \emph{Tümleşmek} (to integrate, to merge into a whole) fiilinden; tüm bileşenlerin tek bir ``bütün'' haline getirildiği devre.

\textbf{merkezi işlem birimi (MİB)} (\emph{CPU, Central Processing Unit})\,: Bilgisayarın hesaplama ve denetim işlevlerini gerçekleştiren ana birim. \emph{Merkez} (center) ve \emph{işlemek} (to process) köklerinden; İngilizce \emph{CPU} kısaltması da yaygın biçimde kullanılır.
\end{terimnotu}

\begin{karsilastirma}[Nesiller Boyunca: Boyut, Güç ve Başarım]
Her nesil, önceki neslin kısıtlarını aşarak yeni olanaklar açmıştır:

\begin{center}
\small
\begin{tabular}{lccc}
\toprule
\textbf{Nesil} & \textbf{Anahtarlama hızı} & \textbf{Tipik güç} & \textbf{Tipik boyut} \\
\midrule
Vakum tüpü & $\sim$\,$\mu$s & $\sim$150\,kW & Oda \\
Transistör & $\sim$100\,ns & $\sim$10\,kW & Dolap \\
Tümleşik devre & $\sim$10\,ns & $\sim$1\,kW & Masa \\
Mikroişlemci & $\sim$1\,ns & $\sim$100\,W & Yonga \\
Çok çekirdekli & $<$1\,ns & 5--300\,W & Yonga \\
\bottomrule
\end{tabular}
\end{center}

Bu tabloda çarpıcı bir örüntü görülür: anahtarlama hızı her nesilde yaklaşık 10 kat artarken, sistem boyutu odadan tek bir yongaya küçülmüştür. Güç tüketimi ise beşinci nesilde artık düşürülemez hale gelmiştir; bu durum, ``güç duvarı'' kavramının ortaya çıkmasına yol açmıştır.
\end{karsilastirma}

\subsection{Moore Yasası ve Transistör Artışı}
\label{sec:moore-yasasi}

Bu beş nesil boyunca yaşanan hızlı gelişimin arkasında, transistör teknolojisindeki üstel ilerleme yatıyor. Intel marka işlemcilerin tarihte gelişimi Tablo~\ref{tab:intel-gelisim}'de veriliyor.

\begin{table}[ht!]
\centering
\caption{Intel işlemcilerin gelişimi}
\label{tab:intel-gelisim}
\footnotesize
\begin{tabular}{llllp{5cm}}
\toprule
\textbf{Yonga} & \textbf{Tarih} & \textbf{Saat Sıklığı} & \textbf{Transistör Sayısı} & \textbf{Açıklama} \\
\midrule
4004       & 1971 & 0,74 MHz   & 2.300   & Bir yongadaki ilk mikroişlemci \\
8008       & 1972 & 0,5--0,8 MHz & 3.500   & İlk 8 bitlik mikroişlemci \\
8080       & 1974 & 2 MHz      & 6.000   & İlk genel amaçlı işlemci \\
8086       & 1978 & 5--10 MHz  & 29.000  & İlk 16 bitlik mikroişlemci \\
8088       & 1979 & 5--8 MHz   & 29.000  & IBM PC'de kullanılan işlemci \\
80286      & 1982 & 8--12 MHz  & 134.000 & Bellek koruması bulunur \\
80386      & 1985 & 16--33 MHz & 275.000 & İlk 32 bitlik mikroişlemci \\
80486      & 1989 & 25--100 MHz& 1,2M    & 8 KB önbelleğe sahip \\
Pentium    & 1993 & 60--200 MHz& 3,1M    & İlk süperskaler x86 (iki boru hattı) \\
Pentium MMX& 1997 & 166--233 MHz& 4,5M   & SIMD buyruk desteği (MMX) \\
Pentium Pro& 1995 & 150--200 MHz& 5,5M   & İki düzey önbellek \\
Pentium II & 1997 & 233--400 MHz& 7,5M   & Pentium Pro mimarisi + MMX \\
Pentium III& 1999 & 500 MHz    & 9,5M    & Güç tasarrufu seviyeleri \\
Pentium 4  & 2000 & 1,5 GHz    & 42M     & NetBurst mikromimari; derin boru hattı \\
Pentium D  & 2005 & 3,2 GHz    & 230M    & İlk 2 çekirdekli masaüstü \\
Core 2 Duo & 2006 & 2,93 GHz   & 291M    & Pentium M'in gelişmiş hali \\
Quad Core  & 2007 & $>$3 GHz   & 410M    & İki tane Core 2 Duo \\
Xeon       & 2009 & 2,26 GHz   & 781M    & Ölçeklenebilirlik odaklı \\
Sandy Bridge & 2011 & 3,4 GHz   & 1,16Mr & AVX, tümleşik GPU \\
Haswell    & 2013 & 3,5 GHz    & 1,4Mr  & AVX2, FMA3 \\
Skylake    & 2015 & 4,0 GHz    & 1,75Mr & DDR4 desteği \\
Coffee Lake& 2018 & 3,6 GHz    & ${\sim}$3Mr & İlk 6 çekirdekli masaüstü \\
Alder Lake & 2021 & 5,2 GHz    & ${\sim}$10Mr & İlk karma mimari (P+E çekirdek) \\
Raptor Lake& 2022 & 5,8 GHz    & ${\sim}$26Mr & 8P+16E çekirdek \\
\bottomrule
\multicolumn{5}{l}{\scriptsize M = milyon; Mr = milyar; ``${\sim}$'' tahminî değer.}
\end{tabular}
\end{table}

Intel'in kurucularından Gordon Moore\index{Moore, Gordon} 1965 yılında bir gözlem ortaya koydu~\cite{moore1965}: Bir yonga üzerindeki transistör sayısı yaklaşık her iki yılda bir ikiye katlanıyor. Moore bu öngörüyü ilk olarak her yıl ikiye katlanma biçiminde ifade etmiş, 1975'te ise süreyi yaklaşık iki yıla revize etmiştir. ``Moore Yasası''\index{Moore yasası} olarak bilinen bu kural, Newton yasaları gibi değişmez bir fizik yasası değil, gözleme dayalı ampirik bir eğilimdir. Buna karşın yarıiletken endüstrisi on yıllar boyunca bu eğilimi bir geliştirme hedefi olarak benimsemiş ve tasarım yol haritalarını buna göre şekillendirmiştir. Tablodaki veriler incelendiğinde, 1971'den 2022'ye kadar transistör sayısının yaklaşık 10 milyon kat arttığı görülür. Bu büyüme hızı, insan uygarlığının başka hiçbir teknoloji alanında gözlemlenmemiştir. Şekil~\ref{fig:moore-yasasi}'de bu büyümenin logaritmik ölçekteki grafiği verilmiştir.

\begin{figure}[H]
\centering
\begin{tikzpicture}[>=Stealth, x=0.22cm, y=1.15cm]
    % Eksen aralıkları: X -> 1970-2025, Y -> log10 ölçeği (10^3 ile 10^11 arası)
    % Y konumu: log10(transistör) - 3  =>  10^3->0, 10^4->1, ..., 10^11->8

    % Y ekseni
    \draw[thick, ->] (1968,0) -- (1968,8.8) node[above, font=\small\bfseries] {Transistör Sayısı};
    % Y ekseni etiketleri
    \foreach \lev/\lab in {0/$10^3$, 1/$10^4$, 2/$10^5$, 3/$10^6$, 4/$10^7$, 5/$10^8$, 6/$10^9$, 7/$10^{10}$, 8/$10^{11}$} {
        \draw (1968,\lev) -- +(-0.8,0);
        \node[left, font=\footnotesize] at (1967.2,\lev) {\lab};
    }
    % Y ekseni ızgara çizgileri
    \foreach \lev in {0,1,...,8} {
        \draw[gray!25, thin] (1968,\lev) -- (2026,\lev);
    }

    % X ekseni
    \draw[thick, ->] (1968,0) -- (2027,0) node[right, font=\small\bfseries] {Yıl};
    % X ekseni etiketleri
    \foreach \yr in {1970,1975,...,2025} {
        \draw (\yr,0) -- ++(0,-0.15);
        \node[below, font=\footnotesize] at (\yr,-0.15) {\yr};
    }

    % Veri noktaları: (yıl, log10(transistör)-3)
    % 4004: 1971, 2300 -> log10(2300)=3.362 -> y=0.362
    % 8080: 1974, 6000 -> log10(6000)=3.778 -> y=0.778
    % 8086: 1978, 29000 -> log10(29000)=4.462 -> y=1.462
    % 80286: 1982, 134000 -> log10(134000)=5.127 -> y=2.127
    % 80386: 1985, 275000 -> log10(275000)=5.439 -> y=2.439
    % 80486: 1989, 1200000 -> log10(1.2M)=6.079 -> y=3.079
    % Pentium: 1993, 3100000 -> log10(3.1M)=6.491 -> y=3.491
    % Pentium III: 1999, 9500000 -> log10(9.5M)=6.978 -> y=3.978
    % Pentium 4: 2000, 42000000 -> log10(42M)=7.623 -> y=4.623
    % Core 2 Duo: 2006, 291000000 -> log10(291M)=8.464 -> y=5.464
    % Quad Core: 2007, 410000000 -> log10(410M)=8.613 -> y=5.613
    % Sandy Bridge: 2011, 1.16e9 -> log10=9.064 -> y=6.064
    % Haswell: 2013, 1.4e9 -> log10=9.146 -> y=6.146
    % Coffee Lake: 2018, 3e9 -> log10=9.477 -> y=6.477
    % Alder Lake: 2021, 1e10 -> log10=10.0 -> y=7.0
    % Raptor Lake: 2022, 2.6e10 -> log10=10.415 -> y=7.415

    % Eğilim çizgisi (yaklaşık Moore Yasası: her 1.5 yılda 2x)
    % 1971'de y=0.362, eğim = log10(2)/1.5 = 0.2007/yıl
    % 2024'te: 0.362 + (2024-1971)*0.2007 = 0.362 + 10.637 = 10.999
    \draw[dashed, sinyalkirmizi, thick] (1971, 0.362) -- (2024, 8.5)
        node[above left, font=\footnotesize, text=sinyalkirmizi] {Moore Yasası};

    % Noktalar ve etiketler
    \filldraw[sinyalmavi] (1971, 0.362) circle (2.5pt)
        node[above right, font=\scriptsize] {4004};
    \filldraw[sinyalmavi] (1974, 0.778) circle (2.5pt)
        node[above right, font=\scriptsize] {8080};
    \filldraw[sinyalmavi] (1978, 1.462) circle (2.5pt)
        node[above right, font=\scriptsize] {8086};
    \filldraw[sinyalmavi] (1982, 2.127) circle (2.5pt)
        node[above right, font=\scriptsize] {80286};
    \filldraw[sinyalmavi] (1985, 2.439) circle (2.5pt)
        node[above left, font=\scriptsize] {80386};
    \filldraw[sinyalmavi] (1989, 3.079) circle (2.5pt)
        node[above right, font=\scriptsize] {80486};
    \filldraw[sinyalmavi] (1993, 3.491) circle (2.5pt)
        node[above right, font=\scriptsize] {Pentium};
    \filldraw[sinyalmavi] (1999, 3.978) circle (2.5pt)
        node[above left, font=\scriptsize] {Pentium III};
    \filldraw[sinyalmavi] (2000, 4.623) circle (2.5pt)
        node[right, font=\scriptsize] {Pentium 4};
    \filldraw[sinyalmavi] (2006, 5.464) circle (2.5pt)
        node[above left, font=\scriptsize] {Core 2 Duo};
    \filldraw[sinyalmavi] (2007, 5.613) circle (2.5pt)
        node[below right, font=\scriptsize] {Quad Core};
    \filldraw[sinyalmavi] (2011, 6.064) circle (2.5pt)
        node[above left, font=\scriptsize] {Sandy Bridge};
    \filldraw[sinyalmavi] (2013, 6.146) circle (2.5pt)
        node[below right, font=\scriptsize] {Haswell};
    \filldraw[sinyalmavi] (2018, 6.477) circle (2.5pt)
        node[below right, font=\scriptsize] {Coffee Lake};
    \filldraw[sinyalmavi] (2021, 7.0) circle (2.5pt)
        node[above left, font=\scriptsize] {Alder Lake};
    \filldraw[sinyalmavi] (2022, 7.415) circle (2.5pt)
        node[right, font=\scriptsize] {Raptor Lake};

    % Noktaları birleştiren çizgi
    \draw[sinyalmavi, thin]
        (1971, 0.362) -- (1974, 0.778) -- (1978, 1.462) --
        (1982, 2.127) -- (1985, 2.439) -- (1989, 3.079) --
        (1993, 3.491) -- (1999, 3.978) -- (2000, 4.623) --
        (2006, 5.464) -- (2007, 5.613) -- (2011, 6.064) --
        (2013, 6.146) -- (2018, 6.477) -- (2021, 7.0) --
        (2022, 7.415);
\end{tikzpicture}
\caption[Intel işlemcilerdeki transistör sayısının artışı]{Yıllara göre Intel işlemcilerdeki transistör sayısının artışı (grafik logaritmik ölçektedir). Kesik çizgi Moore Yasası'nın öngördüğü eğilimi göstermektedir.}
\label{fig:moore-yasasi}
\end{figure}

\begin{bilginot}[Moore Yasası'nın Geleceği]
Moore Yasası, yarım yüzyılı aşkın bir süre endüstrinin yol haritasını belirledi. Ancak transistör boyutları atom ölçeğine yaklaştıkça (3~nm ve altı) fiziksel sınırlar belirginleşiyor. Dennard ölçeklemesinin de sona ermesiyle birlikte saat sıklığını artırma stratejisi sürdürülemez hale geldi ve tasarımcılar çok çekirdekli ve heterojen mimarilere yöneldi. Günümüzde başarım artışı, tek çekirdeğin hızlandırılması yerine koşut hesaplama ve özel amaçlı hızlandırıcılarla (GPU, TPU, NPU) sağlanıyor.
\end{bilginot}

\begin{ornek}[1.1]
\label{orn:moore-hesaplama}
Moore yasasına göre transistör sayısı her 2 yılda ikiye katlanmaktadır.
\begin{enumerate}[(a)]
\item 1971'de 2300 transistörlü Intel 4004'ten başlayarak, 2024'te beklenen transistör sayısını hesaplayın.
\item Gerçek değer (Apple M4 Ultra: $\sim$46 milyar) ile karşılaştırın.
\end{enumerate}

\textbf{Çözüm:}

(a) 1971'den 2024'e kadar 53 yıl $= 26{,}5$ ikiye katlama dönemi geçmiştir:
\[
N = 2300 \times 2^{26{,}5} \approx 2300 \times 94{,}9 \times 10^{6} \approx 218 \times 10^{9} \approx 218 \text{ milyar}
\]

(b) Gerçek değer ($\sim$46 milyar), Moore yasasının öngörüsünün ($\sim$218 milyar) yaklaşık beşte biri kadardır. Bu sapmanın başlıca nedenleri: Dennard ölçeklemesinin 2006 civarında sona ermesi, saat sıklığının sabitlenmesi ve güç duvarının transistör sayısı artışını yavaşlatmasıdır. Yine de 53 yılda $\sim$20 milyon katlık artış, insanlık tarihinin en hızlı teknolojik ilerlemesidir.
\end{ornek}

\section{Bilgisayar Nedir?}
\label{sec:bilgisayar-nedir}

Bilgisayar veri alabilen, aldığı veriler üzerinde işlemler yapan ve sonuç üreten bir makinedir. Bilgisayarlar veri yolu, denetim, bellek, giriş ve çıkış olmak üzere beş ana bileşenden oluşur. Şekil~\ref{fig:bes-bilesen}'te gösterildiği üzere veri yolu ve denetim birimleri işlemciyi oluşturur.

\begin{figure}[H]
\centering
\begin{tikzpicture}[>=Stealth, scale=0.88, every node/.style={scale=0.88},
    bilesen/.style={draw, thick, rounded corners=3pt, minimum height=1.1cm, align=center, font=\small}]

    % === İşlemci çerçevesi ===
    \node[draw, very thick, rounded corners=5pt, minimum width=7.6cm, minimum height=2.0cm,
          fill=blokmavi-acik, dashed] (islemci) {};
    \node[font=\small\bfseries, anchor=north] at ([yshift=2pt]islemci.north) {İşlemci (\emph{Processor})};

    % Veri Yolu (sol iç blok)
    \node[bilesen, fill=blokmavi, minimum width=2.8cm]
          at ([xshift=-1.4cm, yshift=-0.1cm]islemci.center) (veriyolu) {\textbf{Veri Yolu}\\[-2pt]{\scriptsize\itshape Datapath}};
    % Denetim (sağ iç blok)
    \node[bilesen, fill=blokkirmizi-acik, minimum width=2.8cm]
          at ([xshift=1.4cm, yshift=-0.1cm]islemci.center) (denetim) {\textbf{Denetim}\\[-2pt]{\scriptsize\itshape Control}};
    % Veri yolu <-> Denetim iç bağlantısı (ince çizgi, açıklayıcı)
    \draw[-, thin, gray!60] (veriyolu.east) -- (denetim.west);

    % === Bellek (üstte) ===
    \node[bilesen, fill=bloksari, minimum width=3.2cm,
          above=0.7cm of islemci]
          (bellek) {\textbf{Bellek}\\[-2pt]{\scriptsize\itshape Memory}};

    % === Giriş (sol altta) ===
    \node[bilesen, fill=blokyesil-acik, minimum width=2.4cm,
          below left=0.7cm and 0.1cm of islemci]
          (giris) {\textbf{Giriş}\\[-2pt]{\scriptsize\itshape Input}};

    % === Çıkış (sağ altta) ===
    \node[bilesen, fill=blokyesil-acik, minimum width=2.4cm,
          below right=0.7cm and 0.1cm of islemci]
          (cikis) {\textbf{Çıkış}\\[-2pt]{\scriptsize\itshape Output}};

    % === Bağlantı çizgileri ===
    % İşlemci <-> Bellek
    \draw[<->, very thick, sinyalmavi] (islemci.north) -- (bellek.south);

    % Giriş -> İşlemci
    \draw[->, very thick, sinyalmavi] (giris.north) -- ([xshift=-1.2cm]islemci.south);

    % İşlemci -> Çıkış
    \draw[->, very thick, sinyalmavi] ([xshift=1.2cm]islemci.south) -- (cikis.north);

    % Giriş -> Bellek (sol taraftan)
    \draw[->, thick, gray] (giris.west) -- ++(-0.5,0) |- (bellek.west);

    % Bellek -> Çıkış (sağ taraftan)
    \draw[->, thick, gray] (bellek.east) -| ([xshift=0.5cm]cikis.east) -- (cikis.east);

\end{tikzpicture}
\vspace{-2mm}
\caption{Bilgisayarın 5 ana bileşeni}
\label{fig:bes-bilesen}
\end{figure}

\subsection{Bilgisayarın Parçaları}
\label{sec:bilgisayar-parcalari}

Günümüzde ``bilgisayar'' kavramı yalnızca masaüstü ya da dizüstü bilgisayarlarla sınırlı değil. Akıllı telefonlar, tabletler, akıllı saatler, sunucular ve araç içi gömülü sistemler de birer bilgisayar. Boyutları ve kullanım alanları farklı olsa da tüm bu aygıtlar Şekil~\ref{fig:bes-bilesen}'te gösterilen beş temel bileşeni içerir.

\subsubsection*{İşlemci}
Bilgisayarın beyni olan \emph{işlemci}\index{işlemci} (\emph{processor}), programdaki buyrukları yürüten birimdir. Şekil~\ref{fig:bes-bilesen}'te görüldüğü gibi işlemci iki alt birimden oluşur: \textbf{veri yolu}\index{veri yolu} (\emph{datapath}) ve \textbf{denetim}\index{denetim birimi} (\emph{control}). Veri yolu, aritmetik ve mantık işlemlerini gerçekleştiren donanım yollarını (toplayıcılar, yazmaçlar, çoklayıcılar gibi birimleri) barındırır. Denetim birimi ise her buyruğun hangi adımda ne yapacağını belirleyen sinyalleri üretir; bir orkestra şefi gibi veri yolundaki birimlerin doğru zamanda doğru işi yapmasını sağlar. İşlemci tasarımı Bölüm~\ref{chap:islemci-tasarimi}'da, boru hattıyla başarım artırma teknikleri ise Bölüm~\ref{chap:boruhatti}'de ayrıntılı olarak ele alınacaktır.

\subsubsection*{Bellek}
Programların ve verilerin saklandığı bileşene \emph{bellek}\index{bellek} denir. Bilgisayarlarda iki temel bellek türü bulunur: uçucu (\emph{volatile}) ve uçucu olmayan (\emph{non-volatile}). Uçucu bellekler (örneğin DRAM) güç kesildiğinde içeriklerini kaybeder ve program çalışırken veriye hızlı erişim sağlamak için kullanılır. Uçucu olmayan bellekler ise (SSD, sabit disk gibi) güç olmadan da veriyi korur ve kalıcı depolama için tercih edilir. Bellek hiyerarşisi, hız ve maliyet dengesini sağlamak amacıyla farklı türdeki belleklerin katmanlı biçimde düzenlenmesiyle oluşturulur; bu konu Bölüm~\ref{chap:bellek}'de ayrıntılı olarak ele alınacaktır.

\subsubsection*{Giriş aygıtları}
Bilgisayara veri sağlayan birimlere \emph{giriş aygıtı}\index{giriş aygıtı} denir. Klavye ve fare onlarca yıldır en yaygın giriş aygıtları olmuştur; akıllı telefonlarda ise dokunmatik ekran hem giriş hem çıkış işlevi görür. Bunların yanı sıra ivme ölçer, GPS alıcısı ve parmak izi okuyucu gibi çeşitli algılayıcılar da birer giriş aygıtıdır.

Yapay zekâ çağıyla birlikte \emph{mikrofon}\index{mikrofon} ve \emph{kamera}\index{kamera} giderek daha önemli giriş aygıtları hâline geliyor. Büyük dil modelleri ve konuşma tanıma sistemleri sayesinde kullanıcılar bilgisayarla doğal dilde sesli iletişim kurabiliyor. Kameralar ise yalnızca fotoğraf çekmekle kalmayıp, yapay zekâ destekli görüntü anlama, yüz tanıma, hareket algılama ve artırılmış gerçeklik gibi uygulamalara girdi sağlıyor. Bu eğilim, mikrofon ve kameranın yakın gelecekte klavye ve farenin önüne geçeceğine işaret ediyor.

Giriş aygıtlarının hızları birbirinden çok farklıdır. Bir klavye saniyede birkaç bayt üretirken, bir mikrofon saniyede on binlerce ses örneği (\emph{sample}) üretir; örneğin 16 bit, 48\,kHz örnekleme ile saniyede yaklaşık 96\,KB. Yüksek çözünürlüklü bir kamera ise saniyede gigabaytlar düzeyinde ham veri üretebilir. Bu büyük hız farkı, giriş/çıkış alt sisteminin tasarımında can alıcı bir etkendir ve Bölüm~\ref{chap:giris-cikis}'da ayrıntılı olarak ele alınacaktır.

\subsubsection*{Çıkış aygıtları}
İşlenen verilerin kullanıcıya ya da başka bir sisteme iletilmesini sağlayan birimlere \emph{çıkış aygıtı}\index{çıkış aygıtı} denir. Monitör (LCD, OLED) ve yazıcı en bilinen görsel çıkış aygıtlarıdır. Bir beneğin (\emph{pixel}) rengini göstermek için kırmızı, yeşil ve mavi kanalların her birine 8 bit ayrılır; dolayısıyla bir beneğin rengi 24 bitle ifade edilir. 1920$\times$1080 çözünürlükteki bir ekranda yaklaşık 2 milyon benek bulunur ve bunların tamamının saniyede en az 60 kez yenilenmesi gerekir; bu da grafik alt sisteminden yüksek bir veri aktarım hızı talep eder.

Yapay zekânın yaygınlaşmasıyla birlikte \emph{hoparlör}\index{hoparlör} de en az ekran kadar önemli bir çıkış aygıtı hâline gelmiştir. Yapay zekâ destekli konuşma sentezi (\emph{text-to-speech}) teknolojileri artık insandan ayırt edilemeyecek kadar doğal ses üretebiliyor. Akıllı hoparlörler, araç içi asistanlar ve giyilebilir aygıtlar gibi ekranı küçük ya da hiç olmayan sistemlerde ses, birincil çıkış kanalı olarak öne çıkıyor. Böylece giriş tarafında mikrofon ve kamera, çıkış tarafında ise hoparlör ve ekran birlikte çalışarak insanla bilgisayar arasında çok kipli (\emph{multimodal})\index{çok kipli etkileşim} bir etkileşim ortamı oluşturuyor.

\subsubsection*{Ana kart}
Bilgisayarın bileşenlerini birbirine bağlayan ve uyumlu çalışmalarını sağlayan donanıma \emph{ana kart}\index{ana kart} denir. Ana kart üzerinde işlemci yuvası, bellek yuvaları, giriş/çıkış denetleyicileri, genişleme yuvaları ve çeşitli bağlantı arabirimleri bulunur. Bileşenler arasındaki veri akışı, ana kart üzerindeki \emph{yol}\index{yol} (\emph{bus}) adı verilen iletişim hatları aracılığıyla sağlanır. Çağdaş ana kartlarda yüksek hızlı bileşenler için ayrı bir kuzey köprüsü, düşük hızlı bileşenler için ise güney köprüsü kullanılırdı; günümüzde bu işlevlerin çoğu doğrudan işlemci içine entegre edilmiştir.

\subsubsection*{Bilgisayar türleri}
Yukarıda tanımlanan beş temel bileşen (işlemci, bellek, giriş aygıtları, çıkış aygıtları ve ana kart) farklı biçim ve boyutlarda bir araya gelebilir. Masaüstü bilgisayarlar bu bileşenleri ayrı parçalar hâlinde büyük bir kasada barındırırken, dizüstü bilgisayarlar aynı bileşenleri katlanabilir ince bir gövdeye sığdırır. Akıllı telefonlarda ise işlemci, bellek ve birçok denetleyici tek bir yonga üzerinde (\emph{System on Chip}, SoC\index{SoC}) bütünleştirilir. Sunucularda yüksek başarım ve güvenilirlik ön plandadır; gömülü sistemlerde ise düşük güç tüketimi ve küçük boyut öne çıkar. Şekil~\ref{fig:bilgisayar-turleri}'te bu farklı bilgisayar türlerinin örnekleri gösterilmektedir.

\begin{figure}[H]
\centering
\begin{tabular}{@{}c@{\hspace{2mm}}c@{\hspace{2mm}}c@{\hspace{2mm}}c@{\hspace{2mm}}c@{}}
\includegraphics[width=0.18\textwidth]{gorseller/desktop.png} &
\includegraphics[width=0.18\textwidth]{gorseller/laptop.png} &
\includegraphics[width=0.18\textwidth]{gorseller/smartphone.png} &
\includegraphics[width=0.18\textwidth]{gorseller/server.png} &
\includegraphics[width=0.18\textwidth]{gorseller/embedded.png} \\[2pt]
{\small\textbf{Masaüstü}} &
{\small\textbf{Dizüstü}} &
{\small\textbf{Akıllı Telefon}} &
{\small\textbf{Sunucu}} &
{\small\textbf{Gömülü Sistem}} \\[-1pt]
{\scriptsize\itshape Desktop} &
{\scriptsize\itshape Laptop} &
{\scriptsize\itshape Smartphone} &
{\scriptsize\itshape Server} &
{\scriptsize\itshape Embedded}
\end{tabular}
\caption{Farklı bilgisayar türleri}
\label{fig:bilgisayar-turleri}
\end{figure}


\subsection{Gerçek Dünyayla Etkileşen Bilgisayarlar}
\label{sec:fiziksel-bilgisayarlar}

Önceki alt bölümde bilgisayarın beş temel bileşenini tanıdık ve bunların masaüstünden akıllı telefona kadar çeşitli biçimlerde bir araya geldiğini gördük. Bu bileşenler yalnızca insan kullanıcıya hizmet eden aygıtlarla sınırlı değil: günümüzde giderek artan sayıda bilgisayar, gerçek dünyayı doğrudan algılayıp ona müdahale edebiliyor. İnsansı robotlar\index{insansı robot}, sürücüsüz araçlar (otonom)\index{sürücüsüz araç}\index{otonom araç}, insansız hava araçları (İHA)\index{İHA}\index{drone} ve endüstriyel robot kolları\index{endüstriyel robot} bunun en çarpıcı örnekleridir (Şekil~\ref{fig:fiziksel-bilgisayarlar}).

\begin{figure}[H]
\centering
\begin{tabular}{@{}c@{\hspace{2mm}}c@{\hspace{2mm}}c@{\hspace{2mm}}c@{}}
\includegraphics[width=0.23\textwidth]{gorseller/humanoid-robot.jpg} &
\includegraphics[width=0.23\textwidth]{gorseller/surucusuz-arac.jpg} &
\includegraphics[width=0.23\textwidth]{gorseller/iha-drone.jpg} &
\includegraphics[width=0.23\textwidth]{gorseller/endustriyel-robot.jpg} \\[2pt]
{\small\textbf{İnsansı Robot}} &
{\small\textbf{Sürücüsüz Araç}} &
{\small\textbf{İHA (Drone)}} &
{\small\textbf{Endüstriyel Robot}} \\[-1pt]
{\scriptsize\itshape Humanoid Robot} &
{\scriptsize\itshape Autonomous Vehicle} &
{\scriptsize\itshape UAV / Drone} &
{\scriptsize\itshape Industrial Robot}
\end{tabular}
\caption{Gerçek dünyayla etkileşen bilgisayar örnekleri}
\label{fig:fiziksel-bilgisayarlar}
\end{figure}

\subsubsection*{Beş bileşen aynı, giriş/çıkış farklı}
Sayılan sistemlerin her biri özünde bir bilgisayardır ve Şekil~\ref{fig:bes-bilesen}'te gösterilen beş bileşeni eksiksiz barındırır. Onları farklı kılan şey giriş ve çıkış aygıtlarıdır. Geleneksel bir bilgisayarda giriş aygıtı klavye ve fare, çıkış aygıtı ise monitör ve hoparlördür. Oysa bir insansı robotta giriş aygıtları stereo kameralar (gözler), mikrofonlar (kulaklar), LiDAR\index{LiDAR}, kuvvet/tork algılayıcıları ve eklem konum algılayıcılarıdır. Çıkış aygıtları ise bacak ve kol eklemlerini hareket ettiren elektrik motorları ve sürücüleridir (\emph{actuator})\index{eyleyici}. Bir sürücüsüz araçta da durum farklı değildir: kameralar, radar\index{radar}, LiDAR ve ultrasonik algılayıcılar\index{ultrasonik algılayıcı} giriş sağlarken; direksiyon motoru, gaz ve fren denetleyicileri çıkış aygıtlarını oluşturur.

\subsubsection*{Gerçek zamanlılık gereksinimi}
Çevreyle etkileşen sistemlerde en önemli mimari kısıt \emph{gerçek zamanlılık}\index{gerçek zamanlı sistem} (\emph{real-time}) gereksinimidir. Bir masaüstü bilgisayarda fare tıklamasına verilen yanıtın birkaç milisaniye gecikmesi fark edilmeyebilir; ancak saatte 100\,km hızla giden bir sürücüsüz aracın çevreyi algılayıp karar vermesi için yalnızca birkaç on milisaniye süresi vardır. Bu nedenle bu tür sistemlerdeki işlemciler, saniyede milyonlarca algılayıcı verisini birleştirmeli (\emph{sensor fusion})\index{algılayıcı birleştirme}, yapay zekâ modelleriyle çevreyi yorumlamalı ve eyleyicilere komut göndermeli; tüm bunları kesin zaman sınırları içinde yapmalıdır. Bu gereksinim, özel yapay zekâ hızlandırıcılarının (GPU, NPU, TPU gibi) bu platformlara entegre edilmesinin temel nedenidir ve Bölüm~\ref{chap:gpu-hizlandiricilar}'de ayrıntılı olarak ele alınacaktır.

\subsubsection*{Bilgisayar mimarisi açısından önemi}
Gerçek dünyayla etkileşen bilgisayarlar, mimari tasarımı birçok yönden zorlamaktadır: algılayıcılardan gelen yüksek bant genişlikli veri akışları verimli giriş/çıkış alt sistemleri gerektirir; gerçek zamanlı yapay zekâ çıkarımı güçlü koşut işlem kapasitesi talep eder; enerji kaynağının sınırlı olması (pil) düşük güç tüketimini zorunlu kılar; güvenlik açısından can alıcı kararların deterministik zamanda tamamlanması gerekir. Bu gereksinimler, bilgisayar mimarisinin geleceğini şekillendiren en önemli itici güçlerden biri hâline gelmiştir.


\section{İşlemci ve Yonga Üretimi}
\label{sec:islemci}

Bilgisayarın beş ana bileşeni arasında bilgisayarın kalbi sayılan \textbf{işlemci}\index{işlemci} özel bir yere sahiptir. İşlemci, bellekten okunan buyrukları yorumlayıp yürüten, aritmetik ve mantık işlemlerini gerçekleştiren ve tüm sistemi denetleyen merkezi birimdir. Bir bilgisayarın ne kadar hızlı çalışacağı, ne kadar enerji tüketeceği ve hangi iş yüklerini kaldırabileceği büyük ölçüde işlemcisinin tasarımına bağlıdır. Bu nedenle bilgisayar mimarisinin en temel konusu işlemci tasarımıdır ve bu kitabın büyük bölümü işlemcinin iç yapısını, çalışma ilkelerini ve başarımını artırma yöntemlerini ele almaktadır. Şimdi bu can alıcı bileşenin nasıl üretildiğini inceleyelim.

\subsection{İşlemcinin Yapılışı}
\label{sec:islemcinin-yapilisi}
İlk olarak vakum tüpleri\index{vakum tüpü} kullanılarak yapılan işlemci daha sonra transistörlerle yapılmaya başlandı. Yapısı oldukça karmaşık olan işlemcinin yapımında milyonlarca, hatta milyarlarca transistör kullanılabiliyor. Örneğin Apple M4 yongasında yaklaşık 28 milyar, AMD EPYC işlemcisinde ise 50 milyardan fazla transistör bulunur. Bu denli büyük bir karmaşıklığı yönetmek için işlemci tasarımının değişik aşamaları birbirinden soyutlanır. \emph{Soyutlama}\index{soyutlama}, her katmandaki ayrıntıları bir üst katmandan gizleyerek karmaşıklığın denetlenmesini kolaylaştırır. Örneğin bir derleyici, işlemcinin içindeki transistör bağlantılarını bilmek zorunda değildir; yalnızca buyruk kümesi düzeyindeki soyutlamayla çalışır. Benzer biçimde bir işlemci tasarımcısı, transistörlerin fiziksel üretim sürecini değil, mantık kapısı düzeyindeki davranışları temel alır. Şekil~\ref{fig:soyutlama-katmanlari}'da bilgisayar sistemindeki soyutlama katmanları gösterilmektedir.

Bu katmanlı yapı yalnızca karmaşıklığı yönetmekle kalmaz, aynı zamanda farklı uzmanlık alanlarının bağımsız çalışmasını sağlar. Bir yazılım geliştiricisi işlemcinin iç yapısını bilmeden program yazabilir; bir donanım mühendisi de üzerinde hangi yazılımların çalışacağını önceden kestirmek zorunda kalmadan işlemci tasarlayabilir. Bu bağımsızlık, bilgisayar endüstrisinin hızlı gelişiminin temel taşlarından biridir. Örneğin aynı x86 buyruk kümesi mimarisi, 1978'deki 8086 işlemcisinden günümüzdeki milyarlarca transistörlük işlemcilere kadar korunmuş; yazılımlar her yeni nesil donanımda yeniden yazılmadan çalışmaya devam etmiştir. Soyutlama katmanları arasındaki bu \emph{arayüz sözleşmesi}, ilerleyen bölümlerde ayrıntılı olarak ele alınacak olan buyruk kümesi mimarisi kavramının ta kendisidir.

\begin{figure}[H]
\centering
\begin{tikzpicture}[>=Stealth, scale=0.88, every node/.style={transform shape}]
    % Katman stilleri
    \tikzstyle{katman}=[draw, thick, rounded corners=4pt, minimum width=9cm,
        minimum height=0.9cm, font=\small\bfseries, align=center,
        drop shadow={shadow xshift=0.3mm, shadow yshift=-0.3mm, opacity=0.12}]

    % Katmanlar (alttan üste)
    \node[katman, fill=gray!15] (fizik) at (0,0)
        {Fizik / Elektron};
    \node[katman, fill=gray!25] (aygit) at (0,1.1)
        {Aygıt (Transistör)};
    \node[katman, fill=blokmavi!30] (devre) at (0,2.2)
        {Analog / Sayısal Devre};
    \node[katman, fill=blokmavi!50] (mantik) at (0,3.3)
        {Mantık Kapıları};
    \node[katman, fill=blokyesil!40] (mikro) at (0,4.4)
        {Mikromimari (Veri Yolu + Denetim)};
    \node[katman, fill=blokyesil!60] (bkm) at (0,5.5)
        {Buyruk Kümesi Mimarisi (BKM)};
    \node[katman, fill=bloksari!50] (isletim) at (0,6.6)
        {İşletim Sistemi};
    \node[katman, fill=bloksari!70] (uygulama) at (0,7.7)
        {Uygulama Yazılımı};

    % Sol tarafta parantezler ve etiketler (sola açılan, katmanları gruplayan)
    \draw[decorate, decoration={brace, amplitude=8pt}, thick, gray!70]
        (-5.2,-.45) -- (-5.2,2.65) node[midway, left=12pt, font=\normalsize\bfseries, align=right, text=gray!70] {Donanım};
    \draw[decorate, decoration={brace, amplitude=8pt}, thick, gray!70]
        (-5.2,2.85) -- (-5.2,5.95) node[midway, left=12pt, font=\normalsize\bfseries, align=right, text=gray!70] {Mimari};
    \draw[decorate, decoration={brace, amplitude=8pt}, thick, gray!70]
        (-5.2,6.15) -- (-5.2,8.15) node[midway, left=12pt, font=\normalsize\bfseries, align=right, text=gray!70] {Yazılım};

    % Sağ tarafta ok ve etiket (yalnızca yukarı yönlü)
    \draw[->, line width=2pt, sinyalkirmizi!70] (5.2,0) -- (5.2,7.7)
        node[midway, right=6pt, font=\normalsize\bfseries, align=left, text=sinyalkirmizi!70] {Artan\\soyutlama\\düzeyi};
\end{tikzpicture}
\caption[Bilgisayar sistemindeki soyutlama katmanları]{Bilgisayar sistemindeki soyutlama katmanları. Her katman yalnızca bir altındaki katmanın sunduğu arayüzü kullanır.}
\label{fig:soyutlama-katmanlari}
\end{figure}

\subsection{Tümleşik Devreler}
\label{sec:tumlesik-devreler}
Bölüm~\ref{sec:bilgisayar-gelisimi}'de anlatıldığı gibi, tümleşik devrelerin\index{tümleşik devre} keşfi bilgisayar tarihinin dönüm noktalarından biridir. Tümleşik devredeki ana fikir, transistörler, dirençler, kapasitörler ve bunlar arasındaki bağlantıların tamamının aynı yarıiletken malzeme üzerine tek bir süreçle üretilmesidir. Böylece ayrık bileşenlerden oluşan devrelere kıyasla çok daha küçük, hızlı ve güvenilir yongalar\index{yonga} elde ediliyor.

Tümleşik devreler, içerdikleri transistör sayısına göre sınıflandırılır: birkaç düzine transistör barındıran \emph{SSI}\index{SSI} (Small Scale Integration), birkaç yüz transistörlük \emph{MSI}\index{MSI}, on binlerce transistörlük \emph{LSI}\index{LSI} ve yüz binlerin üzerine çıkan \emph{VLSI}\index{VLSI} (Very Large Scale Integration: Çok Büyük Ölçekli Tümleştirme). Günümüzdeki işlemciler milyarlarca transistör içermekte olup artık ``ultra-VLSI'' olarak nitelendirilebilecek bir ölçektedir. Transistör boyutlarının küçülmesi, aynı alana daha fazla transistör sığdırılmasını sağlamış; ancak beraberinde güç tüketimi, ısı yönetimi ve üretim maliyeti gibi yeni mühendislik zorlukları da getirmiştir.

Tümleşik devrelerin etkisi yalnızca boyut ve hız kazanımıyla sınırlı kalmamıştır. Ayrık bileşenlerle kurulan devrelerde her lehim noktası olası bir arıza kaynağıyken, tümleşik devrede bağlantılar fabrikada otomatik olarak oluşturulduğundan güvenilirlik büyük ölçüde artmıştır. Üstelik üretim miktarı arttıkça birim maliyet düşer: günümüzde birkaç milyar transistör içeren bir işlemci yongası, enflasyon etkisi çıkarıldığında 1960'ların birkaç yüz transistörlük devresinden daha ucuza üretilebilmektedir. Bu maliyet--başarım ilişkisi, bilgisayarların yalnızca laboratuvar ve iş dünyasında değil, cep telefonlarından bulaşık makinelerine kadar günlük yaşamın her köşesinde yer almasını olanaklı kılmıştır.

\subsection{Yonga Üretim Süreci}
\label{sec:yonga-uretim-sureci}

\begin{figure}[H]
\centering
\begin{tikzpicture}[>=Stealth, scale=0.92, every node/.style={transform shape},
    adim/.style={draw, thick, rounded corners=4pt, minimum width=2.8cm,
                 minimum height=2.6cm, font=\footnotesize\bfseries, align=center,
                 inner sep=3pt}]

    % ============ ÜST SIRA (soldan sağa) ============

    % 1. Tasarım — devre şeması ikonu
    \node[adim, fill=blokmavi] (tasarim) at (0,0) {};
    \begin{scope}[shift={(tasarim.center)}, yshift=0.2cm, scale=1.3]
        \draw[thin, gray!60] (-0.6,-0.4) grid[step=0.3] (0.6,0.4);
        \draw[thick, sinyalmavi] (-0.6,-0.1) -- (-0.3,-0.1) -- (-0.3,0.2) -- (0.3,0.2) -- (0.3,-0.1) -- (0.6,-0.1);
        \fill[sinyalmavi] (-0.3,-0.1) circle (1.5pt);
        \fill[sinyalmavi] (0.3,0.2) circle (1.5pt);
    \end{scope}

    % 2. Maske Hazırlama — katmanlı dikdörtgenler
    \node[adim, fill=blokmavi] (maske) at (3.5,0) {};
    \begin{scope}[shift={(maske.center)}, yshift=0.2cm, scale=1.3]
        \fill[sinyalmavi!30] (-0.5,-0.25) rectangle (0.5,0.35);
        \fill[sinyalmavi!50, opacity=0.7] (-0.35,-0.15) rectangle (0.35,0.25);
        \draw[thick, sinyalmavi] (-0.2,0.05) rectangle (0.2,0.15);
        \draw[thick, sinyalmavi] (-0.1,-0.1) rectangle (0.4,0.0);
    \end{scope}

    % 3. Fotolitografi — UV ışınları
    \node[adim, fill=blokyesil] (isik) at (7,0) {};
    \begin{scope}[shift={(isik.center)}, yshift=0.15cm, scale=1.3]
        \fill[sinyalyesil!70!yellow] (0,0.45) circle (0.12);
        \foreach \dx in {-0.35, -0.15, 0.05, 0.25} {
            \draw[->, sinyalyesil!70!yellow, thick] (\dx+0.05,0.3) -- (\dx+0.05,-0.2);
        }
        \fill[gray!30] (-0.5,-0.25) rectangle (0.5,-0.15);
        \draw[thick] (-0.5,-0.25) rectangle (0.5,-0.15);
    \end{scope}

    % 4. Dağlama — dalgalı kimyasal çizgiler
    \node[adim, fill=blokyesil] (daglama) at (10.5,0) {};
    \begin{scope}[shift={(daglama.center)}, yshift=0.15cm, scale=1.3]
        \fill[gray!30] (-0.5,-0.3) rectangle (0.5,0.0);
        \fill[white] (-0.2,-0.05) rectangle (0.2,0.05);
        \draw[thick] (-0.5,-0.3) rectangle (0.5,0.0);
        \draw[thick] (-0.2,-0.05) -- (-0.2,0.0);
        \draw[thick] (0.2,-0.05) -- (0.2,0.0);
        \draw[sinyalkirmizi, thick, decorate, decoration={snake, amplitude=1pt, segment length=4pt}]
            (-0.4,0.2) -- (0.4,0.2);
        \draw[sinyalkirmizi, thick, decorate, decoration={snake, amplitude=1pt, segment length=4pt}]
            (-0.35,0.35) -- (0.35,0.35);
    \end{scope}

    % ============ ALT SIRA (sağdan sola) ============

    % 5. Kaplama — aşağı oklar + katman
    \node[adim, fill=bloksari] (kaplama) at (10.5,-3.8) {};
    \begin{scope}[shift={(kaplama.center)}, yshift=0.15cm, scale=1.3]
        \foreach \dx in {-0.3, -0.1, 0.1, 0.3} {
            \draw[->, thick, sinyalmavi!70!black] (\dx,0.35) -- (\dx,-0.05);
        }
        \fill[gray!25] (-0.5,-0.3) rectangle (0.5,-0.2);
        \fill[sinyalmavi!25] (-0.5,-0.2) rectangle (0.5,-0.1);
        \draw[thick] (-0.5,-0.3) rectangle (0.5,-0.1);
    \end{scope}

    % 6. Test — onay işareti
    \node[adim, fill=bloksari] (test) at (7,-3.8) {};
    \begin{scope}[shift={(test.center)}, yshift=0.1cm, scale=1.3]
        \draw[very thick, sinyalyesil] (-0.2,0.0) -- (-0.05,-0.2) -- (0.3,0.3);
        \draw[thick, gray!60] (0,0) circle (0.45);
    \end{scope}

    % 7. Kesim — plaka + testere çizgisi
    \node[adim, fill=blokkirmizi] (kesim) at (3.5,-3.8) {};
    \begin{scope}[shift={(kesim.center)}, yshift=0.1cm, scale=1.3]
        \draw[thick, fill=gray!15] (0,0) circle (0.4);
        \draw[densely dashed, sinyalkirmizi, thin] (-0.35,0.13) -- (0.35,0.13);
        \draw[densely dashed, sinyalkirmizi, thin] (-0.35,-0.13) -- (0.35,-0.13);
        \draw[densely dashed, sinyalkirmizi, thin] (-0.13,-0.35) -- (-0.13,0.35);
        \draw[densely dashed, sinyalkirmizi, thin] (0.13,-0.35) -- (0.13,0.35);
    \end{scope}

    % 8. Paketleme — yonga + pinler
    \node[adim, fill=blokkirmizi-acik] (paket) at (0,-3.8) {};
    \begin{scope}[shift={(paket.center)}, yshift=0.1cm, scale=1.3]
        \fill[gray!40] (-0.3,-0.2) rectangle (0.3,0.2);
        \draw[thick] (-0.3,-0.2) rectangle (0.3,0.2);
        \fill[sinyalmavi!30] (-0.12,-0.08) rectangle (0.12,0.08);
        \foreach \px in {-0.2, -0.05, 0.1, 0.25} {
            \draw[thick] (\px,0.2) -- (\px,0.35);
        }
        \foreach \px in {-0.2, -0.05, 0.1, 0.25} {
            \draw[thick] (\px,-0.2) -- (\px,-0.35);
        }
        \foreach \py in {-0.1, 0.1} {
            \draw[thick] (-0.3,\py) -- (-0.45,\py);
        }
        \foreach \py in {-0.1, 0.1} {
            \draw[thick] (0.3,\py) -- (0.45,\py);
        }
    \end{scope}

    % ============ OKLAR (kesikli çizgiden önce çiz) ============

    % Üst sıra (soldan sağa)
    \draw[->, very thick, sinyalmavi!70!black] (tasarim.east) -- (maske.west);
    \draw[->, very thick, sinyalmavi!70!black] (maske.east) -- (isik.west);
    \draw[->, very thick, sinyalmavi!70!black] (isik.east) -- (daglama.west);

    % Aşağı dönüş (sağ taraftan geçir)
    \draw[->, very thick, sinyalmavi!70!black]
        (daglama.east) -- ++(0.8,0) |- (kaplama.east);

    % Alt sıra (sağdan sola)
    \draw[->, very thick, sinyalmavi!70!black] (kaplama.west) -- (test.east);
    \draw[->, very thick, sinyalmavi!70!black] (test.west) -- (kesim.east);
    \draw[->, very thick, sinyalmavi!70!black] (kesim.west) -- (paket.east);

    % ============ KESİKLİ ÇERÇEVE (okların üstüne, etiketlerin altına) ============
    \draw[dashed, thick, gray!50, rounded corners=5pt]
        ([xshift=-0.5cm, yshift=0.7cm]isik.north west) rectangle
        ([xshift=1.15cm, yshift=-0.7cm]kaplama.south east);

    % ============ ETİKETLER VE NUMARALAR (en üste — kesikli çizginin üstünde) ============
    \node[font=\footnotesize\bfseries, anchor=north] at ([yshift=-0.05cm]tasarim.south) {Tasarım};
    \node[font=\footnotesize\bfseries, fill=white, inner sep=1pt, text=gray!50!black] at ([yshift=0.2cm]tasarim.north) {1};

    \node[font=\footnotesize\bfseries, anchor=north] at ([yshift=-0.05cm]maske.south) {Maske};
    \node[font=\footnotesize\bfseries, fill=white, inner sep=1pt, text=gray!50!black] at ([yshift=0.2cm]maske.north) {2};

    \node[font=\footnotesize\bfseries, anchor=north] at ([yshift=-0.05cm]isik.south) {Fotolitografi};
    \node[font=\footnotesize\bfseries, fill=white, inner sep=1pt, text=gray!50!black] at ([yshift=0.2cm]isik.north) {3};

    \node[font=\footnotesize\bfseries, anchor=north] at ([yshift=-0.05cm]daglama.south) {Dağlama};
    \node[font=\footnotesize\bfseries, fill=white, inner sep=1pt, text=gray!50!black] at ([yshift=0.2cm]daglama.north) {4};

    \node[font=\footnotesize\bfseries, anchor=north] at ([yshift=-0.05cm]kaplama.south) {Kaplama};
    \node[font=\footnotesize\bfseries, fill=white, inner sep=1pt, text=gray!50!black] at ([yshift=0.2cm]kaplama.north) {5};

    \node[font=\footnotesize\bfseries, anchor=north] at ([yshift=-0.05cm]test.south) {Test};
    \node[font=\footnotesize\bfseries, fill=white, inner sep=1pt, text=gray!50!black] at ([yshift=0.2cm]test.north) {6};

    \node[font=\footnotesize\bfseries, anchor=north] at ([yshift=-0.05cm]kesim.south) {Kesim};
    \node[font=\footnotesize\bfseries, fill=white, inner sep=1pt, text=gray!50!black] at ([yshift=0.2cm]kesim.north) {7};

    \node[font=\footnotesize\bfseries, anchor=north] at ([yshift=-0.05cm]paket.south) {Paketleme};
    \node[font=\footnotesize\bfseries, fill=white, inner sep=1pt, text=gray!50!black] at ([yshift=0.2cm]paket.north) {8};

    % Tekrarlama notu (en üstte)
    \node[font=\scriptsize\itshape, fill=white, inner sep=2pt, text=gray!40!black, anchor=south] at (8.75,1.8)
        {Bu adımlar 50+ katman için tekrarlanır};

\end{tikzpicture}
\caption[Yonga yapılma süreci]{Yonga yapılma süreci. Çağdaş bir işlemcinin üretiminde fotolitografi--dağlama--kaplama adımları onlarca katman için tekrarlanır.}
\label{fig:yonga-yapim}
\end{figure}

Yapılması planlanan yonga öncelikle bilgisayar üzerinde plan olarak çizilir ve denenir. Işıklı baskı (photolithography) tekniğiyle tüm yongaların basılması için bir kalıp hazırlanır. Silisyum\index{silisyum} silindirden kristal bir testereyle kesilen ince katmanların üzerine katmanları birbirinden yalıtacak şekilde silisyum oksit kaplanır ve kalıbın ana hatlarının geçilebilmesi için ışığa duyarlı bir tabaka kaplanır. Yonga yapım süreci Şekil~\ref{fig:yonga-yapim}'de gösterilmiştir.

\begin{bilginot}[Silisyum mu, Silikon mu?]
Yarıiletken yongaların temel malzemesi olan kimyasal element İngilizcede \emph{silicon} (Si, atom numarası~14), Türkçede ise \textbf{silisyum} olarak adlandırılır. Günlük dilde sıkça karşılaşılan ``silikon'' sözcüğü ise İngilizce \emph{silicone} kelimesinin karşılığıdır ve silisyum ile oksijenden türetilmiş sentetik bir polimeri ifade eder (mutfak kalıpları, conta malzemeleri vb.). Dolayısıyla yonga üretiminde kullanılan malzeme \emph{silisyum}, ABD'deki ünlü teknoloji vadisinin adı ise doğru çevirisiyle \emph{Silisyum Vadisi}'dir.
\end{bilginot}

Işıklı baskı işlemi son derece hassas bir ortamda gerçekleştirilir. Üretim tesislerindeki \emph{temiz oda}\index{temiz oda} (\emph{cleanroom}) koşullarında havadaki toz partikül sayısı dış ortamdan on binlerce kat azdır; çünkü nanometre ölçeğindeki devrelerde tek bir toz tanesi bile yongayı kullanılamaz hâle getirebilir. Bu nedenle yonga üretim tesisleri (\emph{fab}\index{fab}), dünyada inşa edilen en pahalı endüstriyel yapılar arasında yer alır; günümüzde ileri düğüm teknolojisine sahip bir fabrikanın maliyeti 20 milyar doları aşabilmektedir.

Üretim sürecinin en kritik adımlarından biri, plaka üzerindeki deseni oluşturan fotolitografi basamağıdır. Günümüzde 5\,nm ve altı düğüm teknolojilerinde \emph{aşırı morötesi ışıklı baskı}\index{EUV} (EUV: Extreme Ultraviolet Lithography) kullanılmaktadır. EUV, 13{,}5\,nm dalga boylu ışık kullanarak geleneksel 193\,nm ArF lazerlerinin çözünürlük sınırlarını aşar ve daha ince hatların çizilmesini olanaklı kılar. Ancak EUV makineleri, içinde plazma kaynağı, vakum odaları ve çok katmanlı aynalar barındıran son derece karmaşık sistemlerdir; dünyada yalnızca tek bir firma (ASML) bu makineleri üretebilmektedir.

\begin{figure}[H]
\centering
\begin{tikzpicture}[>=Stealth]
    \def\R{3.8}        % plaka yarıçapı
    \def\ks{0.7}       % kırmık kenar uzunluğu
    \def\kg{0.08}      % kırmıklar arası boşluk (kesim payı)
    \def\adim{0.78}    % ızgara adımı = ks + kg
    \def\ha{0.35}      % yonga yarı kenarı = ks/2

    % --- 1. Silisyum plaka (arka plan) ---
    \fill[gray!12] (0,0) circle (\R);
    \begin{scope}
        \clip (0,0) circle (\R);
        \fill[white] (-0.4,{-\R-0.1}) rectangle (0.4,{-\R+0.25});
    \end{scope}

    % --- 2. Kırmık ızgarası (sadece dört köşesi de dairenin içinde olanlar) ---
    \foreach \i in {-5,-4,...,5} {
        \foreach \j in {-5,-4,...,5} {
            \pgfmathsetmacro{\cx}{\i*\adim}
            \pgfmathsetmacro{\cy}{\j*\adim}
            \pgfmathsetmacro{\da}{sqrt((\cx-\ha)^2+(\cy-\ha)^2)}
            \pgfmathsetmacro{\db}{sqrt((\cx+\ha)^2+(\cy-\ha)^2)}
            \pgfmathsetmacro{\dc}{sqrt((\cx+\ha)^2+(\cy+\ha)^2)}
            \pgfmathsetmacro{\dd}{sqrt((\cx-\ha)^2+(\cy+\ha)^2)}
            \pgfmathsetmacro{\dmax}{max(\da,max(\db,max(\dc,\dd)))}
            \pgfmathparse{\dmax < (\R-0.15) ? 1 : 0}
            \ifnum\pgfmathresult=1
                \fill[sinyalmavi!12] ({\cx-\ha},{\cy-\ha}) rectangle ({\cx+\ha},{\cy+\ha});
                \draw[gray!40, thin] ({\cx-\ha},{\cy-\ha}) rectangle ({\cx+\ha},{\cy+\ha});
            \fi
        }
    }

    % --- 3. Kenar yongaları (kısmi, verim kaybı göstermek için) ---
    \begin{scope}
        \clip (0,0) circle ({\R-0.05});
        \foreach \i in {-5,-4,...,5} {
            \foreach \j in {-5,-4,...,5} {
                \pgfmathsetmacro{\cx}{\i*\adim}
                \pgfmathsetmacro{\cy}{\j*\adim}
                \pgfmathsetmacro{\da}{sqrt((\cx-\ha)^2+(\cy-\ha)^2)}
                \pgfmathsetmacro{\db}{sqrt((\cx+\ha)^2+(\cy-\ha)^2)}
                \pgfmathsetmacro{\dc}{sqrt((\cx+\ha)^2+(\cy+\ha)^2)}
                \pgfmathsetmacro{\dd}{sqrt((\cx-\ha)^2+(\cy+\ha)^2)}
                \pgfmathsetmacro{\dmax}{max(\da,max(\db,max(\dc,\dd)))}
                \pgfmathsetmacro{\dmin}{sqrt(\cx*\cx+\cy*\cy)}
                \pgfmathparse{(\dmax >= (\R-0.15)) && (\dmin < (\R-0.3)) ? 1 : 0}
                \ifnum\pgfmathresult=1
                    \fill[gray!8] ({\cx-\ha},{\cy-\ha}) rectangle ({\cx+\ha},{\cy+\ha});
                    \draw[gray!25, thin, densely dashed] ({\cx-\ha},{\cy-\ha}) rectangle ({\cx+\ha},{\cy+\ha});
                \fi
            }
        }
    \end{scope}

    % --- 4. Plaka konturu (üstte, temiz kenar) ---
    \draw[thick, gray!70] (0,0) circle (\R);
    % Çentik (notch)
    \draw[thick, gray!70] ({-0.35},{-sqrt(\R*\R-0.35*0.35)})
        -- (0,{-\R+0.2})
        -- ({0.35},{-sqrt(\R*\R-0.35*0.35)});

    % --- 5. Bir yongayı vurgula ---
    \draw[very thick, sinyalyesil, fill=sinyalyesil!15]
        ({-\ha},{\adim-\ha}) rectangle ({\ha},{\adim+\ha});

    % --- 6. Etiketler (sağ tarafta, düzenli) ---
    % Kırmık etiketi
    \draw[->, thick, sinyalyesil!70!black] (4.8,2.8) -- (0.45,1.1);
    \node[font=\footnotesize, anchor=west] at (4.9,2.8) {Kırmık (\textit{die})};

    % Hatalı kırmık etiketi (kenar)
    \draw[->, thick, gray!50] (4.8,1.2) -- (3.12,1.56);
    \node[font=\footnotesize, text=gray!60!black, anchor=west] at (4.9,1.4) {Kullanılamayan};
    \node[font=\footnotesize, text=gray!60!black, anchor=west] at (4.9,1.0) {kenar kırmık};

    % Çentik etiketi
    \draw[->, thick, gray!50] (1.8,-4.3) -- (0.15,{-\R+0.22});
    \node[font=\footnotesize, text=gray!60!black, anchor=west] at (1.9,-4.3) {Çentik (\textit{notch})};

    % --- 7. Alt bilgiler ---
    \node[font=\footnotesize\itshape, text=gray!60!black] at (0,-5.0)
        {Silisyum plaka (\textit{wafer})};

    % Çap göstergesi
    \draw[<->, thick, gray!50!black] (-\R,-4.5) -- (\R,-4.5);
    \node[font=\footnotesize, fill=white, inner sep=1pt] at (0,-4.5) {$\varnothing$ 300\,mm};

\end{tikzpicture}
\caption{Bir silisyum plakadan yongaların kesilmesi}
\label{fig:silisyum-plaka}
\end{figure}

\begin{bilginot}[Silisyum plaka neden yuvarlak?]
Öğrenciler sıklıkla ``Plaka neden kare değil, kare olsa daha fazla kırmık sığmaz mı?'' diye sorar. Yanıt üretim sürecinde yatmaktadır. Silisyum plakalar \emph{Czochralski yöntemi}\index{Czochralski yöntemi} ile üretilir: çok saf silisyum bir potada eritilir, dönen bir tohum kristal erimiş silisyuma dokundurularak yavaşça yukarı çekilir. Dönerek çekilme hareketi doğal olarak silindirik bir tek kristal kütle (\emph{ingot}) oluşturur; bu silindirin enine kesitleri daire biçimindedir. Kare plaka üretmek istendiğinde silindirik kütleden kesim sırasında büyük miktarda malzeme israf olurdu. Ayrıca yüksek sıcaklıklı üretim aşamalarında (900--1100\,°C) kare plakanın köşelerinde gerilim yoğunlaşması oluşarak kırılma riski artardı. Yuvarlak geometri ise termal genleşme sırasında gerilimi eşit biçimde dağıtır.
\end{bilginot}


\subsection{Teknoloji Düğümleri ve Gelişimi}
\label{sec:teknoloji-dugumleri}

Bir yonga üretim sürecinin en temel parametresi \emph{teknoloji düğümü}\index{teknoloji düğümü} (\emph{technology node}) olarak adlandırılan ve nanometre (nm) cinsinden ifade edilen boyuttur. Tarihsel olarak bu değer transistörün kapı uzunluğuna karşılık gelirdi; ancak günümüzde gerçek fiziksel boyutlarla doğrudan ilişkisi zayıflamış ve daha çok bir nesil göstergesi hâline gelmiştir. Her yeni teknoloji düğümüne geçişte transistör yoğunluğu yaklaşık iki katına çıkar, bu da aynı alana daha fazla transistör sığdırılmasını sağlar.

\begin{table}[ht!]
\centering
\caption[Teknoloji düğümlerinin gelişimi]{Yıllar içinde teknoloji düğümlerinin gelişimi. Transistör yoğunlukları yaklaşık değerler olup üretim tesisine ve tasarıma göre değişir.}
\label{tab:teknoloji-dugumleri}
\footnotesize
\begin{tabular}{rrrll}
\toprule
\textbf{Düğüm} & \textbf{Yıl} & \textbf{Yoğunluk (MTr/mm$^2$)} & \textbf{Öncü Üretici} & \textbf{Örnek Yonga} \\
\midrule
90\,nm   & 2004 & ${\sim}$15   & Intel     & Pentium 4 Prescott \\
65\,nm   & 2006 & ${\sim}$30   & Intel     & Core 2 Duo \\
45\,nm   & 2008 & ${\sim}$60   & Intel     & Core i7 (Nehalem) \\
32\,nm   & 2010 & ${\sim}$90   & Intel     & Sandy Bridge \\
22\,nm   & 2012 & ${\sim}$160  & Intel     & Ivy Bridge (FinFET) \\
14\,nm   & 2014 & ${\sim}$310  & Intel, Samsung & Skylake, Exynos 7 \\
10\,nm   & 2017 & ${\sim}$530  & TSMC      & Apple A11 \\
7\,nm    & 2018 & ${\sim}$910  & TSMC      & Apple A12, AMD Zen 2 \\
5\,nm    & 2020 & ${\sim}$1\,710 & TSMC    & Apple M1, A14 Bionic \\
3\,nm    & 2022 & ${\sim}$3\,000 & TSMC, Samsung & Apple M3, A17 Pro \\
\bottomrule
\end{tabular}
\end{table}

Tablo~\ref{tab:teknoloji-dugumleri}'den görüldüğü üzere, 2004'ten 2022'ye kadar transistör yoğunluğu yaklaşık 200 kat artmıştır. Ancak bu kazanımlar bedelsiz değildir: teknoloji düğümü küçüldükçe hem plaka maliyeti hem de \emph{maske seti}\index{maske seti} (\emph{mask set}) maliyeti dramatik biçimde artar. Maske seti, yonga üzerindeki her bir katmanın desenini plakaya aktarmak için kullanılan fotolitografi kalıplarının tamamıdır; çağdaş bir işlemcide 80'den fazla maske katmanı bulunabilir ve her maske yüksek duyarlıkta üretilmek zorundadır. Tablo~\ref{tab:plaka-maliyetleri} bu maliyet artışını göstermektedir.

\begin{table}[ht!]
\centering
\caption[Teknoloji düğümlerine göre üretim maliyetleri]{Farklı teknoloji düğümlerinde 300\,mm silisyum plakanın ve maske setinin yaklaşık maliyeti. Maliyetler yonga üretim tesisine ve sipariş hacmine göre değişkenlik gösterir.}
\label{tab:plaka-maliyetleri}
\footnotesize
\begin{tabular}{lrrl}
\toprule
\textbf{Düğüm} & \textbf{Plaka Maliyeti (\$)} & \textbf{Maske Seti (\$)} & \textbf{Tipik Kullanım Alanı} \\
\midrule
28\,nm    & 3\,000--4\,000     & ${\sim}$5\,M    & IoT, otomotiv, gömülü \\
16/14\,nm & 5\,000--6\,000     & ${\sim}$15\,M   & Orta seviye taşınabilir, ağ \\
7\,nm     & 9\,000--10\,000    & ${\sim}$30\,M   & Üst seviye taşınabilir, GPU \\
5\,nm     & 16\,000--17\,000   & ${\sim}$50\,M   & En üst düzey işlemciler \\
3\,nm     & 20\,000+           & ${\sim}$75\,M   & Son nesil yongalar \\
\bottomrule
\multicolumn{4}{l}{\scriptsize M\,=\,milyon; maliyetler tahmini olup üretim tesisine ve hacme göre değişir.}
\end{tabular}
\end{table}

3\,nm düğümünde bir maske setinin maliyeti 75 milyon doları aşabilmektedir. Bu durum, ileri düğümlerde üretimi yalnızca çok yüksek hacimli ürünler (akıllı telefon işlemcileri, veri merkezi yongaları) için ekonomik kılar. Düşük hacimli uygulamalar (IoT algılayıcıları, otomotiv denetleyicileri gibi) maliyet açısından 28\,nm veya 16\,nm gibi olgun düğümlerde üretilmeye devam eder. Bu ekonomik gerçeklik, farklı uygulama alanlarının farklı teknoloji düğümlerinde kalmasına neden olur ve ``en yeni teknoloji her zaman en iyi seçimdir'' yanılgısını çürütür.


\subsection{Üretim Maliyeti ve Verim}
\label{sec:maliyet-verim}

Yonga üretiminin ekonomisini anlamak, mimari kararlar için can alıcıdır. Büyük alanlı yongalar daha fazla transistör barındırabilir ancak verim düşüşü nedeniyle orantısız biçimde pahalılaşır. Tasarımcılar yonga alanı, üretim maliyeti ve verim arasındaki dengeyi gözetmek zorundadır.

Bir yonganın üretim maliyeti aşağıdaki denklemle gösterilebilir:
\begin{equation}
\text{Yonga maliyeti} = \frac{\text{Plaka maliyeti}}{\text{Plakadaki yonga sayısı} \times \text{Verim}}
\label{eq:yonga-maliyet}
\end{equation}

Plakanın alacağı yonga sayısı yaklaşık olarak aşağıdaki eşitlikle gösterilebilir:
\begin{equation}
\text{Plakadaki yonga sayısı} = \frac{\text{Plaka alanı}}{\text{Yonga alanı}}
\label{eq:yonga-sayisi}
\end{equation}

Verim yaklaşık olarak aşağıdaki denklemle modellenir:
\begin{equation}
\text{Verim} = \left( \frac{1}{1 + \text{Birim alan başına hata sayısı} \times \text{Yonga alanı}} \right)^{N}
\label{eq:verim}
\end{equation}
Bu denklemdeki kesirin paydasındaki değerin üssü ($N$) üretim sürecindeki kritik aşama sayısı ile ilişkilidir.

\subsubsection*{Verim Öğrenme Eğrisi}

Yeni bir teknoloji düğümü üretime girdiğinde verim genellikle düşüktür; tipik olarak \%30--50 arasında başlar. Üretim süreci olgunlaştıkça ve hata kaynakları belirlendikçe verim kademeli olarak \%80--90 ve üzerine çıkar. Bu iyileşme sürecine \emph{verim öğrenme eğrisi}\index{verim öğrenme eğrisi} (\emph{yield learning curve}) denir ve bir yonga üretim tesisinin teknolojik olgunluğunun en önemli göstergesidir.

Verim verisi şirketler tarafından ticari sır olarak korunduğundan kesin rakamlar kamuya açıklanmaz; ancak genel eğilimler endüstri analistleri tarafından raporlanır. Örneğin TSMC'nin 7\,nm süreci (N7) 2018'de üretime girdiğinde verimlerin birkaç çeyrek dönem içinde \%80'in üzerine çıktığı tahmin edilmektedir. Denklem~\ref{eq:yonga-maliyet}'den görüldüğü üzere verim yarıya düşerse yonga maliyeti iki katına çıkar; bu nedenle yonga üreticileri yeni düğümlerde verimi artırmak için büyük mühendislik yatırımları yapar.

\begin{ornek}[1.2]
Bir yonga tasarım şirketi, akıllı telefon işlemcisi için 3\,nm teknolojisini kullanmayı düşünüyor. Yonganın kırmık alanı 100\,mm$^2$ olacak. 300\,mm çaplı bir plakanın maliyeti 20\,000\,\$, maske setinin maliyeti 75\,M\,\$, olgun süreçteki verim \%75 olsun. Bir kırmığın üretim maliyetini ve maske seti payını iki farklı üretim hacmi için hesaplayın: (a)~100 milyon adet, (b)~1 milyon adet.

\textbf{Çözüm:}\\
Plaka alanı: $A_\text{plaka} = \pi \times 150^2 \approx 70\,686$\,mm$^2$.

Plakadaki kırmık sayısı (yaklaşık): $70\,686 \div 100 \approx 706$ kırmık. Kenar kayıpları düşüldüğünde kullanılabilir kırmık sayısı yaklaşık 600'dür.

Verimli kırmık sayısı: $600 \times 0{,}75 = 450$ iyi kırmık/plaka.

Plaka başına kırmık maliyeti (Denklem~\ref{eq:yonga-maliyet}): $20\,000 \div 450 \approx 44{,}4$\,\$/kırmık.

\medskip
\textbf{Maske seti payı:}

\smallskip
\begin{tabular}{@{}lrl@{}}
100 milyon adet: & $75\,000\,000 \div 100\,000\,000 = 0{,}75$\,\$/kırmık & $\Rightarrow$ \textbf{toplam $\approx$ 45\,\$/kırmık} \\[4pt]
1 milyon adet:   & $75\,000\,000 \div 1\,000\,000 = 75$\,\$/kırmık     & $\Rightarrow$ \textbf{toplam $\approx$ 119\,\$/kırmık}
\end{tabular}

\medskip
\textbf{Karşılaştırma:} Aynı yonga 28\,nm'de üretilseydi plaka maliyeti ${\sim}$3\,500\,\$, maske seti ${\sim}$5\,M\,\$, verim ${\sim}$\%90 olacaktı. Plaka başına maliyet ${\sim}$6{,}5\,\$/kırmık, 1 milyon adet için maske payı 5\,\$/kırmık, toplam ${\sim}$11{,}5\,\$/kırmık; bu 3\,nm'dekinin onda biridir. Ancak 28\,nm'de aynı alana sığan transistör sayısı 3\,nm'nin yüzde biri kadar olduğundan, aynı işlevsellik çok daha büyük bir kırmık alanı gerektirir.

Bu örnek, ileri teknoloji düğümlerinin neden yalnızca yüksek hacimli ürünlerde ekonomik olduğunu ve maske seti maliyetinin düşük hacimlerde baskın hâle geldiğini göstermektedir.
\end{ornek}

\begin{ornek}[1.3]
Çapı 300\,mm olan bir silisyum plakanın maliyeti 12\,000\,\$ olsun. Üretilecek yonganın alanı 200\,mm$^2$, birim alan başına hata sayısı 0{,}03\,hata/mm$^2$ ve kritik aşama sayısı $N=12$ ise bir yonganın maliyetini hesaplayın.

\textbf{Çözüm:}\\
Plaka alanı: $A_\text{plaka} = \pi \times 150^2 \approx 70\,686$\,mm$^2$.

Plakadaki yonga sayısı: $\displaystyle\frac{70\,686}{200} \approx 353$ yonga.

Verim: $\displaystyle\left(\frac{1}{1 + 0{,}03 \times 200}\right)^{12} = \left(\frac{1}{7}\right)^{12} \approx 7{,}3 \times 10^{-11}$; bu son derece düşük bir verimdir!

Gerçek üretimde birim alan başına hata sayısı çok daha düşüktür. Eğer hata yoğunluğu 0{,}001\,hata/mm$^2$ olsaydı:
\[
\text{Verim} = \left(\frac{1}{1 + 0{,}001 \times 200}\right)^{12} = \left(\frac{1}{1{,}2}\right)^{12} \approx 0{,}112
\]
Bu durumda yonga maliyeti: $\displaystyle\frac{12\,000}{353 \times 0{,}112} \approx 303$\,\$.

Bu örnek, yonga alanının ve hata yoğunluğunun maliyet üzerindeki dramatik etkisini açıkça gösteriyor.
\end{ornek}

\begin{ornek}[1.4]
\label{orn:verim-hesaplama}
300\,mm çapında bir silisyum plakadan 100\,mm$^2$ büyüklüğünde kırmıklar (\emph{die}) kesilmektedir. Plaka başına hata yoğunluğu $d = 0{,}04$ hata/cm$^2$ ve üretim karmaşıklığı parametresi $\alpha = 4$ ise:
\begin{enumerate}[(a)]
\item Plakadan elde edilen yaklaşık kırmık sayısını hesaplayın.
\item Kırmık verimini hesaplayın.
\item Kırmık boyutu 50\,mm$^2$'ye küçültülürse verim nasıl değişir?
\end{enumerate}

\textbf{Çözüm:}

(a) Plaka alanı $A_{\text{plaka}} = \pi \times (150)^2 \approx 70\,686$\,mm$^2$. Kırmık sayısı yaklaşık:
\[
N_{\text{kırmık}} \approx \frac{A_{\text{plaka}}}{A_{\text{kırmık}}} = \frac{70\,686}{100} \approx 707 \text{ kırmık}
\]
(Kenar kayıpları düşüldüğünde uygulamada bu sayı $\sim$600 civarına düşer.)

(b) Verim formülü:
\[
V = \left(1 + \frac{d \times A_{\text{kırmık}}}{\alpha}\right)^{-\alpha} = \left(1 + \frac{0{,}04 \times 1}{4}\right)^{-4} = (1{,}01)^{-4} \approx 0{,}961 = \%96{,}1
\]
(Burada $A_{\text{kırmık}} = 100\text{\,mm}^2 = 1\text{\,cm}^2$.)

(c) $A_{\text{kırmık}} = 50\text{\,mm}^2 = 0{,}5\text{\,cm}^2$ için:
\[
V = \left(1 + \frac{0{,}04 \times 0{,}5}{4}\right)^{-4} = (1{,}005)^{-4} \approx 0{,}980 = \%98{,}0
\]

Kırmık boyutu yarıya düşünce verim \%96,1'den \%98,0'a yükselir ve plaka başına çalışan kırmık sayısı iki kattan fazla artar. Bu, Moore yasasını destekleyen ekonomik motivasyonun önemli bir parçasıdır: daha küçük transistör $\rightarrow$ daha küçük kırmık $\rightarrow$ daha yüksek verim $\rightarrow$ daha düşük birim maliyet.
\end{ornek}


\section{Programların Çalıştırılması}
\label{sec:program-calistirma}

Şimdiye kadar bilgisayarın fiziksel bileşenlerini tanıdık: işlemci, bellek, giriş/çıkış aygıtları ve bunların bir yonga üzerinde nasıl üretildiği. Ancak donanım tek başına anlamsızdır; onu yararlı kılan şey üzerinde çalışan programlardır. Bu bölümde donanımla yazılım arasındaki köprüyü inceleyeceğiz: programlar hangi dilde yazılır, belleğe nasıl yerleştirilir, işlemci tarafından nasıl yorumlanır ve bir bütün olarak nasıl yürütülür?

\subsection{Makine Dili ve Çevirici Dil}
\label{subsec:makine-dili-cevirici}

Bir programı makine üzerinde çalıştırmak için öncelikle makineyle aynı dilin konuşulması gerekir. Bilgisayarların alfabesi 1'ler ve 0'lardan oluşur; bu alfabedeki harflerden oluşan sözcükler ise ikilik tabanda yazılmış sayılardır. Sözcüklerdeki 1 ya da 0 rakamlarından oluşan her bir basamağa ``Binary Digit'' veya kısaca ``Bit''\index{bit} denir. Sekiz bitin bir araya gelmesiyle bir \emph{bayt}\index{bayt} (\emph{byte}) oluşur; bayt, bilgisayarların en temel adreslenebilir veri birimidir. Bilgisayarların anladığı kelimeler 1 ve 0'lardan oluştuğu için üst düzey dilde yazılanların bilgisayar tarafından anlaşılması bunların ikilik düzene çevrilmesiyle mümkün olur.

\emph{Makine dili}\index{makine dili}, işlemcinin doğrudan anlayıp yürütebildiği tek dildir. Her makine dili buyruğu üç temel bilgi taşır: (1) yapılacak işlemi belirleyen \emph{işlem kodu}\index{işlem kodu} (\emph{opcode}), (2) işlemin üzerinde gerçekleştirileceği verilerin konumunu gösteren \emph{işlenen}\index{işlenen} (\emph{operand}) alanları ve (3) buyruğun biçimini tanımlayan ek denetim bitleri. Örneğin RISC-V\index{RISC-V} mimarisinde bir toplama buyruğu 32~bit genişliğindedir ve bu bitler arasında işlem kodu, kaynak yazmaçlarını belirten alanlar ve hedef yazmacı belirten alan yer alır. İlk bilgisayarlarda programcılar bu ikilik kodları elle yazıyordu; yüzlerce satırlık bir programda tek bir bitin yanlış girilmesi saatlerce hata aramayı gerektiriyordu.

Bu zorlukları aşmak için 1950'lerde \emph{çevirici dil}\index{çevirici dil} (\emph{Assembly}\index{Assembly}) geliştirildi. Çevirici dil, ikilik kodlar yerine insanların okuyabileceği kısa anımsatıcılar (\emph{mnemonic}) kullanır: \texttt{add} toplama, \texttt{sub} çıkarma, \texttt{lw} bellekten yükleme anlamına gelir. Her anımsatıcı, tam olarak bir makine buyruğuna karşılık düşer. Çevirici dilde yazılmış kaynak kodu, \emph{çevirici}\index{çevirici} (\emph{Assembler}) adı verilen bir program tarafından makine diline dönüştürülür. Bu dönüşüm bire bir olduğu için çevirici dil, makine diline en yakın programlama düzeyidir. Aşağıdaki örnek bu ilişkiyi somutlaştırır:

\begin{center}
\footnotesize
\begin{tabular}{@{}lll@{}}
\toprule
\textbf{Düzey} & \textbf{Gösterim} & \textbf{Açıklama} \\
\midrule
Çevirici dil & \texttt{add x5, x6, x7} & x6 ile x7'yi topla, sonucu x5'e yaz \\
Makine dili (ikilik) & \texttt{0000000\,00111\,00110\,000\,00101\,0110011} & 32 bitlik R-tipi RISC-V buyruğu \\
Makine dili (onaltılık) & \texttt{0x007302B3} & Aynı buyruğun onaltılık gösterimi \\
\bottomrule
\end{tabular}
\end{center}

Çevirici dilin en önemli avantajı, programcının kaynak ve hedef yazmaçlarını (\texttt{x5}, \texttt{x6}, \texttt{x7}) isimle görmesini sağlaması ve işlem kodlarını anlamlı kısaltmalarla temsil etmesidir. Bununla birlikte çevirici dil, mimariye özgüdür: bir RISC-V işlemcisi için yazılmış çevirici dil kodu, x86 ya da ARM işlemcisinde çalışmaz. Bu sınırlama, ilerleyen alt bölümlerde ele alınacak üst düzey dillerin ve derleyicilerin geliştirilmesinin temel motivasyonlarından biri olmuştur.

\subsection{Von Neumann ve Harvard Mimarileri}
\label{subsec:vonneumann-harvard}

Buyruklar ve veriler bellekte nasıl saklanmalı? Bu soru, bilgisayar mimarisinin en temel tasarım kararlarından birini oluşturur ve tarihsel olarak iki farklı yaklaşım geliştirilmiştir.

\emph{Von Neumann mimarisi}\index{Von Neumann mimarisi}, John Von Neumann'ın 1945'te EDVAC tasarımı için kaleme aldığı raporda~\cite{vonneumann1945} ortaya koyduğu ilkelere dayanır. Bu mimaride buyruklar ve veriler tek bir ortak bellekte saklanır; işlemci aynı yol üzerinden hem buyruğu çeker hem de veriyi okur/yazar. Ortak bellek kullanımının en büyük avantajı esnekliktir: programlar veriyle aynı biçimde depolanabilir, kopyalanabilir ve değiştirilebilir. Bu sayede bir program başka bir programı belleğe yükleyebilir; işletim sistemlerinin ve derleyicilerin çalışma ilkesi tam da budur. Ancak tek yoldan hem buyruk hem veri geçmesi kaçınılmaz bir yarışma yaratır: işlemci aynı anda buyruğu çekip veriyi okuyamaz. Bu sınırlama \emph{Von Neumann darboğazı}\index{Von Neumann darboğazı} olarak bilinir ve günümüzde bile başarımı etkileyen temel sorunlardan biridir.

\emph{Harvard mimarisi}\index{Harvard mimarisi} bu darboğazı ortadan kaldırmayı hedefler. İsmini 1944'te Howard Aiken'in Harvard Üniversitesi'nde tasarladığı Mark~I bilgisayarından alan bu yaklaşımda buyruklar ve veriler fiziksel olarak ayrı belleklerde tutulur ve her birinin kendi yolu vardır. İşlemci, aynı anda buyruk belleğinden bir sonraki buyruğu çekerken veri belleğinden de ihtiyaç duyduğu veriyi okuyabilir. Bu koşut erişim, özellikle gerçek zamanlı sinyal işleme gibi bant genişliğine duyarlı uygulamalarda belirgin bir başarım artışı sağlar. Dezavantajı ise esneklik kaybıdır: buyruk ve veri belleklerinin boyutları tasarım zamanında belirlenmek zorundadır ve biri boşken diğeri doluysa kapasite israf edilir.

\begin{figure}[H]
\centering
\begin{tikzpicture}[>=Stealth, scale=0.85, every node/.style={transform shape}]
    % --- Von Neumann (sol) ---
    \node[font=\small\bfseries, text=sinyalmavi] at (-4.5, 3.5) {Von Neumann Mimarisi};

    \node[blok, fill=blokyesil, minimum width=2.8cm, minimum height=1.2cm] (vn-cpu) at (-4.5, 2)
        {\footnotesize\textbf{İşlemci}\\[-1pt]{\scriptsize Denetim + Veri Yolu}};
    \node[blok, fill=bloksari, minimum width=2.8cm, minimum height=1.2cm] (vn-mem) at (-4.5, -0.5)
        {\footnotesize\textbf{Bellek}\\[-1pt]{\scriptsize Buyruk + Veri}};
    \node[blok, fill=blokmavi!40, minimum width=1.2cm, minimum height=0.8cm] (vn-in) at (-7, 2)
        {\scriptsize\textbf{Giriş}};
    \node[blok, fill=blokmavi!40, minimum width=1.2cm, minimum height=0.8cm] (vn-out) at (-2, 2)
        {\scriptsize\textbf{Çıkış}};

    \draw[okstil, sinyalmavi, <->] (vn-cpu) -- (vn-mem)
        node[midway, right=3pt, font=\footnotesize] {buyruk/veri};
    \draw[okstil, gray, ->] (vn-in) -- (vn-cpu);
    \draw[okstil, gray, ->] (vn-cpu) -- (vn-out);

    % --- Harvard (sağ) ---
    \node[font=\small\bfseries, text=sinyalkirmizi!80!black] at (4.5, 3.5) {Harvard Mimarisi};

    \node[blok, fill=blokyesil, minimum width=2.8cm, minimum height=1.2cm] (h-cpu) at (4.5, 2)
        {\footnotesize\textbf{İşlemci}\\[-1pt]{\scriptsize Denetim + Veri Yolu}};
    \node[blok, fill=bloksari!70, minimum width=2.2cm, minimum height=1cm] (h-imem) at (2.8, -0.5)
        {\scriptsize\textbf{Buyruk}\\[-1pt]\scriptsize\textbf{Belleği}};
    \node[blok, fill=bloksari, minimum width=2.2cm, minimum height=1cm] (h-dmem) at (6.2, -0.5)
        {\scriptsize\textbf{Veri}\\[-1pt]\scriptsize\textbf{Belleği}};
    \node[blok, fill=blokmavi!40, minimum width=1.2cm, minimum height=0.8cm] (h-in) at (2, 2)
        {\scriptsize\textbf{Giriş}};
    \node[blok, fill=blokmavi!40, minimum width=1.2cm, minimum height=0.8cm] (h-out) at (7, 2)
        {\scriptsize\textbf{Çıkış}};

    \draw[okstil, sinyalkirmizi!80!black, <->] (h-cpu.south west) ++(0.3,0) -- (h-imem)
        node[midway, left=3pt, font=\footnotesize] {buyruk};
    \draw[okstil, sinyalkirmizi!80!black, <->] (h-cpu.south east) ++(-0.3,0) -- (h-dmem)
        node[midway, right=3pt, font=\footnotesize] {veri};
    \draw[okstil, gray, ->] (h-in) -- (h-cpu);
    \draw[okstil, gray, ->] (h-cpu) -- (h-out);
\end{tikzpicture}
\caption[Von Neumann ve Harvard mimarileri]{Von Neumann mimarisinde buyruk ve veri aynı belleği paylaşır (tek yol); Harvard mimarisinde ise ayrı belleklerde tutulur (iki bağımsız yol).}
\label{fig:vonneumann-harvard}
\end{figure}

Günümüzde çoğu genel amaçlı işlemci, her iki yaklaşımın avantajlarını birleştiren \emph{değiştirilmiş Harvard mimarisi}\index{değiştirilmiş Harvard mimarisi} kullanır: ana bellekte buyruklar ve veriler birlikte saklanır (Von Neumann esnekliği), ancak işlemcinin içindeki önbelleklerde buyruk önbelleği ile veri önbelleği ayrı tutulur (Harvard bant genişliği). Bu tasarım, Bölüm~\ref{chap:bellek}'de ele alınacak bellek hiyerarşisinin temel yapı taşlarından biridir.

\subsection{RISC ve CISC Tasarım Felsefeleri}
\label{subsec:risc-cisc}

Bir buyruk kümesi mimarisi tasarlanırken karşılaşılan temel soru şudur: buyruklar ne kadar karmaşık olmalı? Tarihsel süreçte bu soruya iki farklı felsefe yanıt vermiştir.

\emph{CISC}\index{CISC} (\emph{Complex Instruction Set Computer}: Karmaşık Buyruk Kümeli Bilgisayar) yaklaşımı, az sayıda buyrukla çok iş yapmayı hedefler. CISC mimarisinde tek bir buyruk, bellekten veri okuma, üzerinde işlem yapma ve sonucu belleğe yazma gibi birden fazla adımı bir arada gerçekleştirebilir. Bu yaklaşım, belleğin pahalı ve yavaş olduğu dönemlerde programları küçük tutmak amacıyla geliştirildi: daha az buyruk demek daha az bellek kullanımı demekti. Intel'in x86\index{x86} ailesi, CISC felsefesinin en yaygın temsilcisidir.

\emph{RISC}\index{RISC} (\emph{Reduced Instruction Set Computer}: İndirgenmiş Buyruk Kümeli Bilgisayar) yaklaşımı ise 1980'lerde Berkeley ve Stanford üniversitelerindeki araştırmalardan doğdu. RISC felsefesinde her buyruk yalnızca bir temel işlem yapar, tüm buyruklar sabit uzunluktadır ve bellek erişimi yalnızca özel yükleme (\texttt{lw}) ve saklama (\texttt{sw}) buyruklarıyla gerçekleştirilir; aritmetik işlemler doğrudan bellekteki veriler üzerinde değil, yazmaçlardaki\index{yazmaç} değerler üzerinde yapılır. Bu ayrım \emph{yükle-sakla mimarisi}\index{yükle-sakla mimarisi} (\emph{load-store architecture}) olarak bilinir. ARM\index{ARM} ve RISC-V\index{RISC-V}, RISC felsefesinin günümüzdeki önemli temsilcileridir.

Farkı somut bir örnekle görelim. ``Bellekteki bir değeri bir yazmacın üstüne ekle'' işlemini düşünelim (bu işlem C dilinde \texttt{A = A + B[i]} gibi bir ifadeye karşılık gelir):

\begin{center}
\footnotesize
\begin{tabular}{@{}p{3.2cm}p{5.5cm}p{5.5cm}@{}}
\toprule
 & \textbf{CISC (x86 benzeri)} & \textbf{RISC (RISC-V benzeri)} \\
\midrule
Buyruk dizisi &
\texttt{add eax, [ebx+ecx*4]} \newline {\scriptsize (Tek buyruk: bellek oku + topla)} &
\texttt{slli t1, t0, 2} \newline
\texttt{add t1, s1, t1} \newline
\texttt{lw t2, 0(t1)} \newline
\texttt{add s0, s0, t2} \newline
{\scriptsize (Dört buyruk: kaydır + adres hesapla + yükle + topla)} \\
\addlinespace
Buyruk sayısı & 1 & 4 \\
Buyruk uzunluğu & Değişken (1--15 bayt) & Sabit (4 bayt) \\
Donanım karmaşıklığı & Yüksek (karmaşık çözücü) & Düşük (yalın boru hattı) \\
\bottomrule
\end{tabular}
\end{center}

Bu örnekte görüldüğü gibi CISC, aynı işi daha az buyrukla yapabilir; ancak her buyruğun uzunluğu ve yürütüm süresi farklı olduğundan donanım tasarımı karmaşıklaşır. RISC ise daha fazla buyruk kullanır; buna karşılık her buyruğun sabit uzunlukta ve öngörülebilir sürede tamamlanması, boru hattı\index{boru hattı} (\emph{pipeline}) tasarımını kolaylaştırır ve saat sıklığının artırılmasına olanak tanır.

\begin{table}[H]
\centering
\caption{RISC ve CISC tasarım felsefelerinin karşılaştırması}
\label{tab:risc-cisc}
\footnotesize
\begin{tabular}{@{}lll@{}}
\toprule
\textbf{Özellik} & \textbf{RISC} & \textbf{CISC} \\
\midrule
Buyruk sayısı         & Az (tipik 50--150)     & Çok (tipik 500--1500+) \\
Buyruk uzunluğu       & Sabit (32 bit)         & Değişken (8--120+ bit) \\
Bellek erişim modeli   & Yalnızca yükle/sakla   & Doğrudan bellekle işlem \\
Yazmaç sayısı         & Çok (32+)              & Az (8--16) \\
Adresleme kipleri     & Az ve yalın            & Çok ve çeşitli \\
Donanım karmaşıklığı  & Düşük                  & Yüksek \\
Derleyici karmaşıklığı& Yüksek                 & Düşük \\
Program boyutu        & Büyük                  & Küçük \\
Örnekler              & RISC-V, ARM, MIPS      & x86, VAX, IBM z/Architecture \\
\bottomrule
\end{tabular}
\end{table}

Günümüzde bu iki felsefe arasındaki sınır eskisi kadar keskin değildir. Çağdaş x86 işlemcileri, karmaşık CISC buyruklarını donanım içinde basit \emph{mikro-işlemlere}\index{mikro-işlem} (\emph{micro-op}) böler ve bu mikro-işlemleri RISC benzeri bir boru hattında yürütür. Öte yandan RISC mimarileri de özelleşmiş buyruklar (vektör uzantıları, atomik işlemler gibi) ekleyerek buyruk kümelerini genişletmiştir. Sonuç olarak bugünün başarım farkları mimari felsefeden çok mikromimari tasarım kararlarına (önbellek boyutuna, boru hattı derinliğine, dal öngörücüsü kalitesine) bağlıdır.

\begin{bilginot}[RISC-V: Bu Kitabın Buyruk Kümesi]
Bu kitapta buyruk kümesi mimarisi olarak açık kaynaklı \textbf{RISC-V}\index{RISC-V} kullanılıyor. RISC\index{RISC}, \emph{Reduced Instruction Set Computer}\index{RISC!açılım} (İndirgenmiş Buyruk Kümeli Bilgisayar) kısaltmasıdır ve 1980'lerde Berkeley ile Stanford'da geliştirilen tasarım felsefesini ifade eder. Berkeley'de bu felsefe çerçevesinde art arda beş nesil işlemci tasarlandı: RISC-I\index{RISC-I} (1981), RISC-II\index{RISC-II} (1983), SOAR (1984), SPUR (1988) ve son olarak 2010'da başlayan proje. Beşinci nesil olduğu için adındaki ``V'' Romen rakamı 5'tir.

RISC-V ``\textbf{risk beş}'' diye okunur (\emph{risk five}); ``risk-vee'' okuyuşu yanlıştır. Türkçe ek alırken son ses /ş/ sert ünsüz olduğundan sert ünsüz uyumu uygulanır: RISC-V'\emph{te}, RISC-V'\emph{ten}, RISC-V'\emph{in}, RISC-V'\emph{e}. RISC-V'in neden tercih edildiği, tescilli mimarilerle karşılaştırması ve endüstrideki yeri Bölüm~\ref{sec:neden-riscv}'da ayrıntılı olarak ele alınıyor.
\end{bilginot}

\subsection{Üst Düzey Diller ve Derleme Zinciri}
\label{subsec:ust-duzey-diller}

Mimari tasarımındaki gelişmeler ve buyrukların çeviricilerle makine koduna çevrilmesi kullanıcılara büyük kolaylık sağladı. Ancak bilimsel çalışmalar, modellemeler ve karmaşık problemler için daha karmaşık programlama yöntemlerine gereksinim duyuldu ve bu gereksinimi karşılamak için üst düzey diller oluşturuldu. Üst düzey dillerde neredeyse metin gibi okunup anlaşılabilecek kadar anlamlı yapılar kullanılır. Üst düzey bir dilde yazılmış kodları çevirici dile dönüştüren \emph{derleyici}\index{derleyici} (\emph{Compiler}) programlar yazıldı. Programcı \texttt{A = B + C;} yazdığında, derleyici bu kodu \texttt{add A, B, C} haline getirir ve çevirici de bu kodu \texttt{1000\allowbreak 1010\allowbreak 1010\allowbreak 1010\allowbreak 1100\allowbreak 0100\allowbreak 0010\allowbreak 0010} olarak makine diline çevirir. Şekil~\ref{fig:ust-duzey-cevrim}'da üst düzey bir dilde yazılmış bir programın nasıl en alt düzeye çevrildiği bir RISC-V\index{RISC-V} işlemcisi örneğiyle gösteriliyor.

\begin{figure}[H]
\centering
\begin{tikzpicture}[>=Stealth, node distance=1.8cm]
    % Üst düzey program kutusu
    \node[blok, fill=blokyesil, minimum width=7cm, minimum height=1.4cm] (ust)
        {\textbf{Üst Düzey Program}\\[2pt]\texttt{A = B + C;}};

    % Derleyici oku
    \node[below=0.6cm of ust, font=\small\itshape, text=sinyalmavi] (derleyici) {Derleyici};

    % Çevirici dil kutusu
    \node[blok, fill=blokmavi, minimum width=7cm, minimum height=1.4cm, below=0.6cm of derleyici] (cevirici)
        {\textbf{Çevirici Dil (Assembly)}\\[2pt]\texttt{add A, B, C}};

    % Çevirici oku
    \node[below=0.6cm of cevirici, font=\small\itshape, text=sinyalmavi] (ceviriciprg) {Çevirici (Assembler)};

    % Makine dili kutusu
    \node[blok, fill=bloksari, minimum width=7cm, minimum height=1.4cm, below=0.6cm of ceviriciprg] (makine)
        {\textbf{Makine Dili}\\[2pt]\texttt{1000\,1010\,1010\,1010\,1100\,0100\,0010\,0010}};

    % Oklar
    \draw[okstil, sinyalmavi] (ust) -- (derleyici);
    \draw[okstil, sinyalmavi] (derleyici) -- (cevirici);
    \draw[okstil, sinyalmavi] (cevirici) -- (ceviriciprg);
    \draw[okstil, sinyalmavi] (ceviriciprg) -- (makine);
\end{tikzpicture}
\caption{Bir üst düzey programın alt düzey denetim işaretlerine çevrim aşamaları}
\label{fig:ust-duzey-cevrim}
\end{figure}

\subsection{Yazılım Katmanları}
\label{subsec:yazilim-katmanlari}

Üst düzey dillerin ortaya çıkması, programların yazılmasını ve okunmasını büyük ölçüde hızlandırdı. Farklı ihtiyaçlara uygun diller geliştirildi: bilimsel hesaplamalar için Fortran (1957), ticari uygulamalar için COBOL (1959), sistem programlama için C (1972), nesne yönelimli geliştirme için Java (1995) ve C++ (1985). Ancak tek başına bir programlama dili yeterli değildir; dil yalnızca programları ifade etmenin bir yoludur. Programların çalışması için altlarında katmanlı bir yazılım altyapısı gerekir.

Yazılımlar kullanılma amaçlarına göre \emph{sistem yazılımları}\index{sistem yazılımı} ve \emph{uygulama yazılımları}\index{uygulama yazılımı} olmak üzere iki gruba ayrılır:

\begin{itemize}[nosep, leftmargin=*]
    \item \emph{Uygulama yazılımları}, kullanıcının doğrudan etkileştiği programlardır: tarayıcılar, metin düzenleyiciler, oyunlar, hesap çizelgeleri, yapay zekâ uygulamaları gibi. Kullanıcı bir uygulamayı çalıştırdığında, aslında alttaki tüm katmanları dolaylı olarak kullanır.
    \item \emph{Sistem yazılımları}, donanım ile uygulama yazılımları arasında köprü kurar ve donanım kaynaklarını yönetir. Bu gruptaki en önemli iki bileşen \emph{işletim sistemi}\index{işletim sistemi} ve \emph{derleyici}\index{derleyici}dir.
\end{itemize}

\emph{İşletim sistemi}\index{işletim sistemi} (\emph{Operating System}), bilgisayarın kaynaklarını (işlemci zamanı, bellek alanı, giriş/çıkış aygıtları) uygulamalar arasında paylaştırır ve düzenler. Birden fazla programın aynı anda çalışabilmesi (çoklu görev), dosya sisteminin yönetilmesi ve kullanıcıyla etkileşim sağlanması işletim sisteminin sorumluluğundadır. Linux, Windows, macOS, Android ve iOS günümüzün en yaygın işletim sistemleridir.

\emph{Derleyici}\index{derleyici} (\emph{Compiler}), üst düzey dilde yazılmış kaynak kodu önce çevirici dile, ardından çevirici aracılığıyla makine diline dönüştüren programdır. Derleyicinin görevi yalnızca çeviri yapmak değildir; aynı zamanda programı eniyileyerek (döngü açma, ölü kod eleme, yazmaç ataması gibi tekniklerle) donanımdan en yüksek başarımı elde etmeye çalışır. Bu nedenle derleyici kalitesi, aynı kaynak kodunun farklı donanımlarda ne kadar hızlı çalışacağını doğrudan etkiler.

Bu katmanlar arasındaki ilişki Şekil~\ref{fig:yazilim-katmanlari}'da gösterilmektedir. Her katman yalnızca bir altındaki katmanın sunduğu \emph{arayüzü} kullanır; altındaki katmanın iç ayrıntılarını bilmek zorunda değildir. Bu yapı, Şekil~\ref{fig:soyutlama-katmanlari}'da ele alınan soyutlama ilkesinin yazılım tarafındaki yansımasıdır.

\begin{figure}[H]
\centering
\begin{tikzpicture}[>=Stealth, scale=0.88, every node/.style={transform shape}]
    \def\katmanaraligi{1.5}  % kutular arası dikey mesafe
    \tikzstyle{ykatman}=[draw, thick, rounded corners=4pt,
        text width=10cm, minimum height=1.2cm, font=\small, align=center,
        drop shadow={shadow xshift=0.3mm, shadow yshift=-0.3mm, opacity=0.12}]

    \node[ykatman, fill=bloksari!70] (app) at (0, 4*\katmanaraligi)
        {Uygulama Yazılımları\\{\footnotesize(Tarayıcı, Oyun, Yapay Zekâ Uygulaması, \ldots)}};
    \node[ykatman, fill=bloksari!45] (lib) at (0, 3*\katmanaraligi)
        {Kütüphaneler\\{\footnotesize(Standart kütüphane, grafik, matematik, \ldots)}};
    \node[ykatman, fill=blokyesil!50] (os) at (0, 2*\katmanaraligi)
        {İşletim Sistemi\\{\footnotesize(Süreç yönetimi, bellek yönetimi, dosya sistemi, sürücüler)}};
    \node[ykatman, fill=blokmavi!40] (fw) at (0, 1*\katmanaraligi)
        {Aygıt Yazılımı / BIOS\\{\footnotesize(Donanım başlatma, düşük düzey denetim)}};
    \node[ykatman, fill=gray!25] (hw) at (0, 0)
        {Donanım\\{\footnotesize(İşlemci, Bellek, Giriş/Çıkış Aygıtları)}};

    % Sağ tarafta parantez ve etiketler (brace sağa açılacak şekilde: üstten alta)
    \draw[decorate, decoration={brace, amplitude=8pt}, thick, gray!70]
        (5.8, 4*\katmanaraligi+0.6) -- (5.8, 3*\katmanaraligi-0.6)
        node[midway, right=10pt, font=\footnotesize, align=left, text=gray!70] {Uygulama\\yazılımları};
    \draw[decorate, decoration={brace, amplitude=8pt}, thick, gray!70]
        (5.8, 2*\katmanaraligi+0.6) -- (5.8, 1*\katmanaraligi-0.6)
        node[midway, right=10pt, font=\footnotesize, align=left, text=gray!70] {Sistem\\yazılımları};
\end{tikzpicture}
\caption{Yazılım katmanları ve donanım arasındaki ilişki}
\label{fig:yazilim-katmanlari}
\end{figure}

\emph{Kütüphaneler}\index{kütüphane} (\emph{library}), sık kullanılan işlevlerin bir araya getirildiği hazır kod koleksiyonlarıdır. Bir programcı matematiksel bir fonksiyon hesaplamak ya da ekrana görüntü çizmek istediğinde ilgili kütüphaneyi çağırır; aynı kodu sıfırdan yazmak zorunda kalmaz. Kütüphaneler, yazılım geliştirme hızını artırmanın yanı sıra test edilmiş ve eniyilenmiş kodun tekrar tekrar kullanılmasını sağlar.

Bu katmanlı yapı, bilgisayar mimarisi açısından can alıcı bir sonuç doğurur: bir programın başarımı yalnızca donanıma değil, derleyicinin eniyileme kalitesine, işletim sisteminin kaynak yönetim politikalarına ve kütüphanelerin verimliliğine de bağlıdır. Mimari tasarımcılar, buyruk kümesini ve donanımı bu yazılım katmanlarıyla uyumlu biçimde geliştirmek zorundadır. Bu bölümde tanıttığımız komut-derleme-çalıştırma zincirini RISC-V buyruklarıyla ayrıntılı biçimde Bölüm~\ref{chap:buyruk-kumesi}'te, çevirici dil programlamayı ise Bölüm~\ref{chap:programlama}'te ele alacağız.

\subsection{Buyruk Yürütüm Döngüsü}
\label{subsec:buyruk-yurutum-dongusu}

Derleyici ve çevirici yoluyla makine diline çevrilen buyrukların yürütüm döngüsü Şekil~\ref{fig:buyruk-yurutum}'da veriliyor. İlk olarak programın saklandığı yerden sıradaki buyruk çekilir ve kullanılan buyruk kümesinin tanımlamalarına göre yapılacak işlem ve buyruğun büyüklüğü belirlenir. İşlenecek veriler bulundukları konumlardan okunur, verilerin üzerinde işlem yapılır, sonuçlar ya da durumlar daha sonra kullanılmak üzere saklanır ve bir sonraki buyruk bellekten çekilmek üzere, işlenen buyruğun bellekteki konumunun adresini tutan Program Sayacı\index{program sayacı} yazmacının değeri güncellenir. Programın son buyruğuna kadar bu döngü sürer.

\begin{figure}[H]
\centering\vspace{-4pt}
\begin{tikzpicture}[>=Stealth, node distance=1.8cm]
    % Altıgen/dairesel yerleşim açıları: 90, 30, -30, -90, -150, 150
    \def\radius{2.6cm}
    % Düğümler
    \node[akisislem, minimum width=2.5cm, inner sep=3pt] (cek) at (90:\radius)
        {Buyruğu Çek\\[-1pt]{\footnotesize\itshape (Fetch)}};
    \node[akisislem, minimum width=2.5cm, inner sep=3pt] (coz) at (30:\radius)
        {Buyruğu Çöz\\[-1pt]{\footnotesize\itshape (Decode)}};
    \node[akisislem, minimum width=2.5cm, inner sep=3pt] (oku) at (-30:\radius)
        {Verileri Oku\\[-1pt]{\footnotesize\itshape (Read)}};
    \node[akisislem, minimum width=2.5cm, inner sep=3pt] (yurut) at (-90:\radius)
        {İşlemi Yürüt\\[-1pt]{\footnotesize\itshape (Execute)}};
    \node[akisislem, minimum width=2.5cm, inner sep=3pt] (yaz) at (-150:\radius)
        {Sonucu Yaz\\[-1pt]{\footnotesize\itshape (Write Back)}};
    \node[akisislem, minimum width=2.5cm, inner sep=3pt] (ps) at (150:\radius)
        {PS Güncelle\\[-1pt]{\footnotesize\itshape (Update PC)}};

    % Oklar (saat yönünde)
    \draw[akisok, sinyalmavi] (cek) -- (coz);
    \draw[akisok, sinyalmavi] (coz) -- (oku);
    \draw[akisok, sinyalmavi] (oku) -- (yurut);
    \draw[akisok, sinyalmavi] (yurut) -- (yaz);
    \draw[akisok, sinyalmavi] (yaz) -- (ps);
    \draw[akisok, sinyalmavi] (ps) -- (cek);

    % Ortadaki etiket
    \node[font=\small\bfseries, text=sinyalmavi, align=center] at (0,0) {Buyruk\\Döngüsü};
\end{tikzpicture}\vspace{-4pt}
\caption{Buyruk yürütüm döngüsü}
\label{fig:buyruk-yurutum}
\end{figure}

\begin{bilginot}[Getir--Çöz--Yürüt ve Bulaşık Makinesi Benzetmesi]
Bir bulaşık makinesini düşünün: kirli tabaklar konulur (getir), makine programını seçer (çöz) ve yıkama başlar (yürüt). Her seferinde tek bir tabak yıkamak yerine, birden çok tabağı sırayla makineye vermek verimliliği artırır; tıpkı boru hattının (\emph{pipeline}) birden çok buyruğu çakıştırarak (\emph{overlap}) işlemesi gibi. Boru hattı kavramı Bölüm~7'de ayrıntılı olarak ele alınacaktır.
\end{bilginot}

\section{Bilgisayar Mimarisini Şekillendiren Etkenler}
\label{sec:sekillendiren-etkenler}

Bilgisayar, birbirini tamamlayan pek çok ayrı parçadan oluşur. Bu parçaların birlikte verimli ve hatasız çalışması için birimlerin organizasyonu oldukça önemli. Bir bilgisayarın mimarisini şekillendiren beş temel etken Şekil~\ref{fig:mimari-etkenler}'de gösteriliyor.

\begin{figure}[H]
\centering
\begin{tikzpicture}[>=Stealth, scale=0.88, every node/.style={transform shape},
    % Zihin haritası stilleri
    dalstil/.style={line width=1.6pt, rounded corners=2pt},
    etken/.style={draw, thick, rounded corners=5pt, minimum width=2.4cm,
        minimum height=0.9cm, font=\small\bfseries, align=center, drop shadow={
        shadow xshift=0.4mm, shadow yshift=-0.4mm, opacity=0.15}},
    yaprakkutu/.style={draw=none, rounded corners=3pt, inner sep=2.5pt,
        font=\scriptsize, align=left},
    dalcik/.style={semithick, rounded corners=3pt},
]
    % Merkez düğüm
    \node[draw=sinyalmavi!80!black, line width=1.8pt, rounded corners=7pt,
          minimum width=3.2cm, minimum height=1.5cm,
          fill=blokmavi, font=\normalsize\bfseries, align=center,
          drop shadow={shadow xshift=0.6mm, shadow yshift=-0.6mm, opacity=0.2}]
          (merkez) {Bilgisayar\\Mimarisi};

    % ── 1. TEKNOLOJİ (üst) ──
    \node[etken, fill=blokyesil, draw=sinyalyesil!70!black] (tek) at (0,3.5)
        {Teknoloji};
    \draw[dalstil, sinyalyesil!70!black] (merkez.north) -- (tek.south);
    % Yapraklar
    \node[yaprakkutu, fill=blokyesil-acik, anchor=east] (tek1) at (-2.2,4.6)
        {Transistör boyutu\\küçülmesi};
    \node[yaprakkutu, fill=blokyesil-acik, anchor=west] (tek2) at (2.2,4.6)
        {Güç yoğunluğu\\ve ısınma};
    \draw[dalcik, sinyalyesil!60] (tek.west) -| (tek1.south);
    \draw[dalcik, sinyalyesil!60] (tek.east) -| (tek2.south);

    % ── 2. PROGRAMLAMA DİLLERİ (sol üst) ──
    \node[etken, fill=bloksari, draw=sinyalkirmizi!60!black] (prd) at (-3.8,1.6)
        {Programlama\\[-1pt]Dilleri};
    \draw[dalstil, sinyalkirmizi!60!black] (merkez) -- (prd);
    % Yapraklar
    \node[yaprakkutu, fill=bloksari-acik, anchor=east] (prd1) at (-6.5,2.1)
        {Dil yapılarının\\donanım desteği};
    \node[yaprakkutu, fill=bloksari-acik, anchor=east] (prd2) at (-6.5,1.1)
        {Derleme\\eniyilemesi};
    \draw[dalcik, sinyalkirmizi!40!bloksari] (prd.west) -- ++(-0.25,0) |- (prd1.east);
    \draw[dalcik, sinyalkirmizi!40!bloksari] (prd.west) -- ++(-0.25,0) |- (prd2.east);

    % ── 3. UYGULAMA (sol alt) ──
    \node[etken, fill=blokkirmizi, draw=sinyalkirmizi!70!black] (uyg) at (-3.8,-1.6)
        {Uygulama};
    \draw[dalstil, sinyalkirmizi!70!black] (merkez) -- (uyg);
    % Yapraklar
    \node[yaprakkutu, fill=blokkirmizi-acik, anchor=east] (uyg1) at (-6.5,-1.1)
        {Bilimsel hesaplama\\ve yapay zekâ};
    \node[yaprakkutu, fill=blokkirmizi-acik, anchor=east] (uyg2) at (-6.5,-2.1)
        {Gömülü sistemler\\ve mobil};
    \draw[dalcik, sinyalkirmizi!50] (uyg.west) -- ++(-0.25,0) |- (uyg1.east);
    \draw[dalcik, sinyalkirmizi!50] (uyg.west) -- ++(-0.25,0) |- (uyg2.east);

    % ── 4. İŞLETİM SİSTEMLERİ (sağ alt) ──
    \node[etken, fill=blokmavi, draw=sinyalmavi!70!black] (isl) at (3.8,-1.6)
        {İşletim\\[-1pt]Sistemleri};
    \draw[dalstil, sinyalmavi!70!black] (merkez) -- (isl);
    % Yapraklar
    \node[yaprakkutu, fill=blokmavi-acik, anchor=west] (isl1) at (6.5,-1.1)
        {Bellek koruması\\ve sanal bellek};
    \node[yaprakkutu, fill=blokmavi-acik, anchor=west] (isl2) at (6.5,-2.1)
        {Bağlam değiştirme\\ve kesmeler};
    \draw[dalcik, sinyalmavi!50] (isl.east) -- ++(0.25,0) |- (isl1.west);
    \draw[dalcik, sinyalmavi!50] (isl.east) -- ++(0.25,0) |- (isl2.west);

    % ── 5. GEÇMİŞ (sağ üst) ──
    \node[etken, fill=blokgri, draw=gray!70!black] (tar) at (3.8,1.6)
        {Geçmiş\\[-1pt](Tarih)};
    \draw[dalstil, gray!60!black] (merkez) -- (tar);
    % Yapraklar
    \node[yaprakkutu, fill=blokgri-acik, anchor=west] (tar1) at (6.5,2.1)
        {Geçmişe\\uyumluluk};
    \node[yaprakkutu, fill=blokgri-acik, anchor=west] (tar2) at (6.5,1.1)
        {x86'nın 1978'den\\beri korunması};
    \draw[dalcik, gray!50] (tar.east) -- ++(0.25,0) |- (tar1.west);
    \draw[dalcik, gray!50] (tar.east) -- ++(0.25,0) |- (tar2.west);

    % ── Etkileşim yayları (etkenler arası dolaylı ilişki) ──
    \draw[dashed, gray!40, thin]
        (tek.west) to[bend right=20] (prd.north);
    \draw[dashed, gray!40, thin]
        (isl.south) to[bend left=20] (uyg.south);
    \draw[dashed, gray!40, thin]
        (tar.north) to[bend right=20] (tek.east);
\end{tikzpicture}
\caption{Bilgisayar Mimarisini Şekillendiren Etkenler}
\label{fig:mimari-etkenler}
\end{figure}

\begin{itemize}
    \item \textbf{Teknoloji:} Yeni üretim süreçleri ve malzeme bilimi ilerlemeleri, transistör boyutunu küçültüp işlem gücünü artırırken aynı zamanda güç tüketimi ve ısınma gibi yeni tasarım kısıtları getiriyor. Teknolojideki her sıçrama mimarinin yeniden düşünülmesini gerektiriyor.
    \item \textbf{Programlama dilleri:} Bilgisayarın programlama dillerini desteklemesi için işlemcinin, dilin gereksinim duyduğu işlemleri verimli biçimde gerçekleştiren buyrukları içerecek şekilde tasarlanması gerekiyor.
    \item \textbf{Uygulama:} Farklı kullanım alanları (bilimsel hesaplama, yapay zekâ, gömülü sistemler, taşınabilir uygulamalar) farklı mimari gereksinimlere yol açar. Alana özgü hızlandırıcıların ortaya çıkması bu etkinin somut bir sonucu.
    \item \textbf{İşletim sistemleri:} Donanımın, işletim sisteminin gereksinim duyduğu bellek koruması, bağlam değiştirme ve kesme mekanizmalarını desteklemesi gerekiyor. İşletim sistemi ile donanım birlikte evrilir.
    \item \textbf{Geçmiş (Tarih):} \emph{Geçmişe uyumluluk}\index{geçmişe uyumluluk} (\emph{backward compatibility}) ilkesi, daha önce yazılmış yazılımların yeni donanımda da çalışmasını zorunlu kılar. Bu ilke mimariyi tarihe bağlı kılar; x86 buyruk kümesinin 1978'den bu yana korunması bunun en bilinen örneği.
\end{itemize}

Bu etkenler arasında \emph{geçmişe uyumluluk} özellikle güçlü bir kısıt oluşturur. Intel, 8086'dan bu yana her yeni işlemcide önceki neslin tüm yazılımlarını çalıştırma taahhüdünü sürdürmüştür. Devasa bir yazılım ekosistemine sahip olmanın getirdiği bu rekabet avantajı, şirket içinde \emph{altın kelepçe}\index{altın kelepçe}\index{golden handcuffs|see{altın kelepçe}} (\emph{golden handcuffs}) olarak adlandırılır: milyarlarca dolarlık mevcut yazılım yatırımı müşterileri x86 ekosistemine bağlarken, aynı zamanda mimarların eski buyrukları ve davranışları sonsuza dek desteklemesini zorunlu kılar. Bu durum, teknik açıdan daha verimli tasarım seçeneklerini engellese bile ticari açıdan vazgeçilmez görülmektedir. Geçmişe uyumluluğun buyruk kümesi tasarımına etkileri Bölüm~\ref{chap:buyruk-kumesi}'te ayrıntılı olarak ele alınacaktır.

\subsection{Dört Tasarım İlkesi}
\label{sec:tasarim-ilkeleri}

Bilgisayar mimarisi alanında onlarca yıllık deneyimden süzülmüş dört temel tasarım ilkesi vardır. Bu ilkeleri kitap boyunca tekrar tekrar göreceğiz; burada kısaca tanıyalım.

\begin{enumerate}
    \item \emph{Küçük olan hızlıdır.}\index{tasarım ilkeleri!küçük olan hızlıdır} Bir birim ne kadar küçükse içindeki teller o kadar kısa, sinyal gecikmesi o kadar düşük olur. Yazmaç öbeği, önbellek, çoklayıcı: hepsi gereksinimi karşılayacak en küçük boyutta tutulmalıdır.
    \item \emph{Olağan durumu hızlandır.}\index{tasarım ilkeleri!olağan durumu hızlandır} Nadir gerçekleşen durumlar için karmaşık donanım eklemek yerine, en sık karşılaşılan durumu hızlı kılmak toplam başarıma çok daha fazla katkı sağlar. Bu ilke, Bölüm~\ref{chap:basarim}'de ele alacağımız Amdahl Yasası ile doğrudan bağlantılıdır.
    \item \emph{İyi tasarım ödünleşme gerektirir.}\index{tasarım ilkeleri!ödünleşme} Her tasarım kararının bir getirisi, bir de bedeli vardır. Daha geniş bir veri yolu aktarım hızını artırır ama daha fazla silikon alanı tüketir; daha büyük bir önbellek bulma oranını yükseltir ama erişim süresini uzatır. Burada ``en iyi'' diye mutlak bir seçenek yoktur; önemli olan gereksinime göre getiri--bedel çözümlemesi yaparak artıda kalmaktır.
    \item \emph{Yalınlık düzenden gelir.}\index{tasarım ilkeleri!yalınlık düzenden gelir} Aynı boyutta buyruklar, tutarlı alan düzenleri, az sayıda kural dışı durum: düzenli bir yapı donanımı yalınlaştırır. Yalın donanım daha küçük, daha hızlı ve daha az hata eğilimlidir.
\end{enumerate}

\begin{bilginot}[Tasarım İlkeleri Yalnızca Donanım İçin mi?]
Bu dört ilke yalnızca silikon üzerinde değil, aslında herhangi bir sistem tasarımında geçerlidir. Bir takım küçükse daha çevik hareket eder (küçük olan hızlıdır). Her gün yaptığınız bir işe yatırım yapmak, yılda bir yaptığınız işe yatırım yapmaktan çok daha verimlidir (olağan durumu hızlandır). Aldığınız her kararın bir artısı, bir de eksisi vardır; hiçbir seçenek mutlak iyi değildir (ödünleşme). İstisnası az olan düzenli bir sistem, yönetmesi ve uygulaması kolay bir sistemdir (yalınlık düzenden gelir).

Bunların hepsinden önce ise sıfırıncı adım gelir: \emph{gereksinim belirleme}. Bir işi neden yaptığınızı ortaya koymadan, ne kadar başarılı olduğunuzu ölçemezsiniz. ``En iyi uçak hangisidir?'' sorusunu düşünelim: F-16 mı, Boeing 747 mi? Yanıt, o uçakla ne yapmak istediğinize bağlıdır. Bomba atıp hızla kaçacaksanız F-16, yüzlerce yolcuyu ucuza uzak mesafeye taşıyacaksanız Boeing 747. Mühendislikte de hayatta da her şey gereksinimle başlar.\footnote{Bu konunun daha geniş bir tartışması için ``Mühendisliği Hayata Uygulamak'' başlıklı canlı yayını izleyebilirsiniz: \url{https://www.youtube.com/live/S85FnepzrGo}}
\end{bilginot}

\section{Güç Tüketimi Eğilimleri}
\label{sec:guc-tuketimi}

Sayısal devrelerde güç tüketimi\index{güç tüketimi} giderek daha fazla önem kazanan bir tasarım ölçütü. Taşınabilir bilgisayarların kullanımının artmasıyla önem kazanan pil ömrü, sayısal devrelerin daha az güç tüketecek biçimde tasarlanmasını gerektiriyor.

Bir işlemcinin güç tüketimini ifade etmek için yaygın olarak kullanılan ölçüt \textbf{TDP}\index{TDP} (\emph{Thermal Design Power}: Soğutma Tasarım Gücü) değeridir. TDP, işlemcinin tipik iş yükü altında ürettiği en yüksek ısı miktarını watt cinsinden belirtir ve soğutma sisteminin bu ısıyı uzaklaştırabilecek kapasitede tasarlanması gerektiğini gösterir. TDP değeri işlemcinin tüketebileceği mutlak en yüksek güç değil, soğutma tasarımı için referans alınan üst sınırdır. Tablo~\ref{tab:sogutma-sistemleri} farklı soğutma çözümlerinin kapasitesini ve yaklaşık maliyetini karşılaştırmaktadır.

\begin{table}[ht!]
\centering
\caption[Soğutma çözümlerinin TDP kapasitesi ve maliyeti]{İşlemci soğutma çözümlerinin yaklaşık ısı uzaklaştırma kapasitesi, maliyeti ve tipik kullanım alanları. Değerler tipik ürünlere göre verilmiştir; gerçek başarım kasaya, hava akışına ve ortam sıcaklığına bağlı olarak değişir.}
\label{tab:sogutma-sistemleri}
\footnotesize
\begin{tabular}{lrrp{3.6cm}p{3.2cm}}
\toprule
\textbf{Soğutma Türü} & \textbf{Kapasite (W)} & \textbf{Maliyet (\$)} & \textbf{Açıklama} & \textbf{Tipik Kullanım} \\
\midrule
Stok soğutucu          & 45--65    & 0 (dahil)    & Alüminyum blok + pervane & Ofis, web, hafif kullanım \\
Kule tipi hava (giriş) & 120--150  & 25--40       & Tek kule, 4 ısı borusu, pervane & Oyun, yazılım geliştirme \\
Kule tipi hava (üst)   & 200--250  & 80--110      & Çift kule, 6+ ısı borusu, 2 pervane & Üst seviye oyun, içerik üretimi \\
HBA sıvı (240\,mm)     & 200--250  & 70--120      & Kapalı döngü, 240\,mm radyatör & Oyun, iş istasyonu \\
HBA sıvı (360\,mm)     & 250--300  & 100--180     & Kapalı döngü, 360\,mm radyatör & Hız aşırtma, çok çekirdekli \\
Özel döngü sıvı        & 300+      & 300--600+    & Açık döngü, kullanıcı kurulumu & Aşırı hız aşırtma, çoklu GPU \\
\bottomrule
\end{tabular}
\end{table}

Tabloda geçen \textbf{HBA}\index{HBA} (Hepsi Bir Arada) terimi, pompası, radyatörü ve sıvı bağlantıları fabrikada birleştirilmiş kapalı döngü sıvı soğutma sistemlerini ifade eder. Kullanıcının sıvı doldurmasına veya bağlantı yapmasına gerek kalmaz; kutudan çıktığı gibi takılır. Hava soğutmalı çözümlerde ise ısı, metal kanatçıklar (\emph{heatsink}) aracılığıyla yayılır ve bir veya birden fazla \emph{pervane}\index{pervane} (\emph{fan}) ile üflenen hava akışıyla uzaklaştırılır.

TDP'si 65\,W olan bir işlemci stok soğutucuyla çalışabilirken, 125\,W ve üzeri TDP'ye sahip işlemciler kule tipi veya sıvı soğutma gerektirir. Günümüzde üst seviye masaüstü işlemcilerinin TDP değerleri 200\,W'ı aşabildiğinden, soğutma sistemi tasarımı ve maliyeti önemli bir mühendislik kısıtı hâline gelmiştir.

Şekil~\ref{fig:tdp-trend} Intel işlemcilerinin yıllara göre TDP değerlerini göstermektedir.%
\begin{figure}[H]
\centering
\begin{tikzpicture}[>=Stealth]
    % Sabitler
    \def\bw{0.35}    % çubuk yarı genişliği
    \def\sp{1.25}    % çubuklar arası mesafe
    \def\ysc{0.044}  % watt başına cm

    % Y ekseni
    \draw[thick, ->] (-0.3,0) -- (-0.3,{130*\ysc+0.3})
        node[above, font=\small\bfseries] {TDP (W)};
    \foreach \w in {0,25,50,75,100,125} {
        \draw (-0.3,{\w*\ysc}) -- +(-0.15,0);
        \node[left, font=\scriptsize] at (-0.45,{\w*\ysc}) {\w};
        \draw[gray!15, thin] (-0.3,{\w*\ysc}) -- ({11*\sp+0.8},{\w*\ysc});
    }

    % X ekseni
    \draw[thick] (-0.3,0) -- ({11*\sp+0.8},0);

    % Çubuklar — eşit aralıklı, idx=0..11
    % TDP'ye orantılı ısı gradyanı: düşük TDP → mavi (serin), yüksek TDP → kırmızı (sıcak)
    % idx / yıl / TDP / ad
    \foreach \idx/\yr/\tdp/\ad in {
        0/1971/1/4004,
        1/1978/3/8086,
        2/1989/5/80486,
        3/1993/15/Pentium,
        4/1995/30/Pentium Pro,
        5/1999/35/Pentium III,
        6/2000/75/Pentium 4,
        7/2004/115/Prescott,
        8/2006/65/Core 2 Duo,
        9/2011/95/Sandy Bridge,
        10/2015/91/Skylake,
        11/2021/125/Alder Lake} {
        % TDP'ye göre renk: 0 W → mavi (serin), 125 W → kırmızı (sıcak)
        \pgfmathtruncatemacro{\blueratio}{max(0, min(100, 100 - \tdp*100/125))}
        \fill[sinyalmavi!\blueratio!sinyalkirmizi, draw=gray!50, thin]
            ({\idx*\sp-\bw},0) rectangle ({\idx*\sp+\bw},{\tdp*\ysc});
        % Yıl (altta)
        \node[font=\scriptsize, below] at ({\idx*\sp},-0.08) {\yr};
        % İşlemci adı ve watt değeri (çubuğun üstünde, eğik, sola kaydırılmış)
        \node[font=\scriptsize, rotate=50, anchor=west, xshift=-9pt, align=center]
            at ({\idx*\sp},{\tdp*\ysc+0.25}) {\ad\\[-1pt]{\tiny\tdp\,W}};
    }

    % Güç duvarı oku (Prescott'tan Core 2 Duo'ya)
    \draw[sinyalkirmizi, very thick, ->, densely dashed]
        ({7.5*\sp},{110*\ysc}) -- ({8*\sp-\bw-0.05},{65*\ysc+0.15});
    \node[font=\footnotesize\bfseries, text=sinyalkirmizi, fill=white, inner sep=1pt]
        at ({8.2*\sp},{120*\ysc}) {Güç Duvarı};

\end{tikzpicture}
\caption[İşlemcilerin TDP değerlerinin tarihsel gelişimi]{Intel işlemcilerin yıllara göre TDP değerleri. Çubuk rengi TDP değeriyle orantılıdır: mavi serin (düşük güç), kırmızı sıcak (yüksek güç). 2004'teki Prescott'un 115\,W'a ulaşması güç duvarını somutlaştırmış, Core~2~Duo ile verimlilik odaklı tasarıma geçilmiştir.}
\label{fig:tdp-trend}
\end{figure}
%
Şekilden görüldüğü üzere, işlemcilerin güç tüketimi 1971'deki 4004'ün 1\,W altındaki değerinden 2004'te Pentium~4 Prescott ile 115\,W'a yükselmiştir. Intel'in 2006'da Core~2~Duo ile mimari değişikliğe gitmesi, saat sıklığı yerine verimliliğe odaklanıldığının açık göstergesidir.

\subsection{Devingen Güç Tüketimi}
\label{sec:devingen-guc}

Devingen güç tüketimi\index{devingen güç tüketimi}, mantık devrelerini oluşturan transistörlerin açılıp kapanarak işlem gerçekleştirdiği sırada ortaya çıkan güç tüketimidir:
\begin{equation}
P_{\text{devingen}} = C \times V_{dd}^2 \times f
\label{eq:devingen-guc}
\end{equation}
Burada $C$ sürülen yükü, $V_{dd}$ kaynak gerilimini, $f$ ise açılıp kapanma sıklığını gösterir. Denklemden görüldüğü üzere güç tüketimi kaynak gerilimiyle karesel orantılı olduğundan, $V_{dd}$'nin düşürülmesi güç tüketimini azaltmanın en etkili yoludur.

\begin{bilginot}[Kaynak Geriliminin Tarihsel Evrimi]
İlk mikroişlemcilerde (1970'ler--1980'ler) kaynak gerilimi 5{,}0~V idi. Üretim teknolojileri geliştikçe transistör boyutları küçüldü ve kaynak gerilimi de buna koşut olarak düşürüldü:
\begin{itemize}[nosep]
    \item 1{,}0~\textmu m teknoloji (1980'ler): $V_{dd} = 5{,}0$~V
    \item 0{,}35~\textmu m teknoloji (1990'lar): $V_{dd} = 3{,}3$~V
    \item 0{,}18~\textmu m teknoloji (2000'ler): $V_{dd} = 1{,}8$~V
    \item 65~nm teknoloji (2006): $V_{dd} = 1{,}2$~V
    \item 22~nm teknoloji (2012): $V_{dd} = 0{,}8$~V
    \item 5~nm teknoloji (2020'ler): $V_{dd} \approx 0{,}7$~V
\end{itemize}
Kaynak gerilimi 5{,}0~V'tan 0{,}7~V'a düştüğünde, yalnızca bu değişikliğin etkisiyle devingen güç tüketimi yaklaşık 50 kat azalır $(5^2 / 0{,}7^2 \approx 51)$. Ancak düşen gerilimle birlikte transistör eşik gerilimine yaklaşıldığından sızıntı akımları artar ve devrelerin gürültüye karşı dayanıklılığı azalır.
\end{bilginot}

Güç ($P$) birim zamanda harcanan enerji miktarını, yani anlık tüketimi, ifade ederken, \emph{enerji} ($E$) belirli bir süre boyunca harcanan toplam miktarı gösterir. Bir işlemcinin güç tüketimi yüksek olsa bile görevi kısa sürede bitirirse toplam enerji tüketimi düşük kalabilir; tam tersine, düşük güçlü bir işlemci uzun süre çalışarak daha fazla enerji harcayabilir. Taşınabilir aygıtlarda pil ömrünü belirleyen güç değil enerjidir; benzer biçimde veri merkezlerinin elektrik faturasını da toplam enerji tüketimi belirler.

Enerji tüketimi şu denklemle ifade edilir:
\begin{equation}
E = P \times t = C \times V_{dd}^2 \times f \times t
\label{eq:enerji}
\end{equation}
Burada $t$ işlemin süresini gösterir. Enerji tasarrufu hem güç tüketimini düşürerek hem de işlemi daha kısa sürede tamamlayarak sağlanabilir. Bu nedenle ``yavaş ama düşük güçlü'' bir işlemci her zaman ``hızlı ama yüksek güçlü'' bir işlemciden daha verimli değildir; enerji açısından her iki boyutun da birlikte değerlendirilmesi gerekir.

\begin{ornek}[1.5]
\label{orn:guc-hesaplama}
Bir işlemcinin kapasitif yükü $C = 20$\,nF, besleme gerilimi $V = 1{,}0$\,V ve saat sıklığı $f = 3$\,GHz ise:
\begin{enumerate}[(a)]
\item Devingen güç tüketimini hesaplayın.
\item Gerilim \%10 azaltılırsa (\,$V = 0{,}9$\,V) güç tüketimi ne kadar düşer?
\item Hem gerilim \%10 hem de saat sıklığı \%10 azaltılırsa güç tasarrufu ne olur?
\end{enumerate}

\textbf{Çözüm:}

(a) Devingen güç: $P = C \times V^2 \times f = 20 \times 10^{-9} \times 1{,}0^2 \times 3 \times 10^{9} = 60$\,W.

(b) $V = 0{,}9$\,V için: $P = 20 \times 10^{-9} \times 0{,}81 \times 3 \times 10^{9} = 48{,}6$\,W. Güç \%19 azalmıştır. Gerilim yalnızca \%10 düşürülmesine karşın güç \%19 azalmıştır; bunun nedeni güç tüketiminin gerilimin \emph{karesiyle} orantılı olmasıdır.

(c) $V = 0{,}9$\,V, $f = 2{,}7$\,GHz: $P = 20 \times 10^{-9} \times 0{,}81 \times 2{,}7 \times 10^{9} = 43{,}74$\,W. Toplam güç tasarrufu: $(60 - 43{,}74)/60 = \%27{,}1$.

Bu örnek, gerilim düşürmenin (\emph{voltage scaling}) neden en etkili güç azaltma yöntemi olduğunu açıkça göstermektedir. Dennard ölçeklemesi sona erince transistör boyutu küçülse bile gerilim artık orantılı biçimde düşürülemez hale geldi; bu durum güç duvarını ortaya çıkardı.
\end{ornek}


\subsection{Durağan Güç Tüketimi}
\label{sec:duragan-guc}

Mantık devrelerini oluşturan transistörlerin hiçbir anlamlı iş yapmadıkları halde, kaynak geriliminden toprağa akım sızdırmaları nedeniyle oluşan güç tüketimine \emph{durağan güç tüketimi}\index{durağan güç tüketimi} denir:
\begin{equation}
P_{\text{durağan}} = I_{\text{sızıntı}} \times V_{dd}
\label{eq:duragan-guc}
\end{equation}

Transistör boyutları küçüldükçe sızıntı akımı artar. Yükselen transistör sayıları ve düşen transistör boyutları durağan güç tüketimini giderek artırıyor. Tablo~\ref{tab:sizinti-trend} durağan güç tüketiminin toplam güç içindeki payının teknoloji düğümleriyle nasıl değiştiğini göstermektedir.

\begin{table}[ht!]
\centering
\caption[Durağan güç oranının teknoloji düğümlerine göre değişimi]{Durağan (sızıntı) güç tüketiminin toplam işlemci güç tüketimi içindeki yaklaşık payı.}
\label{tab:sizinti-trend}
\footnotesize
\begin{tabular}{lrl}
\toprule
\textbf{Teknoloji Düğümü} & \textbf{Sızıntı Payı (\%)} & \textbf{Dönem} \\
\midrule
250\,nm  & 2--5    & 1990'ların sonu \\
130\,nm  & ${\sim}$10   & 2001 \\
90\,nm   & 20--25  & 2004 \\
65\,nm   & 30--35  & 2006 \\
45\,nm   & 35--40  & 2008 \\
32\,nm   & 40--50  & 2010 \\
22\,nm   & 35--40  & 2012 \\
\bottomrule
\multicolumn{3}{l}{\scriptsize 22\,nm'deki düşüş FinFET teknolojisinin etkisini yansıtır.}
\end{tabular}
\end{table}

Tablodan görüldüğü üzere, 250\,nm'den 32\,nm'ye geçişte sızıntı akımının toplam güç içindeki payı yaklaşık 10 kat artmıştır. Ancak Intel'in 22\,nm düğümünde \textbf{FinFET}\index{FinFET} (üç boyutlu kapılı, \emph{Fin Field-Effect Transistor}) yapısına geçmesiyle sızıntı akımı kısmen kontrol altına alınmıştır. Günümüzdeki ileri düğümlerde (5\,nm ve altı) \textbf{GAA}\index{GAA} (\emph{Gate-All-Around}) yapılar sızıntıyı daha da azaltmayı hedeflemektedir.


\subsection{Güç Tüketimini Azaltma Yöntemleri}
\label{sec:guc-azaltma}

Güç tüketiminin birimi Watt'tır. Saat sıklığı giderek yükselen işlemcilerdeki güç tüketimini azaltmanın en kolay yolu kaynak gerilimini düşürmek. Ancak kaynak geriliminin düşmesi devre gecikmelerini ve saklanan elektrik sinyallerinin gürültüden etkilenme olasılığını artırır.

Güç tüketimini azaltmak için kullanılan bir diğer yöntem ise kullanılmayan birimlerin kaynak gerilimini kesmek ya da saat sıklığını düşürmek. Güç tüketimi işlemci tasarımcıları tarafından tasarımın en önemli kısıtlarından biri olarak görüldüğü için aşılması gereken bu engele \emph{Güç Duvarı}\index{güç duvarı} adı verildi. Tek bir çekirdekle çok yüksek oranda işlem yapılmasını kısıtlayan bu güç duvarı, tasarımcıları aynı anda daha fazla işlem yapılmasını sağlayan çok çekirdekli\index{çok çekirdekli} tasarımlara yönlendirdi. Güç tüketimi kısıtının boru hattı derinliğini ve işlemci mikromimarisini nasıl biçimlendirdiğini Bölüm~\ref{chap:islemci-tasarimi}'da ve Bölüm~\ref{chap:boruhatti}'de ayrıntılı olarak göreceğiz.

\begin{bilginot}[GHz Yarışından Verimlilik Yarışına]
2000'li yılların başında işlemci üreticileri saat sıklığını artırarak başarım kazanımı sağlama yarışına girmişti. Ancak Pentium~4 serisiyle saat sıklığı 3{,}8~GHz'e ulaştığında güç tüketimi ve ısınma sorunları bu yaklaşımın sürdürülemez olduğunu gösterdi. Günümüzde başarım ölçütü olarak ``performans/watt'' birimi öne çıkıyor. ARM\index{ARM} tabanlı işlemcilerin taşınabilir aygıtlarda ve hatta sunucularda (ör.\ AWS Graviton) yaygınlaşması, güç verimli tasarımın öneminin somut bir göstergesi. RISC-V'in modüler yapısı da tasarımcılara güç bütçesine göre özelleştirilmiş çekirdekler oluşturma esnekliği sunuyor.
\end{bilginot}

\begin{ornek}[1.6]
\label{orn:dennard-guc-duvari}
Dennard ölçeklemesi\index{Dennard ölçeklemesi} döneminde transistör boyutu küçüldükçe kaynak gerilimi ve saat sıklığı da buna uyumlu biçimde değişiyordu. Aşağıdaki iki kuşak işlemciyi karşılaştırarak güç duvarının kökenini nicel olarak gösterin.

\begin{center}
\small
\begin{tabular}{lcc}
\toprule
& \textbf{Kuşak~1 (130\,nm)} & \textbf{Kuşak~2 (65\,nm)} \\
\midrule
Transistör boyutu ($s$)  & 1 (referans) & $0{,}5\times$ \\
Kaynak gerilimi ($V_{dd}$) & 1{,}2\,V & ? \\
Saat sıklığı ($f$) & 2\,GHz & ? \\
Kapasitif yük ($C$) & 20\,nF & ? \\
\bottomrule
\end{tabular}
\end{center}

\begin{enumerate}[(a)]
\item Dennard ölçeklemesi geçerli olsaydı Kuşak~2'nin $V_{dd}$, $f$ ve $C$ değerleri ne olurdu? Devingen güç tüketimi nasıl değişirdi?
\item Uygulamada Kuşak~2 işlemcisinin gerilimi yalnızca 1{,}0\,V'a düşürülebildi ($0{,}83\times$), saat sıklığı ise 3\,GHz'e ($1{,}5\times$) çıkarıldı. Kapasitif yük $0{,}5\times$ oranında küçüldü. Gerçek güç tüketimini hesaplayın ve ideal durumla karşılaştırın.
\end{enumerate}

\textbf{Çözüm:}

(a) Dennard ölçeklemesinde boyut $s$ kat küçüldüğünde: gerilim $s$ kat düşer, saat sıklığı $1/s$ kat artar, kapasitif yük $s$ kat azalır. $s = 0{,}5$ için:

$V_{dd} = 1{,}2 \times 0{,}5 = 0{,}6$\,V, \quad $f = 2 \times 2 = 4$\,GHz, \quad $C = 20 \times 0{,}5 = 10$\,nF.

İdeal güç: $P = C \times V_{dd}^2 \times f = 10 \times 10^{-9} \times 0{,}6^2 \times 4 \times 10^{9} = 14{,}4$\,W.

Kuşak~1'in gücü: $P_1 = 20 \times 10^{-9} \times 1{,}2^2 \times 2 \times 10^{9} = 57{,}6$\,W.

İdeal ölçeklemede güç \textbf{$14{,}4 / 57{,}6 = 0{,}25 = s^2$} oranında düşer. Transistör sayısı 4 kat artmasına karşın güç aynı kalır; Dennard ölçeklemesinin büyüsü budur.

(b) Gerçek değerlerle: $P_2 = 10 \times 10^{-9} \times 1{,}0^2 \times 3 \times 10^{9} = 30$\,W.

\begin{center}
\small
\begin{tabular}{lccc}
\toprule
& \textbf{İdeal (Dennard)} & \textbf{Gerçek} & \textbf{Fark} \\
\midrule
Gerilim & 0{,}6\,V & 1{,}0\,V & $1{,}67\times$ yüksek \\
Saat sıklığı & 4\,GHz & 3\,GHz & $0{,}75\times$ düşük \\
Güç tüketimi & 14{,}4\,W & 30\,W & $2{,}08\times$ yüksek \\
\bottomrule
\end{tabular}
\end{center}

Gerilimin ideal değerin $1{,}67$ katında kalması, güç tüketimini $1{,}67^2 \approx 2{,}8$ kat artırır. Saat sıklığının düşük kalması bunu kısmen telafi etse de sonuçta güç tüketimi idealin 2 katını aşar. Her yeni kuşakta bu fark birikir ve sonunda güç duvarına ulaşılır. 65\,nm sonrasında Dennard ölçeklemesinin çöküşü, endüstriyi tek çekirdekli hız artışından çok çekirdekli tasarımlara yönelmeye zorlamıştır.
\end{ornek}

%% ============================================================
%% BÖLÜM 1.8 — GÜVENİLİRLİK
%% ============================================================
\section{Güvenilirlik}
\label{sec:guvenilirlik}

Bir bilgisayar sisteminden beklenen yalnızca hızlı çalışması değil, \emph{doğru} ve \emph{kesintisiz} çalışmasıdır. Bir süper bilgisayar saniyede katrilyon işlem yapabilir ama hesaplama sonuçları yanlışsa ya da sistem saatte bir çökerse hiçbir kullanıcıya faydası yoktur. Güç tüketimi gibi güvenilirlik de günümüzde başarım kadar önemli bir tasarım ölçütüdür; özellikle sunucu, veri merkezi, uzay ve otomotiv uygulamalarında.

\subsection{Arıza, Hata ve Bozukluk}
\label{subsec:ariza-hata-bozukluk}

Güvenilirlik çözümlemesinde üç kavram dikkatli biçimde birbirinden ayrılmalıdır:

\begin{enumerate}
    \item \textbf{Bozukluk}\index{bozukluk (\emph{fault})} (\emph{fault}): Sistemin bir bileşenindeki fiziksel veya mantıksal kusur. Örneğin bir bellek hücresine kozmik ışının çarpması ya da bir lehim noktasının çatlaması bir bozukluktur.
    \item \textbf{Hata}\index{hata (\emph{error})} (\emph{error}): Bozukluğun sisteme yansıması. Kozmik ışının çarptığı bellek hücresindeki bitin 0'dan 1'e dönmesi bir hatadır.
    \item \textbf{Arıza}\index{arıza (\emph{failure})} (\emph{failure}): Hatanın kullanıcı tarafından gözlemlenebilir bir yanlış sonuca veya hizmet kesintisine dönüşmesi. Bellekteki hatalı bit bir program çökmesine neden olursa bu bir arızadır.
\end{enumerate}

Bu zincir şöyle özetlenebilir:
%
\[
\text{Bozukluk} \;\longrightarrow\; \text{Hata} \;\longrightarrow\; \text{Arıza}
\]
%
Her bozukluk hataya, her hata arızaya dönüşmek zorunda değildir. \emph{Hata düzeltme kodlaması}\index{hata düzeltme kodu}\index{ECC|see{hata düzeltme kodu}} (\emph{ECC}: Error Correcting Code) hatayı arızaya dönüşmeden yakalayabilir; yedekli hesaplama ise bir birimin bozukluğunu diğerleriyle maskeleyebilir. Güvenilirlik mühendisliğinin amacı bu zinciri mümkün olan en erken noktada kırmaktır.

ECC'nin temel çalışma ilkesi şudur: veri saklanırken veya iletilirken sözcüğe fazladan \emph{denetim bitleri} eklenir. Bu bitler, belirli matematiksel kodlama kurallarına göre üretilir ve okuma sırasında denetlenir. Tek bitlik bir hata oluşmuşsa ECC devresi hatayı hem \emph{saptayıp} hem de \emph{düzeltebilir}; çift bitlik hatalar ise saptanabilir ama düzeltilemez (SECDED\index{SECDED}: \emph{Single Error Correction, Double Error Detection}). Bu mekanizma özellikle sunucu belleklerinde yaygın olarak kullanılır ve kozmik ışın kaynaklı geçici bozuklukların büyük çoğunluğunu kullanıcı farkına bile varmadan giderir. Hata düzeltme kodlamasının ayrıntılı matematiksel temeli ve farklı kodlama yöntemleri Bölüm~\ref{chap:bellek}'de ele alınacaktır.

\subsection{Bozukluk Türleri}
\label{subsec:bozukluk-turleri}

Bozukluklar süresine göre üç gruba ayrılır:
\begin{itemize}
    \item \textbf{Geçici bozukluk}\index{geçici bozukluk} (\emph{transient fault}): Anlık bir olay sonucu oluşur ve kendiliğinden kaybolur. Kozmik ışının neden olduğu bit hataları, gerilim dalgalanmaları ve elektromanyetik girişim başlıca örnekleridir. Uzay uygulamalarında radyasyon kaynaklı geçici bozukluklar özellikle sık görülür.
    \item \textbf{Aralıklı bozukluk}\index{aralıklı bozukluk} (\emph{intermittent fault}): Belirli koşullar altında tekrar tekrar ortaya çıkan ama koşullar değişince kaybolan bozukluktur. Yüksek sıcaklıkta temas kaybeden bir lehim noktası veya belirli bir veri örüntüsünde hatalı davranan bir bellek hücresi bu türe örnektir. Tanılanması en zor bozukluk türüdür.
    \item \textbf{Kalıcı bozukluk}\index{kalıcı bozukluk} (\emph{permanent fault}): Bir bileşenin geri dönüşü olmayan biçimde hasar görmesi. Yanmış bir transistör, aşınmış bir disk yüzeyi veya kopmuş bir bağlantı kalıcı bozukluklara örnektir.
\end{itemize}

\subsection{Güvenilirlik Ölçütleri}
\label{subsec:guvenilirlik-olcutleri}

Bir sistemin güvenilirliğini nicel olarak ifade etmek için üç temel ölçüt kullanılır:

\begin{itemize}
    \item \textbf{MTTF}\index{MTTF} (\emph{Mean Time To Failure}: Ortalama Arızaya Kadar Geçen Süre): Onarılamayan bir bileşenin çalışmaya başlamasından ilk arızaya kadar geçen ortalama süredir.

    \item \textbf{MTTR}\index{MTTR} (\emph{Mean Time To Repair}: Ortalama Onarım Süresi): Arıza gerçekleştikten sonra sistemin onarılıp yeniden çalışır duruma getirilmesi için geçen ortalama süredir.

    \item \textbf{MTBF}\index{MTBF} (\emph{Mean Time Between Failures}: Ortalama Arızalar Arası Süre): Onarılabilir bir sistemde iki ardışık arıza arasında geçen ortalama süredir:
\end{itemize}
%
\begin{equation}
\text{MTBF} = \text{MTTF} + \text{MTTR}
\label{eq:mtbf}
\end{equation}

Bu ölçütlerden türetilen iki önemli kavram daha vardır. Birincisi \textbf{kullanılabilirlik}\index{kullanılabilirlik} (\emph{availability}), sistemin çalışır durumda olduğu zamanın toplam zamana oranıdır:
%
\begin{equation}
\text{Kullanılabilirlik} = \frac{\text{MTTF}}{\text{MTTF} + \text{MTTR}} = \frac{\text{MTTF}}{\text{MTBF}}
\label{eq:kullanilabilirlik}
\end{equation}
%
Örneğin MTTF'si 999~saat ve MTTR'si 1~saat olan bir sistemin kullanılabilirliği $999/1000 = \%99{,}9$ olur. Veri merkezlerinde kullanılabilirlik hedefi genellikle ``beş dokuz'' (\%99,999) olarak ifade edilir; bu, yılda en fazla yaklaşık 5~dakika kesinti anlamına gelir.

İkincisi \textbf{FIT}\index{FIT} (\emph{Failures In Time}), bir milyar aygıt-saatinde beklenen arıza sayısıdır:
%
\begin{equation}
\text{FIT} = \frac{10^9}{\text{MTTF (saat cinsinden)}}
\label{eq:fit}
\end{equation}
%
FIT birimi özellikle yarıiletken endüstrisinde yaygın olarak kullanılır. Bir bellek yongasının FIT değeri 100 ise, bu yonga bir milyar saat çalışmada ortalama 100~kez arıza verir; eşdeğer olarak, 10 milyon saat çalışmada 1~kez arıza beklenir.

Birden fazla bileşenden oluşan bir sistemde, herhangi bir bileşenin arızası tüm sistemi durduruyorsa (seri bağlantı), sistemin MTTF'si her bir bileşenin arıza oranlarının toplamının tersidir:
%
\begin{equation}
\text{MTTF}_{\text{sistem}} = \frac{1}{\displaystyle\sum_{i=1}^{N} \dfrac{1}{\text{MTTF}_i}} = \frac{1}{\displaystyle\sum_{i=1}^{N} \lambda_i}
\label{eq:mttf-sistem}
\end{equation}
%
Burada $\lambda_i = 1/\text{MTTF}_i$ bileşen $i$'nin arıza oranıdır. $N$ özdeş bileşenden oluşan bir sistemde bu ifade $\text{MTTF}_{\text{sistem}} = \text{MTTF}_{\text{bileşen}} / N$ biçimine sadeleşir. Bu denklem can alıcı bir sonuç doğurur: bir veri merkezinde 10\,000 sunucu varsa ve her birinin MTTF'si 100\,000~saat ise, merkezin MTTF'si yalnızca 10~saat olur; ortalama her 10 saatte bir sunucu arızası beklenir.

\begin{ornek}[1.7]
Bir sunucu üç temel bileşenden oluşuyor: işlemci, bellek ve disk. Her bileşenin MTTF değeri şöyledir:
\begin{center}
\small
\begin{tabular}{@{}lcc@{}}
\toprule
\textbf{Bileşen} & \textbf{MTTF (saat)} & \textbf{FIT} \\
\midrule
İşlemci  & 500\,000 & 2\,000 \\
Bellek   & 200\,000 & 5\,000 \\
Disk     & 100\,000 & 10\,000 \\
\bottomrule
\end{tabular}
\end{center}

\noindent Herhangi bir bileşenin arızası sunucuyu durduruyor (seri bağlantı). Ortalama onarım süresi (MTTR) 2~saattir.\\[4pt]
(a) Sunucunun MTTF'sini hesaplayın.\\
(b) FIT değerini bulun.\\
(c) Kullanılabilirliği (\%) hesaplayın.\\
(d) Aynı sunucudan 500 tane bulunan bir veri merkezinin MTTF'sini bulun.\\[6pt]
\textbf{Çözüm:}\\[2pt]
\textbf{(a)} Seri bağlantıda toplam arıza oranı:
\[
\lambda_{\text{toplam}} = \frac{1}{500\,000} + \frac{1}{200\,000} + \frac{1}{100\,000} = 2 + 5 + 10 = 17 \times 10^{-6}\;\text{arıza/saat}
\]
\[
\text{MTTF}_{\text{sunucu}} = \frac{1}{\lambda_{\text{toplam}}} = \frac{1}{17 \times 10^{-6}} \approx 58\,824\;\text{saat} \approx 6{,}7\;\text{yıl}
\]
Dikkat: en zayıf halka olan disk (100\,000~saat) sunucunun MTTF'sini büyük ölçüde belirliyor.

\textbf{(b)} FIT değeri:
\[
\text{FIT} = \frac{10^9}{\text{MTTF}} = \frac{10^9}{58\,824} \approx 17\,000
\]
Bu değer, bileşenlerin FIT toplamına ($2\,000 + 5\,000 + 10\,000 = 17\,000$) eşittir; seri bağlantıda FIT değerleri doğrudan toplanır.

\textbf{(c)} Kullanılabilirlik:
\[
\text{MTBF} = \text{MTTF} + \text{MTTR} = 58\,824 + 2 = 58\,826\;\text{saat}
\]
\[
\text{Kullanılabilirlik} = \frac{\text{MTTF}}{\text{MTBF}} = \frac{58\,824}{58\,826} \approx \%99{,}9966
\]
Bu ``dört dokuz''un üzerinde ama ``beş dokuz'' (\%99,999) hedefinin altında bir değerdir.

\textbf{(d)} Veri merkezinin MTTF'si:
\[
\text{MTTF}_{\text{merkez}} = \frac{\text{MTTF}_{\text{sunucu}}}{500} = \frac{58\,824}{500} \approx 117{,}6\;\text{saat} \approx 4{,}9\;\text{gün}
\]
Veri merkezinde ortalama her 5~günde bir sunucu arızası beklenir. Bu nedenle büyük veri merkezleri yedekleme ve otomatik yük aktarma mekanizmaları olmadan çalışamaz.
\end{ornek}

\subsection{Güvenilirlik, Dayanıklılık ve Gürbüzlük}
\label{subsec:guvenilirlik-dayaniklilik-gurbuzluk}

Günlük dilde birbiriyle karıştırılan üç kavramı kesin biçimde ayırmak gerekir:
\begin{itemize}
    \item \textbf{Güvenilirlik}\index{güvenilirlik} (\emph{reliability}): Sistemin belirli bir süre boyunca doğru çalışma olasılığı. MTTF ve MTBF ile ölçülür. ``Bu sunucu iki yıl arızasız çalışır mı?'' sorusunun yanıtıdır.
    \item \textbf{Dayanıklılık}\index{dayanıklılık} (\emph{durability}): Fiziksel aşınma ve çevresel koşullara direnme kapasitesi. Bir SSD'nin yazma dayanıklılığı (ör.\ 600\,TBW: Terabytes Written) veya bir askeri bilgisayarın titreşim dayanımı bu kavramla ifade edilir.
    \item \textbf{Gürbüzlük}\index{gürbüzlük} (\emph{robustness}): Beklenmedik koşullar altında kabul edilebilir biçimde çalışmaya devam edebilme yeteneği. Gerilim dalgalanmalarında veri kaybetmeyen bir dosya sistemi veya geçersiz girdide çökmek yerine hata iletisi veren bir yazılım gürbüzdür.
\end{itemize}

Bu üç kavram birbirini tamamlar: güvenilir bir sistem uzun süre doğru çalışır; dayanıklı bir sistem fiziksel koşullara direnir; gürbüz bir sistem beklenmedik durumlarla başa çıkar. İyi bir tasarım üçünü birden gözetir. Güvenilirliği donanım düzeyinde destekleyen önbellek ECC bitleri ve hata düzeltme mekanizmalarını Bölüm~\ref{chap:bellek}'de, giriş/çıkış sistemlerinin güvenilirlik tasarımını ise Bölüm~\ref{chap:giris-cikis}'da ele alacağız.

\begin{bilginot}[Mars'ta Güvenilirlik]
NASA'nın Mars keşif aracı \emph{Curiosity}\index{Curiosity}, radyasyona dayanıklı\index{radyasyona dayanıklılık} (\emph{radiation-hardened}) bir RAD750\index{RAD750} işlemcisi kullanır. Bu işlemci yalnızca 200~MHz'de çalışır; dünya üzerindeki herhangi bir akıllı telefondan yüzlerce kat yavaştır. Ancak Mars yüzeyinde kozmik ışınlara ve aşırı sıcaklık değişimlerine ($-$130°C ile $+$70°C arası) yıllarca dayanabilir. Uzay mühendisliğinde başarım ölçütü ``saniye başına buyruk'' değil, ``milyonlarca kilometre uzakta on yıl boyunca hatasız çalışma''dır.
\end{bilginot}


%% ============================================================
%% BÖLÜM 1.9 — NEDEN RISC-V?
%% ============================================================
\section{Neden RISC-V?}
\label{sec:neden-riscv}

Günümüzde masaüstü ve sunucu pazarına Intel ile AMD'nin \textbf{x86}\index{x86} mimarisi, taşınabilir ve gömülü pazara ise ARM Holdings'in \textbf{ARM}\index{ARM} mimarisi hâkim. Her iki buyruk kümesi mimarisi de tescilli (proprietary) olup kullanımı lisans anlaşmalarına ve ücretlere bağlı. Bu durum yeni işlemci tasarlamak isteyen girişimciler, araştırmacılar ve ülkeler için yüksek bir giriş bariyeri oluşturuyor.

2010 yılında Kaliforniya Üniversitesi Berkeley'de Krste Asanovi\'{c} ve David Patterson\index{Patterson, David} öncülüğünde geliştirilen \textbf{RISC-V}\index{RISC-V}, bu soruna köklü bir çözüm sunuyor. RISC-V, herhangi bir lisans ücreti gerektirmeyen, tamamen açık kaynaklı bir buyruk kümesi mimarisi. Başlangıçta RISC-V Foundation tarafından yönetilen proje, 2020'de \emph{RISC-V International} çatısı altında küresel bir standart haline geldi.

RISC-V'in öne çıkan özellikleri şöyle sıralanabilir:
\begin{itemize}
    \item \textbf{Açık kaynak ve ücretsiz:} Bir üniversite öğrencisinden çok uluslu bir şirkete kadar herkes RISC-V tabanlı işlemci tasarlayıp üretebilir.
    \item \textbf{Modüler tasarım:} Temel tamsayı buyruk kümesi (RV32I/RV64I) üzerine çarpma (M), atomik (A), kayan nokta (F, D) ve vektör (V) gibi isteğe bağlı uzantılar eklenebilir. Bu modülerlik, gömülü bir denetleyiciden süper bilgisayar çekirdeğine kadar geniş bir yelpazede özelleştirme olanağı tanır.
    \item \textbf{Temiz ve çağdaş:} 1970'lerin geriye uyumluluk yükünden arınmış, günümüz derleyici ve işletim sistemi gereksinimlerine göre tasarlanmış yalın bir mimari.
    \item \textbf{Eğitim dostu:} Karmaşık tarihsel eklentiler içermediği için öğrencilerin bir işlemciyi sıfırdan tasarlayıp benzetmesi oldukça kolaylaşıyor.
\end{itemize}

\begin{figure}[H]
\centering
\begin{tikzpicture}[>=Stealth,
    baslik/.style={font=\small\bfseries, minimum height=0.7cm, minimum width=3.2cm, align=center},
    hucre/.style={font=\footnotesize, minimum height=0.6cm, minimum width=3.2cm, align=center, draw, thin}]

    % Sütun başlıkları (marka renkleri: Intel mavi, ARM turuncu, RISC-V yeşil)
    \node[baslik, fill=sinyalmavi!70, text=white] (hx86) at (0, 0)   {x86};
    \node[baslik, fill=orange!70]   (harm) at (3.5, 0)   {ARM};
    \node[baslik, fill=sinyalyesil!70, text=white]  (hrv)  at (7, 0)     {RISC-V};

    % Satır başlıkları (sol taraf)
    \def\satirlar{{"Lisanslama","Kullanım Alanı","Genişletilebilirlik","Tasarım Maliyeti"}}

    \foreach \i/\lab in {1/Lisanslama, 2/Kullanım Alanı, 3/Genişletilebilirlik, 4/Tasarım Maliyeti} {
        \node[font=\footnotesize\itshape, anchor=east] at (-1.7, {-\i*0.7}) {\lab};
    }

    % x86 sütunu (Intel mavi)
    \node[hucre, fill=sinyalmavi!15] at (0, -0.7)  {Tescilli};
    \node[hucre, fill=sinyalmavi!15] at (0, -1.4)  {Masaüstü, Sunucu};
    \node[hucre, fill=sinyalmavi!15] at (0, -2.1)  {Sabit (eski yük)};
    \node[hucre, fill=sinyalmavi!15] at (0, -2.8)  {Yüksek};

    % ARM sütunu (turuncu)
    \node[hucre, fill=orange!15] at (3.5, -0.7)  {Tescilli};
    \node[hucre, fill=orange!15] at (3.5, -1.4)  {Mobil, Gömülü};
    \node[hucre, fill=orange!15] at (3.5, -2.1)  {Kısmi};
    \node[hucre, fill=orange!15] at (3.5, -2.8)  {Orta};

    % RISC-V sütunu (yeşil)
    \node[hucre, fill=sinyalyesil!15] at (7, -0.7)  {Açık Kaynak};
    \node[hucre, fill=sinyalyesil!15] at (7, -1.4)  {Giderek genişleyen};
    \node[hucre, fill=sinyalyesil!15] at (7, -2.1)  {Tam modüler};
    \node[hucre, fill=sinyalyesil!15] at (7, -2.8)  {Ücretsiz};

\end{tikzpicture}
\caption{x86, ARM ve RISC-V buyruk kümesi mimarilerinin karşılaştırması}
\label{fig:isa-karsilastirma}
\end{figure}

\begin{karsilastirma}[BKM Karşılaştırması: Kullanım Alanları]
Günümüzde üç büyük buyruk kümesi mimarisi farklı alanlarda baskın konumdadır:

\begin{center}
\small
\begin{tabular}{lp{3cm}p{3cm}p{3cm}}
\toprule
 & \textbf{x86-64} & \textbf{ARM} & \textbf{RISC-V} \\
\midrule
Baskın alan & Masaüstü, sunucu & Mobil, gömülü & Gömülü, araştırma \\
Lisans & Kapalı (Intel/AMD) & Ücretli lisans & Açık ve ücretsiz \\
Buyruk boyutu & Değişken (1--15\,B) & Sabit (32\,bit) & Sabit (32\,bit) \\
Geriye uyumluluk & 1978'den beri & AArch32/AArch64 & Yük yok \\
Pazar payı (2025) & $\sim$\%85 masaüstü & $\sim$\%95 mobil & Hızla büyüyor \\
\bottomrule
\end{tabular}
\end{center}

RISC-V'in en büyük avantajı, tarihsel yükten arınmış olarak sıfırdan tasarlanmış olmasıdır. x86'nın 1978'den bu yana sürdürmek zorunda olduğu geriye uyumluluk, mimariyi karmaşıklaştırmış ve güç tüketimini artırmıştır. ARM, taşınabilir pazarda hegemonya kurmuş olmakla birlikte lisans modeli küçük girişimler için engel oluşturabilmektedir. RISC-V ise açık kaynak modeli sayesinde hem üniversiteler hem de şirketler için erişilebilir bir seçenektir.
\end{karsilastirma}

Şekil~\ref{fig:isa-karsilastirma}'te üç önemli buyruk kümesi mimarisinin temel boyutlarda karşılaştırması veriliyor. Günümüzde SiFive, Alibaba T-Head, Andes ve NVIDIA gibi firmalar RISC-V tabanlı ticari işlemciler üretiyor. Avrupa İşlemci Girişimi (European Processor Initiative) de yüksek başarımlı hesaplama alanında RISC-V'e yatırım yapıyor. Akademik dünyada ise BOOM, Rocket ve CVA6 gibi açık kaynaklı RISC-V çekirdekleri araştırmacılara gerçek bir işlemci tasarımı üzerinde deney yapma olanağı sunuyor.

RISC-V'in benimsenmesi yalnızca akademik dünyayla sınırlı kalmadı; endüstride de hızla ivme kazandı. Şekil~\ref{fig:riscv-buyume}'te birikimli RISC-V çekirdek sayısının yıllara göre artışı gösterilmektedir. ARM mimarisinin 10~milyar çekirdeğe ulaşması 17~yıl sürmüşken, RISC-V bu eşiğe yalnızca 12~yılda ulaştı. Dikkat çekici bir nokta: NVIDIA tek başına 2024'te 1~milyar RISC-V çekirdeğini piyasaya sürdü; bunlar GPU denetleyicileri, güç yönetimi ve güvenlik alt sistemleri gibi görevlerde kullanılmaktadır.

\begin{figure}[H]
\centering
\begin{tikzpicture}
\begin{axis}[
    width=0.85\textwidth,
    height=7.5cm,
    ybar,
    bar width=0.55cm,
    ymin=0, ymax=25,
    ylabel={Birikimli çekirdek (milyar)},
    ylabel style={font=\small},
    symbolic x coords={2018,2019,2020,2021,2022,2023,2024,2025},
    xtick=data,
    xticklabel style={font=\small, yshift=-2pt},
    yticklabel style={font=\small},
    ytick={0,5,10,15,20,25},
    nodes near coords,
    nodes near coords style={font=\scriptsize\bfseries},
    every node near coord/.append style={yshift=5pt},
    enlarge x limits=0.1,
    grid=major,
    grid style={gray!20},
    axis line style={thick},
    tick style={thick},
]
\addplot[fill=sinyalyesil!40, draw=sinyalyesil!80!black] coordinates {
    (2018,1) (2019,2) (2020,4) (2021,7) (2022,10) (2023,13) (2024,17) (2025,20)
};
\end{axis}
% Açıklama notu
\node[font=\scriptsize, text=gray!70, anchor=north west] at (0.3,-0.8)
    {Kaynak: RISC-V International, Semico Research (2022--2025 verileri; 2025 tahmini)};
\end{tikzpicture}
\caption[RISC-V çekirdek sayısının yıllara göre büyümesi]{Piyasaya sürülen birikimli RISC-V çekirdek sayısı (milyar). 2022'de 10~milyar eşiği aşıldı; 2025 itibarıyla tahmini 20~milyar çekirdeğe ulaşıldı.}
\label{fig:riscv-buyume}
\end{figure}

\begin{bilginot}[Meta ve RISC-V: Veri Merkezine Giden Yol]
RISC-V'in yalnızca gömülü sistemler ve mikrodenetleyiciler için bir mimari olduğu algısı, son yıllarda hızla değişmektedir. Meta\index{Meta} (Facebook), Eylül 2025'te RISC-V tabanlı sunucu ve yapay zekâ yongaları geliştiren \emph{Rivos}\index{Rivos} girişimini satın aldı. Rivos, 3{,}1\,GHz'lik 64-bit RISC-V işlemci çekirdekleri ve büyük dil modelleri için eniyilenmiş özel bir GPU tasarlıyordu. Meta'nın bu hamlesi, daha önce Amazon'un Annapurna Labs'ı\index{Annapurna Labs} alarak ARM tabanlı Graviton\index{Graviton} işlemcilerini geliştirmesine benzetilmektedir; ancak burada tercih edilen mimari ARM değil, RISC-V'tir. Meta'nın yapay zekâ çıkarım hızlandırıcısı MTIA\index{MTIA} zaten RISC-V tabanlı işlem elemanları kullanmaktaydı; Rivos satın almasıyla şirket, veri merkezlerinde RISC-V'i çok daha geniş ölçekte konumlandırmayı hedeflemektedir. Bu gelişme, açık kaynaklı bir buyruk kümesi mimarisinin en yüksek başarım gerektiren iş yüklerinde bile rekabetçi olabileceğini gösteren güçlü bir işarettir.
\end{bilginot}

Bu kitapta RISC-V'in tercih edilmesinin başlıca nedeni, öğrencilerin buyruk kümesinden işlemci tasarımına, boru hattından bellek hiyerarşisine kadar tüm kavramları açık ve erişilebilir bir mimari üzerinde uygulayabilmeleridir. RISC-V buyruk kümesinin ayrıntıları Bölüm~\ref{chap:buyruk-kumesi}'te ele alınacaktır.


%% ============================================================
%% BÖLÜM 1.10 — YAPAY ZEKA VE BİLGİSAYAR MİMARİSİ
%% ============================================================
\section{Bilgisayar Mimarisinin Yeni Ufku: Yapay Zekâ}
\label{sec:yapay-zeka-ufku}

Bilgisayar mimarisi alanı, 2010'lardan itibaren yapay zekâ\index{yapay zekâ} (YZ) uygulamalarının muazzam hesaplama gereksinimleriyle yeni bir dönüşüm sürecine girdi. Uzun yıllar boyunca genel amaçlı işlemciler (MİB) hemen her iş yükü için yeterli görülürken, günümüzde derin öğrenme\index{derin öğrenme} modelleri gibi hesaplama yoğun uygulamalar alana özgü hızlandırıcılar olmaksızın verimli biçimde çalıştırılamaz hale geldi.


\subsection{Genel Amaçlıdan Alana Özgü Mimarilere}
\label{sec:genel-alana-ozgu}

Geleneksel bilgisayar mimarisinde tek bir işlemci tüm görevleri (metin işleme, bilimsel hesaplama, grafik oluşturma) üstlenirdi. Ancak transistör bütçesinin artması ve güç duvarının ortaya çıkması, tasarımcıları farklı bir soruya yönlendirdi: \emph{``Her şeyi biraz iyi yapan bir yonga mı, belirli bir şeyi çok iyi yapan bir yonga mı tasarlamalıyız?''}

Bu sorunun yanıtı \emph{alana özgü mimari}\index{alana özgü mimari} (Domain-Specific Architecture: DSA) kavramını doğurdu. Günümüzde yapay zekâ iş yükleri bu dönüşümün en güçlü itici gücü; çünkü derin öğrenme modelleri birkaç temel özellik sergiliyor:
\begin{itemize}
    \item \textbf{Yoğun matris aritmetiği:} Sinir ağlarının her katmanı büyük matris çarpımları ve toplama işlemleri gerektirir.
    \item \textbf{Yüksek koşutluk:} Milyonlarca bağımsız çarpma--toplama işlemi eş zamanlı yürütülebilir.
    \item \textbf{Düşük duyarlıkta yeterli doğruluk:} Eğitim 16 bitle, çıkarım ise 8 hatta 4 bitle yapılabilir; 64 bitlik kayan nokta gerektiren bilimsel hesaplamalardan çok farklıdır.
\end{itemize}
Bu özellikler, genel amaçlı bir MİB yerine bu işlemlere göre özelleştirilmiş bir yonganın çok daha verimli olabileceğini gösteriyor.


\subsection{Yapay Zekâ Hızlandırıcıları: GPU, TPU, NPU}
\label{sec:yz-hizlandiricilar-giris}

Yapay zekâ çağının üç temel donanım yapı taşı şöyle özetlenebilir:

\begin{itemize}
    \item \textbf{GPU\index{GPU} (Grafik İşlem Birimi):} Aslen üç boyutlu grafik işleme için tasarlanmış olan GPU'lar, binlerce küçük çekirdeğiyle yoğun koşut hesaplamalara olanak tanır. NVIDIA'nın CUDA\index{CUDA} platformu, GPU'ları yapay zekâ araştırmacıları için birincil hesaplama aracına dönüştürdü. Günümüzün büyük dil modellerinin eğitimi binlerce GPU içeren kümeler üzerinde haftalarca sürebiliyor.

    \item \textbf{TPU\index{TPU} (Tensör İşlem Birimi):} Google tarafından geliştirilen TPU'lar, \emph{sistolik dizi}\index{sistolik dizi} mimarisiyle matris çarpımını doğrudan donanımda gerçekleştirir. Bu yaklaşım, genel amaçlı bir çekirdeğe kıyasla çok daha yüksek enerji verimliliği sağlıyor.

    \item \textbf{NPU\index{NPU} (Sinir Ağı İşlem Birimi):} Apple Neural Engine, Qualcomm Hexagon ve Google Tensor yongalarında bulunan NPU'lar, taşınabilir ve uç aygıtlarda düşük güçle yapay zekâ çıkarımı gerçekleştirmek için tasarlandı. Fotoğraf iyileştirmeden gerçek zamanlı çeviriye kadar pek çok özellik bu birimlere dayanıyor.
\end{itemize}
Bu hızlandırıcıların ayrıntılı mimarisi Bölüm~\ref{chap:gpu-hizlandiricilar}'de kapsamlı olarak incelenecek.

\begin{figure}[H]
\centering
\begin{tikzpicture}[>=Stealth,
    donem/.style={draw, thick, rounded corners=5pt, minimum height=2.2cm,
                  minimum width=4.2cm, align=center, font=\small},
    bariyer/.style={draw, very thick, dashed, sinyalkirmizi, minimum height=3cm,
                    minimum width=0.6cm, align=center, font=\footnotesize\bfseries,
                    text=sinyalkirmizi, fill=blokkirmizi-acik, rounded corners=2pt}]

    % Sol dönem: Genel Amaçlı MİB
    \node[donem, fill=blokmavi] (cpu) at (0,0)
        {\textbf{Genel Amaçlı}\\[2pt]\textbf{MİB Dönemi}\\[4pt]
         {\footnotesize 1970'ler--2000'ler}\\[2pt]
         {\scriptsize Tek çekirdek, saat sıklığı artışı}};

    % Bariyer: Güç Duvarı
    \node[bariyer] (bariyer) at (5.3,0) {Güç\\Duvarı\\(2006)};

    % Sağ dönem: Heterojen Hesaplama
    \node[donem, fill=blokyesil, minimum width=5cm] (het) at (10.5,0)
        {\textbf{Heterojen Hesaplama}\\[2pt]\textbf{Dönemi}\\[4pt]
         {\footnotesize 2010'lar--günümüz}};

    % Alt bileşenler (heterojen dönem)
    \node[draw, rounded corners=2pt, fill=blokmavi-acik, font=\scriptsize\bfseries,
          minimum width=1cm] at (8.8,-1.8) {MİB};
    \node[draw, rounded corners=2pt, fill=blokyesil-acik, font=\scriptsize\bfseries,
          minimum width=1cm] at (10,-1.8) {GPU};
    \node[draw, rounded corners=2pt, fill=bloksari, font=\scriptsize\bfseries,
          minimum width=1cm] at (11.2,-1.8) {TPU};
    \node[draw, rounded corners=2pt, fill=blokkirmizi-acik, font=\scriptsize\bfseries,
          minimum width=1cm] at (12.4,-1.8) {NPU};

    % Oklar
    \draw[->, very thick, sinyalmavi!70!black] (cpu.east) -- (bariyer.west);
    \draw[->, very thick, sinyalyesil!70!black] (bariyer.east) -- (het.west);

    % Zaman çizelgesi çizgisi (alt)
    \draw[thick, gray!50] (-2.5,-2.5) -- (14,-2.5);
    \foreach \yr/\xp in {1970/-2, 1980/-0.5, 1990/1, 2000/2.5, 2010/7, 2020/12} {
        \draw[gray!50] (\xp,-2.5) -- (\xp,-2.7);
        \node[font=\tiny, text=gray!60] at (\xp,-2.9) {\yr};
    }

\end{tikzpicture}
\caption[Bilgisayar mimarisinde paradigma değişimi]{Bilgisayar mimarisinde paradigma değişimi: güç duvarı genel amaçlı MİB dönemini sonlandırmış, yerini heterojen hesaplamaya bırakmıştır.}
\label{fig:paradigma-degisimi}
\end{figure}


\subsection{Mimari Eğitimine Etkisi}
\label{sec:degisim-egitim}

Alana özgü hızlandırıcılar ne kadar farklı görünürse görünsün, hepsi bu kitapta öğreneceğiniz temel kavramlar üzerine inşa edilmiş. Bir GPU'nun çekirdek yapısı, boru hattı ve önbellek ilkeleriyle anlaşılır. Bir TPU'nun sistolik dizisi, sayısal aritmetik temelleri olmadan kavranamaz. Bellek hiyerarşisi ve bant genişliği hesapları her mimari için geçerli.

Bu nedenle yapay zekâ çağı, bilgisayar mimarisi eğitimini önemsizleştirmek bir yana, daha da \emph{can alıcı} hale getirdi. Buyruk kümesi tasarımı (Bölüm~\ref{chap:buyruk-kumesi}), işlemci mikromimarisi (Bölüm~\ref{chap:islemci-tasarimi}), boru hattı (Bölüm~\ref{chap:boruhatti}) ve bellek hiyerarşisi (Bölüm~\ref{chap:bellek}) gibi konular, hem geleneksel MİB'leri hem de yeni nesil hızlandırıcıları anlamanın temelini oluşturuyor.

RISC-V'in modüler genişletilebilirlik yapısı bu bağlamda özel bir avantaj sunuyor: Öğrenciler, temel RV32I çekirdeğine özel matris çarpımı buyrukları ekleyerek basit bir yapay zekâ hızlandırıcısını prototipleyebilir. Bu deneyim, alana özgü mimari tasarımının özünü kavramak için eşsiz bir fırsat sağlıyor.

\begin{bilginot}[Büyük Dil Modelleri ve Donanım]
ChatGPT ve benzeri büyük dil modelleri\index{büyük dil modeli} (Large Language Models: LLM), on milyarlarca parametreden oluşan yapay sinir ağları. GPT-4 düzeyinde bir modelin eğitimi, on bin GPU'dan oluşan bir küme üzerinde aylarca sürebiliyor ve milyonlarca dolar enerji maliyeti gerektiriyor. Modelin kullanıcılara sunulması (çıkarım) ise daha düşük güçle mümkün, ancak milyonlarca eş zamanlı istek düşünüldüğünde bu da büyük ölçekli bir donanım altyapısı gerektiriyor. Transformer mimarisindeki \emph{dikkat mekanizması}\index{dikkat mekanizması} (attention), bellek bant genişliği açısından son derece yoğun bir işlem olup yeni nesil donanım tasarımlarını doğrudan şekillendiriyor.
\end{bilginot}


% -----------------------------------------------------------------
% Yanılgılar ve Tuzaklar
% -----------------------------------------------------------------
\section{Yanılgılar ve Tuzaklar}
\label{sec:giris-yanilgilar}

Bilgisayar mimarisine giriş aşamasında bazı kavramlar sezgisel olarak doğru görünse de gerçekte yanıltıcıdır. Bu bölümde en yaygın yanlış kanıları ele alacağız.

\subsubsection*{Yanılgı: İşlemci hızı bilgisayarın hızını tek başına belirler}
Bir bilgisayarın başarımı yalnızca işlemciye bağlı değildir. Bellek hızı, depolama erişim süresi, veri yolu bant genişliği ve yazılımın kalitesi de en az işlemci kadar belirleyicidir. Örneğin yavaş bir sabit disk (HDD) kullanan güçlü bir işlemci, hızlı bir SSD'ye sahip daha düşük saat sıklıklı bir sistemden uygulamada daha yavaş çalışabilir.

\subsubsection*{Yanılgı: Daha fazla transistör her zaman daha iyi başarım demektir}
Moore Yasası transistör sayısının artacağını öngörür, ancak bu artış doğrudan başarıma dönüşmez. Dennard ölçeklemesinin\index{Dennard ölçeklemesi} sona ermesiyle birlikte eklenen transistörler daha fazla güç tüketir ve ısınır. Bu nedenle çağdaş işlemciler transistörleri saat sıklığını artırmak yerine daha çok çekirdeğe, önbelleğe ve özelleşmiş birimlere ayırır.

\subsubsection*{Yanılgı: CISC mimarileri RISC'ten yavaştır}
Alt Bölüm~\ref{subsec:risc-cisc}'te ele alındığı gibi, 1980'lerdeki RISC devrimiyle birlikte ``basit buyruklar hızlıdır'' sloganı yerleşmiş olsa da günümüzdeki tablo çok daha karmaşıktır. Çağdaş x86 işlemcileri CISC buyrukları donanım içinde mikro-işlemlere (RISC benzeri) çevirir ve boru hattından geçirir. Apple M serisi (ARM/RISC) ile Intel Core serisi (x86/CISC) arasındaki başarım farkları artık mimari felsefeden çok mikromimari tasarım kararlarına bağlıdır.

\subsubsection*{Tuzak: Tek bir ölçütle bilgisayar seçmek}
GHz, çekirdek sayısı, RAM miktarı gibi tek bir özelliğe bakarak bilgisayar karşılaştırmak yanıltıcıdır. Bir bilgisayarın gerçek başarımı iş yüküne bağlıdır ve birden fazla ölçütün birlikte değerlendirilmesini gerektirir. Pazarlama kampanyaları genellikle tek bir sayıyı öne çıkarır; mühendislik ise bütünsel bir bakış açısı gerektirir.

\subsubsection*{Tuzak: ``Açık kaynak'' ile ``ücretsiz'' kavramlarını karıştırmak}
RISC-V açık kaynaklı bir buyruk kümesi mimarisidir; bu, spesifikasyonun serbestçe kullanılabileceği anlamına gelir. Ancak bir RISC-V işlemcisinin tasarlanması, üretilmesi ve doğrulanması ciddi mühendislik yatırımı gerektirir. Açık kaynak, giriş engelini düşürür ama geliştirme maliyetini sıfırlamaz.


\section{Özet}
\label{sec:giris-ozet}

Bilgisayar mimarisi çok hızlı gelişen bir mühendislik alanıdır. Elektrik-elektronik ve bilgisayar mühendislerinin bir eşgüdüm içinde oluşturduğu değişik soyutlama katmanlarından oluşur. Bir bilgisayarın veri yolu, denetim, bellek, giriş ve çıkıştan oluşan beş ana bileşeni vardır; veri yolu ve denetimi içinde barındıran işlemci en önemli parçasıdır. Programlar üst düzey dillerden derleyiciler aracılığıyla makine diline dönüştürülür ve işletim sistemi, kütüphaneler ile aygıt yazılımından oluşan katmanlı bir yazılım altyapısı üzerinde çalışır. Buyrukların ve verilerin bellekte saklanma biçimine göre Von Neumann (ortak bellek) ve Harvard (ayrı bellekler) olmak üzere iki temel mimari model bulunur; günümüzde çoğu işlemci her ikisinin avantajlarını birleştiren değiştirilmiş Harvard modelini kullanır. Buyrukların tasarımında ise RISC ve CISC olmak üzere iki temel yaklaşım vardır; bu kitapta açık kaynaklı bir RISC mimarisi olan RISC-V kullanılmaktadır.

Bilgisayar tarihinde vakum tüplerinden transistörlere, tümleşik devrelerden mikroişlemcilere ve günümüzün çok çekirdekli heterojen sistemlerine uzanan beş nesil ayırt edilebilir. Moore Yasası'nın öngördüğü transistör artışı onlarca yıl boyunca endüstrinin yol haritasını belirledi, ancak Dennard ölçeklemesinin sona ermesiyle güç duvarı ortaya çıktı. Güç tüketimi, günümüzde işlemci tasarımının en temel kısıtlarından biri. Bilgisayar başarımının nasıl ölçüldüğü ve değerlendirildiği Bölüm~\ref{chap:basarim}'de ayrıntılı olarak ele alınıyor.

RISC-V'in açık kaynak, modüler ve eğitim dostu yapısı, onu hem bu kitabın hem de dünya genelinde giderek artan sayıda mimari projenin tercih ettiği buyruk kümesi haline getirdi. Yapay zekâ ise bilgisayar mimarisinin yeni ufku: GPU, TPU ve NPU gibi alana özgü hızlandırıcılar, genel amaçlı MİB'lerin tek başına karşılayamadığı devasa hesaplama gereksinimlerini verimli biçimde karşılamak için tasarlanıyor. Bu hızlandırıcıların mimarisi Bölüm~\ref{chap:gpu-hizlandiricilar}'de kapsamlı olarak incelenecek.

İlerleyen bölümlerde başarım (Bölüm~\ref{chap:basarim}), buyruk kümesi mimarisi (Bölüm~\ref{chap:buyruk-kumesi}), sayısal aritmetik (Bölüm~\ref{chap:aritmetik}), işlemci tasarımı (Bölüm~\ref{chap:islemci-tasarimi}), boru hattı (Bölüm~\ref{chap:boruhatti}), bellek hiyerarşisi (Bölüm~\ref{chap:bellek}), giriş/çıkış (Bölüm~\ref{chap:giris-cikis}), çok çekirdekli sistemler (Bölüm~\ref{chap:cok-cekirdekli}) ve GPU ile hızlandırıcılar (Bölüm~\ref{chap:gpu-hizlandiricilar}) ele alınacak.


\section*{Alıştırmalar}
\addcontentsline{toc}{section}{Alıştırmalar}
\label{sec:giris-alistirmalar}

\begin{alistirma}[\kolay]
Aşağıdaki sözcüklerin anlamları nedir? (a) Von Neumann Mimarisi, (b) Harvard Mimarisi, (c) CISC, (d) RISC.
\end{alistirma}

\begin{alistirma}[\kolay]
Aşağıdaki kavramları açıklayın: (a) Makine Dili, (b) Derleyici, (c) Çevirici, (d) Programlama Dili, (e) İşletim Sistemi.
\end{alistirma}

\begin{alistirma}[\kolay]
Buyruk yürütüm döngüsünü açıklayın.
\end{alistirma}

\begin{alistirma}[\kolay]
Bir bilgisayarı oluşturan beş ana bileşen nedir?
\end{alistirma}

\begin{alistirma}[\orta]
Geçmişin bilgisayar mimarisine nasıl bir etkisi vardır? Daha önce tasarlanmış işlemcilere uyum mutlaka gerekli midir?
\end{alistirma}

\begin{alistirma}[\kolay]
Aşağıdaki kavramları açıklayın: (a) Vakum tüpü, (b) Transistör, (c) Tümleşik devre, (d) Yonga, (e) Silisyum ve plaka, (f) Moore yasası.
\end{alistirma}

\begin{alistirma}[\orta]
Moore Yasası nedir? Moore Yasası diğer bilim dallarındaki gibi mutlak bir gerçeği mi tanımlar? Açıklayın.
\end{alistirma}

\begin{alistirma}[\orta]
Eğer 2001 yılının Ekim ayında üretilen bir yongada 100 transistör varsa ve yongayı üreten şirket Moore yasasını izliyorsa, şirketin 2007 yılının Ekim ayında ürettiği yeni yongada kaç transistör olur?
\end{alistirma}

\begin{alistirma}[\orta]
Bir üniversitedeki öğrenci sayısı şu anda yaklaşık 2000'dir. Eğer bu sayı Moore yasasında belirlenmiş olan yonga başına düşen transistör sayısının artış hızıyla artarsa üniversitenin mevcudunun 32000 kişi olması için ne kadar zaman geçmesi gerekir?
\end{alistirma}

\begin{alistirma}[\orta]
Plaka maliyetiyle bir yonganın alanı arasındaki bağıntıyı gösterin ve buna göre yonga alanının tasarımdaki önemini yorumlayın.
\end{alistirma}

\begin{alistirma}[\zor]
Günümüz bilgisayar endüstrisinde Moore Yasası'nın yavaşlamasının yarattığı zorlukları ve bu zorluklara karşı geliştirilen mimari çözümleri tartışın. Güç duvarı kavramını açıklayın.
\end{alistirma}

\begin{alistirma}[\orta]
Devingen ve durağan güç tüketimi arasındaki farkları açıklayın. Aynı işlemcinin daha düşük çalışma gerilimiyle çalıştırılması devingen güç tüketimini nasıl etkiler?
\end{alistirma}

\begin{alistirma}[\zor]
RISC ve CISC mimarilerinin avantaj ve dezavantajlarını karşılaştırın. RISC-V'in açık kaynak olmasının eğitim ve endüstri açısından sağladığı faydaları tartışın.
\end{alistirma}

\begin{alistirma}[\zor]
Moore yasası fiilen sona erdi mi, yoksa yalnızca yavaşladı mı? Aşağıdaki senaryoları değerlendirin ve her birinin avantaj ile sınırlılıklarını tartışın:
\begin{enumerate}[(a)]
    \item Transistör boyutlarını küçültmeye devam etmek (EUV litografi, GAA transistörler).
    \item Üçüncü boyuta genişlemek (3B yonga istifleme, HBM bellek).
    \item Alana özgü hızlandırıcılar tasarlamak (GPU, TPU, FPGA).
    \item Yeni hesaplama paradigmaları geliştirmek (kuantum bilgisayarlar, nöromorfik yongalar).
\end{enumerate}
Sizce gelecek on yılda bilgisayar başarımı artışının birincil kaynağı hangisi olacaktır? Yanıtınızı gerekçelendirin.
\end{alistirma}

\begin{alistirma}[\zor]
Yapay zekâ iş yüklerinin (özellikle derin öğrenme eğitimi ve çıkarımı) geleneksel MİB mimarisi üzerindeki taleplerini açıklayın. GPU, TPU ve NPU gibi hızlandırıcıların neden geliştirildiğini mimari tasarım perspektifinden tartışın.
\end{alistirma}

\begin{alistirma}[\orta]
Genel amaçlı işlemci (MİB) ile alana özgü hızlandırıcı (GPU, TPU) arasındaki temel tasarım felsefesi farkını ``gecikme odaklı'' (latency-oriented) ve ``işlem hacmi odaklı'' (throughput-oriented) kavramlarıyla açıklayın.
\end{alistirma}

\begin{alistirma}[\zor]
x86, ARM ve RISC-V buyruk kümesi mimarilerini lisanslama modeli, kullanım alanları ve genişletilebilirlik açısından karşılaştırın. Açık kaynak bir buyruk kümesi mimarisinin ülkelerin teknolojik bağımsızlığı için neden önemli olabileceğini tartışın.
\end{alistirma}

%% ===================================================================
%% Aşağıdaki sorular BİL 361 dersinin geçmiş yıllardan sınav
%% sorularından uyarlanmıştır.
%% ===================================================================

\begin{alistirma}[\kolay]
Aşağıdaki cümlelerde boş bırakılan yerleri doldurun ya da doğru/yanlış (D/Y) olarak yanıtlayın.
\begin{enumerate}[(a)]
    \item Saat vuruş sıklığı, saatin iki yükselen darbesi arasındaki geçen süreye denir. (D/Y)
    \item 2 bayt adresleme yapılan bir mimaride buyruklar 128 bit genişliğindeyse, sonraki buyruğa erişmek için program sayacını \rule{2cm}{0.4pt} kadar artırmalıyız.
    \item \rule{3cm}{0.4pt} mimarisinde buyruk ve veri aynı bellekte bulunur.
    \item İşlemcilerde saat sıklığını devredeki \rule{3cm}{0.4pt} belirler.
    \item İşlemcideki bellek gecikmeleri her zaman sabittir. (D/Y)
    \item Bilgisayarınızın anabelleği \rule{2cm}{0.4pt} teknolojisini kullanır.
    \item Von Neumann mimarisinde veri ve buyruk bellekleri ayrı iken, Harvard mimarisinde ortaktır. (D/Y)
\end{enumerate}
\end{alistirma}

\begin{alistirma}[\orta]
Aşağıdaki ifadelerin her biri için doğruysa D, yanlışsa Y yazın. Yanlış olanlarda gerekçenizi belirtin.
\begin{enumerate}[(a)]
    \item Moore Yasası fiziksel bir yasa değil, gözleme dayalı bir eğilimdir.
    \item RISC mimarisinde buyruklar değişken uzunluktadır.
    \item Harvard mimarisinde buyruk ve veri bellekleri birbirinden ayrıdır.
    \item Aynı buyruk kümesini kullanan iki farklı işlemci, aynı programı farklı sürelerde çalıştırabilir.
\end{enumerate}
\end{alistirma}

\begin{alistirma}[\kolay]
Aşağıdaki boşlukları doldurun.
\begin{enumerate}[(a)]
    \item Bir bilgisayarı oluşturan beş ana bileşen: \rule{2cm}{0.4pt}, \rule{2cm}{0.4pt}, \rule{2cm}{0.4pt}, giriş ve çıkıştır.
    \item Programcı ile donanım arasında bir sözleşme niteliği taşıyan soyutlama katmanına \rule{3cm}{0.4pt} denir.
    \item Von Neumann mimarisinde buyruklar ve veriler \rule{3cm}{0.4pt} bellekte saklanır.
    \item Tüm bileşenleri bir arada tutan bağlantıya \rule{2cm}{0.4pt} denir.
\end{enumerate}
\end{alistirma}

\begin{alistirma}[\zor]
Yeni bir buyruk kümesi tasarlıyorsanız, aşağıdaki tasarım seçeneklerinin her birinin olası artı ve eksilerini tartışın:
\begin{enumerate}[(a)]
    \item Yazmaç sayısının artırılması
    \item Buyruk uzunluğunun sabit mi yoksa değişken mi olması
    \item Anlık değer alanının büyütülmesi
    \item Buyruk türü sayısının artırılması
\end{enumerate}
\end{alistirma}

%% ===================================================================
%% Sayısal Sorular — Yonga Üretimi
%% ===================================================================

\begin{alistirma}[\zor]
Çapı 300\,mm olan bir silisyum plakasının maliyeti 16\,000\,\$ olsun. Üretilecek yonganın kırmık alanı 120\,mm$^2$, birim alan başına hata sayısı 0{,}002\,hata/mm$^2$ ve kritik aşama sayısı $N = 10$ ise:
\begin{enumerate}[(a)]
    \item Plakadan elde edilebilecek yaklaşık kırmık sayısını hesaplayın.
    \item Verimi hesaplayın.
    \item Bir kırmığın üretim maliyetini bulun.
    \item Kırmık alanı iki katına (240\,mm$^2$) çıkarılsaydı, maliyet nasıl değişirdi? Hesaplayın ve yorumlayın.
\end{enumerate}
\end{alistirma}

\begin{alistirma}[\zor]
Bir yonga üreticisi aynı tasarımı iki farklı teknoloji düğümünde üretmeyi karşılaştırıyor:

\smallskip
\begin{tabular}{@{}lcc@{}}
\toprule
 & \textbf{28\,nm} & \textbf{5\,nm} \\
\midrule
Plaka maliyeti & 3\,500\,\$ & 17\,000\,\$ \\
Kırmık alanı & 250\,mm$^2$ & 80\,mm$^2$ \\
Verim & \%88 & \%72 \\
Maske seti maliyeti & 5\,M\,\$ & 60\,M\,\$ \\
\bottomrule
\end{tabular}
\smallskip

300\,mm çaplı plaka kullanıldığını varsayın.
\begin{enumerate}[(a)]
    \item Her iki düğüm için plaka başına düşen kırmık maliyetini hesaplayın.
    \item Her iki düğüm için maske seti payını 10 milyon adetlik ve 500 bin adetlik üretim hacimleri için ayrı ayrı hesaplayın.
    \item Hangi üretim hacminde 5\,nm ekonomik olarak anlamlı hâle gelir? Tartışın.
\end{enumerate}
\end{alistirma}

\begin{alistirma}[\zor]
Yeni bir 3\,nm üretim hattı başlatılıyor. Verim öğrenme eğrisi aşağıdaki gibi ilerliyor:

\smallskip
\begin{tabular}{@{}lccccc@{}}
\toprule
\textbf{Çeyrek dönem} & Ç1 & Ç2 & Ç3 & Ç4 & Ç5 \\
\midrule
\textbf{Verim (\%)} & 35 & 48 & 62 & 74 & 82 \\
\bottomrule
\end{tabular}
\smallskip

Plaka maliyeti 20\,000\,\$, kırmık alanı 100\,mm$^2$ ve plaka çapı 300\,mm ise:
\begin{enumerate}[(a)]
    \item Her çeyrek dönem için kırmık maliyetini hesaplayın.
    \item Kırmık maliyetinin Ç1'den Ç5'e kadar yüzde kaç düştüğünü bulun.
    \item Bir şirket ürününü Ç1'de piyasaya sürmek ile Ç3'ü beklemek arasında karar veriyor. Her iki senaryonun avantaj ve dezavantajlarını maliyet ve rekabet açısından tartışın.
\end{enumerate}
\end{alistirma}

\begin{alistirma}[\orta]
Aşağıdaki iki işlemcinin kırmık alanı ve plaka bilgileri verilmiştir:

\smallskip
\begin{tabular}{@{}lcc@{}}
\toprule
 & \textbf{İşlemci A} & \textbf{İşlemci B} \\
\midrule
Kırmık alanı & 50\,mm$^2$ & 400\,mm$^2$ \\
Hata yoğunluğu & \multicolumn{2}{c}{0{,}001\,hata/mm$^2$} \\
Kritik aşama sayısı ($N$) & \multicolumn{2}{c}{12} \\
\bottomrule
\end{tabular}
\smallskip

\begin{enumerate}[(a)]
    \item Her iki işlemcinin verimini hesaplayın.
    \item 300\,mm çaplı plakadan her işlemciden kaç tane sağlam kırmık elde edilir?
    \item Verim farkının büyük kırmıklar için neden daha dramatik olduğunu açıklayın.
\end{enumerate}
\end{alistirma}

%% ===================================================================
%% Sayısal Sorular — Güç Tüketimi ve Enerji
%% ===================================================================

\begin{alistirma}[\orta]
Bir işlemci 1{,}0\,V kaynak geriliminde ve 3\,GHz saat sıklığında çalışırken devingen güç tüketimi 80\,W'tır.
\begin{enumerate}[(a)]
    \item Kaynak gerilimi 0{,}8\,V'a ve saat sıklığı 2{,}4\,GHz'e düşürülürse devingen güç tüketimi kaç watt olur?
    \item Bu değişiklikle güç tüketimi yüzde kaç azalır?
    \item Saat sıklığının düşürülmesi başarımı nasıl etkiler? Güç--başarım ödünleşimini tartışın.
\end{enumerate}
\end{alistirma}

\begin{alistirma}[\zor]
Aynı görevi çalıştıran iki işlemciyi karşılaştırın:

\smallskip
\begin{tabular}{@{}lcc@{}}
\toprule
 & \textbf{İşlemci X} & \textbf{İşlemci Y} \\
\midrule
Güç tüketimi & 120\,W & 40\,W \\
Görev süresi & 5\,s & 18\,s \\
\bottomrule
\end{tabular}
\smallskip

\begin{enumerate}[(a)]
    \item Her iki işlemci için toplam enerji tüketimini ($E = P \times t$) hesaplayın.
    \item Hangi işlemci enerji açısından daha verimlidir?
    \item Bir veri merkezinde bu görev günde 10\,000 kez çalıştırılıyorsa, yıllık enerji farkı kaç kWh olur?
    \item ``Düşük güçlü işlemci her zaman daha az enerji harcar'' ifadesinin neden yanıltıcı olduğunu bu örnekle açıklayın.
\end{enumerate}
\end{alistirma}

\begin{alistirma}[\orta]
Bir taşınabilir işlemcinin durağan (sızıntı) güç tüketimi $P_{\text{durağan}} = I_{\text{sızıntı}} \times V_{dd}$ formülüyle hesaplanır. İşlemci 0{,}75\,V kaynak geriliminde çalışmaktadır ve toplam sızıntı akımı 2{,}5\,A'dir.
\begin{enumerate}[(a)]
    \item Durağan güç tüketimini hesaplayın.
    \item Devingen güç tüketimi 4\,W ise durağan gücün toplam güç içindeki yüzdesini bulun.
    \item Bu işlemci bir akıllı telefonda kullanılıyor ve telefon bekleme modundayken (devingen güç $\approx 0$) yalnızca durağan güç tüketiyor. 5\,000\,mAh kapasiteli, 3{,}8\,V'luk bir pilde saklanan enerji kaç Wh'tir? Telefon yalnızca sızıntı gücüyle kaç saat dayanır?
\end{enumerate}
\end{alistirma}

%% ===================================================================
%% Sayısal Sorular — Moore Yasası
%% ===================================================================

\begin{alistirma}[\orta]
Moore Yasası'na göre bir yongadaki transistör sayısı yaklaşık her 2 yılda ikiye katlanır.
\begin{enumerate}[(a)]
    \item 2010 yılında 1 milyar transistör içeren bir işlemci ailesi Moore Yasası'nı izlerse, 2024'te kaç transistör içerir?
    \item Gerçekte Apple M4 (2024) yaklaşık 28 milyar transistör içermektedir. Bu değer Moore Yasası tahminiyle uyumlu mu? Tartışın.
    \item Moore Yasası'nın ``yasa'' olarak adlandırılması fizik yasalarıyla karşılaştırıldığında neden yanıltıcıdır?
\end{enumerate}
\end{alistirma}

%% ===================================================================
%% Sayısal Sorular — Güvenilirlik
%% ===================================================================

\begin{alistirma}[\orta]
Bozukluk (\emph{fault}), hata (\emph{error}) ve arıza (\emph{failure}) kavramları arasındaki farkı bir bellek modülü örneği üzerinden açıklayın. Hata düzeltme kodlarının (ECC) bu zincirdeki rolünü tartışın.
\end{alistirma}

\begin{alistirma}[\kolay]
Geçici, aralıklı ve kalıcı bozuklukları tanımlayın. Her biri için birer gerçek dünya örneği verin. Hangisinin tanılanması en zordur? Nedenini açıklayın.
\end{alistirma}

\begin{alistirma}[\zor]
Bir gömülü sistemde dört özdeş bileşen seri bağlı olarak çalışmaktadır. Her bileşenin MTTF'si 250\,000~saattir.
\begin{enumerate}[(a)]
    \item Sistemin MTTF'sini hesaplayın.
    \item Her bileşenin FIT değerini ve sistemin toplam FIT değerini bulun.
    \item MTTR 4~saat ise sistemin kullanılabilirliğini hesaplayın.
    \item Bu kullanılabilirlik ``beş dokuz'' (\%99,999) hedefini karşılıyor mu? Karşılamıyorsa, bileşen MTTF'si en az ne olmalıdır?
\end{enumerate}
\end{alistirma}

\begin{alistirma}[\zor]
Bir veri merkezi 2\,000 sunucudan oluşmaktadır. Her sunucunun FIT değeri 20\,000'dir ve ortalama onarım süresi 1{,}5~saattir.
\begin{enumerate}[(a)]
    \item Tek bir sunucunun MTTF'sini saat ve yıl cinsinden bulun.
    \item Veri merkezinin tamamı için MTTF'yi hesaplayın. Ortalama ne sıklıkla bir sunucu arızası beklenir?
    \item Veri merkezinin kullanılabilirliğini hesaplayın.
    \item Sunucu sayısı 10\,000'e çıkarılırsa MTTF nasıl değişir? Büyük ölçekli sistemlerde yedeklemenin neden zorunlu olduğunu tartışın.
\end{enumerate}
\end{alistirma}

\begin{alistirma}[\zor]
Bir otonom araç işlemcisi için MTTF hedefi en az 1\,000\,000~saat (yaklaşık 114~yıl) olarak belirlenmiştir.
\begin{enumerate}[(a)]
    \item Bu hedefe karşılık gelen FIT değerini hesaplayın.
    \item İşlemci üç alt birimden oluşuyor: hesaplama birimi (FIT\,=\,200), bellek denetleyicisi (FIT\,=\,150) ve giriş/çıkış birimi (FIT\,=\,100). Sistem seri bağlıysa toplam FIT ve MTTF değerlerini bulun. Hedef karşılanıyor mu?
    \item Uzay uygulamalarındaki güvenilirlik gereksinimleri otomotivden nasıl farklıdır? Mars'taki Curiosity aracını örnek vererek tartışın.
\end{enumerate}
\end{alistirma}

\begin{alistirma}[\orta]
Güvenilirlik (\emph{reliability}), dayanıklılık (\emph{durability}) ve gürbüzlük (\emph{robustness}) kavramlarını bir askeri dizüstü bilgisayar tasarımı bağlamında karşılaştırın. Her kavram için bu bilgisayarda alınması gereken bir tasarım kararı örneği verin.
\end{alistirma}
